% ; - part core ; -
\maketitle
\noteprovenance

\begin{abstract}
We prove that an increasing sequence of positive integers satisfying
$a_n^2/a_{n+1}=1+3/n+o(n^{-3})$ has an irrational reciprocal sum.
Rationality would force integer numerators of the reciprocal tails to
agree eventually with a cubic polynomial. Divisibility, a calculation in
a cubic number field and congruences modulo seven exclude that polynomial.
For rational reciprocal sums, we also develop sufficient conditions for
the eventual recurrence $a_{n+1}=a_n^2-a_n+1$ under the more general limit
$a_{n+1}/a_n^2\to1$. They control either increases of an integer tail
numerator or a convergent sum at steps setting new maxima. The paper
gives the underlying recurrences, their relation to earlier work, and
examples separating numerical growth bounds from the exact recurrences. The cubic
proof and the further written arguments are distinguished from the
specific theorems formalised in Lean. None of the additional bounds is
derived for every rational reciprocal sum in Erd\H{o}s Problem~\#243.
\end{abstract}

\section{Introduction}\label{long243:sec:problem}
Sylvester's sequence $2,3,7,43,1807,\ldots$ is generated by
$a_{n}=a_{n-1}^{2}-a_{n-1}+1$ and has reciprocal sum $1$; the recurrence is
equivalent to the telescoping identity
\[
 \frac1{a_{n-1}-1}=\frac1{a_{n-1}}+\frac1{a_n-1}.
\]
Consequently the tail beginning with $1/a_{n-1}$ equals $1/(a_{n-1}-1)$.
Erd\H{o}s Problem~\#243 asserts that this is the only way a sequence of that
growth can have a rational reciprocal sum:

\begin{problem}[Erd\H{o}s \#243]\label{long243:res:problem}
Let $1\le a_1<a_2<\cdots$ be a sequence of integers with
\[
 \lim_{n\to\infty}\frac{a_n}{a_{n-1}^{2}}=1
 \qquad\text{and}\qquad
 \sum\frac{1}{a_n}\in\Q .
\]
Then $a_n=a_{n-1}^{2}-a_{n-1}+1$ for all sufficiently large $n$.
\end{problem}

\noindent See~\cite[p.~64]{erdosgraham1980} and~\cite[p.~105]{erdos1988}.
The supplied catalogue snapshot lists the problem as
open~\cite{erdosproblems}.
Below we index the same recurrence as
$a_{n+1}=a_n^{2}-a_n+1$, which is the shift the formal sources use; the two
forms are the same statement. The polynomial $a^2-a+1$ is called the
\lword{Erdos243/ReciprocalTailRigidity.lean}{37}{sylvesterNext}{Sylvester
successor}.

The hypothesis is asymptotic and the conclusion is exact.  Clearing a rational
tail as $C_n/D_n$ exposes that distinction in the integer
$E_n=D_n-(a_n-1)C_n$: a Sylvester tail has $E_n=0$, while quadratic growth
only gives $E_n/C_n\to0$.  Integrality alone cannot turn this relative
estimate into $E_n=0$ when $C_n$ grows.  The exact dynamics retain extra
information: divisors persist in the unreduced denominator $D_n$,
and coprimality restricts the later reduced numerators once the gcd has
stabilised. When cancellation continues, persistence in the reduced
denominator must be proved separately. These are the divisibility
properties used by the bounded-negative and record-crossing arguments.

The cubic proof is independent of the arguments about bounded increments.
For the bounded-increment proof accompanying the short note,
Sections~\ref{long243:sec:state} and~\ref{long243:sec:defect} give the
recurrences and zero-error argument;
Sections~\ref{long243:sec:barrier} and~\ref{long243:sec:bounded} supply
the gcd stabilisation and CRT contradiction.
Sections~\ref{long243:sec:lcmrecords} and~\ref{long243:sec:records}
use two different integer numerators, respectively obtained by clearing
denominators with an LCM and by reducing a fraction to lowest terms.
Section~\ref{long243:sec:mass} gives the independent, elementary argument
for a convergent sum of relative increases. The constant and periodic
cases appear in Sections~\ref{long243:sec:constant}
and~\ref{long243:sec:periodic}.

Section~\ref{long243:sec:priorwork} compares the results with prior work
and identifies their formalisation status. Section~\ref{long243:sec:open}
separates questions about reciprocal tails from examples that merely
avoid prescribed moduli. Appendices~\ref{long243:app:index}
and~\ref{long243:app:residue} give source locations and a finite residue
calculation.

We use $z_+=\max(z,0)$. Deleting a finite prefix is harmless for an eventual
recurrence, but not for the coefficients in a higher-order rate. In
particular, the index $n$ in $1+3/n+o(n^{-3})$ is kept fixed throughout
the cubic proof. The recurrence sections also use zero-based indexing,
with that convention stated where the integer sequences are introduced.

\medskip
\noindent\textbf{Keywords.} irrationality; Ahmes series; Sylvester's sequence;
unit fractions; Lean~4.
\qquad
\textbf{MSC 2020.} 11J72 (primary); 11B37, 11D68, 68V20 (secondary).

\section{Irrationality at the cubic rate}\label{long243:sec:cubicrate}

The rate in the next theorem is much more specific than
$a_{n+1}/a_n^2\to1$: both the coefficient $3/n$ and the error
$o(n^{-3})$ are required. It excludes a Sylvester tail, whose deviation
from $1$ is $O(1/a_n)$. The coefficient $3$ also puts it outside the bounded-increment
criterion: the increments of the product ratio grow like a positive
multiple of $n^2$. We first construct the integer numerators that rationality would
require, and then prove that they cannot have the resulting cubic form.
The little-$o$ condition is used in the integer finite-difference argument
below; it cannot simply be replaced there by $O(n^{-3})$.

\begin{theorem}[cubic-rate irrationality]\label{long243:res:cubicrate}
If strictly increasing positive integers satisfy
$a_n^2/a_{n+1}=1+3/n+o(n^{-3})$, then $\sum_n1/a_n$ is irrational.
\end{theorem}
\leannote{Lean: \lproof{ErdosProblems/Erdos243/PaperCompleteR21/GaloisSignFlipClosure.lean}{270}{cubic_rate_irrationality_of_chebotarev}. Conditional on Chebotarev's density theorem; see the \hyperref[long243:sec:coverage]{coverage section}.}

For a concrete sequence satisfying the hypothesis, take
\[
 a_1=8,\qquad a_{n+1}=\left\lceil\frac{n a_n^2}{n+3}\right\rceil
 \quad(n\ge1).
\]
It begins $8,16,103,5305,\ldots$. The inequality $a_{n+1}\ge a_n^2/4$
gives $a_n\ge4\cdot2^{2^{n-1}}\ge8$ by induction. Hence
$a_{n+1}\ge a_n^2/4\ge2a_n$, so the sequence is strictly increasing.
Rounding upwards gives
\[
 0\le1+\frac3n-\frac{a_n^2}{a_{n+1}}<\frac{16}{a_n^2}=o(n^{-3}).
\]
Thus its reciprocal sum is irrational. The proof of the theorem uses the
following arithmetic obstruction.

\begin{theorem}[rising-factorial cubic exclusion]\label{long243:res:cubicexclusion}
Let $a,C,D:\N\to\Z_{>0}$ satisfy
\begin{equation}\label{long243:eq:cubicorbit}
 C_{n+1}=a_nC_n-D_n,\qquad D_{n+1}=a_nD_n .
\end{equation}
Then for every $A\in\Q_{>0}$ and every $B\in\Q$,
\[
 \liminf_{X\to\infty}
 \frac{\#\{n\le X:C_n\ne A\,n(n+1)(n+2)+B\}}{X}>0 .
\]
In particular no such orbit satisfies $C_n=A\,n(n+1)(n+2)+B$ for all
large $n$.
\end{theorem}
\leannote{Lean: \lproof{ErdosProblems/Erdos243/PaperCompleteR21/GaloisSignFlipClosure.lean}{256}{cubic_exclusion_of_chebotarev}. Conditional on Chebotarev's density theorem; see the \hyperref[long243:sec:coverage]{coverage section}.}

The irrationality theorem needs only exclusion of eventual equality.
Theorem~\ref{long243:res:cubicexclusion} proves more: disagreement with
the proposed cubic has positive lower density. The lower bound may depend
on the orbit and the cubic; it is not a uniform numerical constant.

The exclusion has three steps. Agreement on arbitrarily late blocks of
four consecutive indices gives a uniform bound on the common divisor of
numerator and denominator. Dividing out its eventual value reduces the
proposed cubic to
$m\,n(n+1)(n+2)/6+c$, with $m$ a positive integer and $c=\pm1$.
Next, the two-step numerator recurrence forces a square condition at
primes dividing the middle of three consecutive numerators. Chebotarev's
theorem turns this into a square in a cubic number field, and a trace
calculation forces $m=12$. Finally, the two remaining cubics fail the
exact recurrence modulo seven, at different specified residue classes of
the index. The lemmas below carry out these steps in that order.

\paragraph{The integer tail.}
Suppose $\sum_{n\ge1}1/a_n=p/q$ with positive integers $p,q$, and write
$x_n=\sum_{k\ge n}1/a_k$ for the tail.  Put
\[
 D_n=q\prod_{1\le k<n}a_k,\qquad C_n=D_nx_n .
\]
Then $D_n$ is a positive integer, and
\[
 C_n=D_n\Bigl(\frac pq-\sum_{1\le k<n}\frac1{a_k}\Bigr)
    =p\prod_{1\le k<n}a_k-q\sum_{1\le k<n}\ \prod_{\substack{1\le j<n\\ j\ne k}}a_j
\]
is a positive integer as well.  From $x_n=1/a_n+x_{n+1}$,
\[
 C_{n+1}=D_{n+1}x_{n+1}=a_nD_n\Bigl(x_n-\frac1{a_n}\Bigr)=a_nC_n-D_n,
 \qquad D_{n+1}=a_nD_n,
\]
which is \eqref{long243:eq:cubicorbit}.  These are the tail variables of Koizumi's
Lemma~4 \cite[pp.~11--12]{koizumi2025}, up to a common positive factor.
Section~\ref{long243:sec:state} constructs these variables together with
the centred error $E_n=D_n-(a_n-1)C_n$, which this section does not need.  The subsequent cubic exclusion uses only these integer recurrences;
it does not assume a reciprocal sum. To apply its zero-based statement
without shifting the rate hypothesis, adjoin $a_0=1$, $D_0=q$ and
$C_0=p+q$. The recurrences at $0$ then give $D_1=q$ and $C_1=p$,
while every original index $n\ge1$ is unchanged. The auxiliary term
$a_0$ need not satisfy the increasing-sequence hypothesis of the
irrationality theorem: the exclusion requires only positive multipliers.

\subsection{Reduction under the zero lower-density assumption}

For $S\subseteq\N$ write $\underline d(S)=\liminf_X\#(S\cap[1,X])/X$.
The condition $\underline d(S)=0$ only requires these proportions to
approach zero along a sequence of cutoffs; it does not require their
limit to exist. The argument below uses exactly this weaker assumption.
For a sequence $F$, let $\Delta F_n=F_{n+1}-F_n$; higher powers of
$\Delta$ mean repeated forward differences. This operator is distinct
from the later indexed quantity $\Delta_n$, the Sylvester defect.

\begin{lemma}[periodic obstruction principle]\label{long243:res:periodicobstruction}
Let $h,L$ be positive integers and let $j_1,\ldots,j_L$ be fixed integer
offsets. Suppose that for every sufficiently large $n$ in one residue
class modulo $h$ at least one of $n+j_1,\ldots,n+j_L$ lies in $S$.
Then $\underline d(S)\ge1/(Lh)$.
\end{lemma}
\leannote{Lean: \lproof{ErdosProblems/Erdos243/PaperCompleteR11/DensityTransport.lean}{156}{fixed_offsets_periodic_lowerDensity}.}

\begin{proof}
There are $X/h+O(1)$ relevant starting indices below $X$. Their witnesses
lie below $X+O(1)$ because the offsets are fixed, and each element of $S$
serves as a witness for at most $L$ starting indices. Hence
$L\#(S\cap[1,X])\ge X/h-O(1)$, which gives the claimed lower density.
\end{proof}

Fix $A\in\Q_{>0}$, $B\in\Q$, and suppose for contradiction that
\[
 P(n)=A\,n(n+1)(n+2)+B,
 \qquad S=\{n:C_n\ne P(n)\},
 \qquad \underline d(S)=0 .
\]
Everything through the end of Section~\ref{long243:sec:modseven} runs under that
assumption.

\begin{lemma}[gcd stabilisation and the primitive shape]\label{long243:res:gcdshape}
Write $G_n=\gcd(C_n,D_n)$.  Then $6A$ is a positive integer, $G_n$ divides
$6A$ for every $n$, and $G_n$ is eventually equal to a positive integer
$g$.  On the tail where $G_n=g$, put $u_n=C_n/g$, $v_n=D_n/g$ and
$Q(n)=P(n)/g$.  Then
\begin{equation}\label{long243:eq:primitivetail}
 u_{n+1}=a_nu_n-v_n,\qquad v_{n+1}=a_nv_n,\qquad \gcd(u_n,v_n)=1,
 \qquad\gcd(u_n,u_{n+1})=1,
\end{equation}
and there are $m\in\Z_{>0}$ and $c\in\{-1,1\}$ with
\begin{equation}\label{long243:eq:Qmc}
 Q(n)=\frac m6n(n+1)(n+2)+c .
\end{equation}
\end{lemma}
\leannote{Lean: \lproof{ErdosProblems/Erdos243/PaperCompleteR21/CubicProfileGcdShape.lean}{103}{cubic_profile_gcd_stabilisation_and_primitive_shape}.}

\begin{proof}
Equation~\eqref{long243:eq:cubicorbit} gives $G_n\mid G_{n+1}$ and $G_n\mid C_t$ for
every $t\ge n$.  Since $\underline d(S)=0$, beyond every threshold there is a block of four
consecutive indices disjoint from $S$: otherwise every block of four would meet
$S$ and $\underline d(S)\ge1/4$.  On such a block $\Delta^3C_t=\Delta^3P(t)=6A$, and
$G_n$ divides each of the four values, hence divides $6A$.  So $6A$
is a positive integer and the divisibility chain $(G_n)$ is bounded and
stabilises at some $g$.

On the primitive tail, $\gcd(u_n,v_n)=1$ by construction, and a prime dividing
$u_n$ and $u_{n+1}$ would divide $v_n=a_nu_n-u_{n+1}$, which gives
$\gcd(u_n,u_{n+1})=1$.

The polynomial $Q$ is integer-valued: it takes an integer value at every
integer argument, although its coefficients need not all be integers. Indeed,
let $N$ clear the denominators of its coefficients, so $NQ\in\Z[X]$ and
$Q(n+N)-Q(n)\in\Z$ for every $n$.  A non-integral value at one argument would
therefore force an exception at every late argument in one residue class modulo
$N$, and Lemma~\ref{long243:res:periodicobstruction} with $L=1$ would contradict
$\underline d(S)=0$.  Now $Q(-1)=Q(0)=B/g$, so $c:=B/g\in\Z$, and
$m:=Q(1)-Q(0)=6A/g$ is a positive integer, which is \eqref{long243:eq:Qmc} with
$c\in\Z$.

It remains to exclude $c=0$ and every $c$ with a prime factor.  Let $\ell$ be a
prime dividing $c$, or any prime if $c=0$.  For $n\equiv-1\pmod{6\ell}$, write
$n+1=6\ell k$; then
\[
 \frac{n(n+1)(n+2)}6=\ell k(6\ell k-1)(6\ell k+1),
 \qquad
 \frac{(n+1)(n+2)(n+3)}6=\ell k(n+2)(n+3),
\]
so $\ell$ divides both $Q(n)$ and $Q(n+1)$.  Agreement at both indices would
give $\ell\mid\gcd(u_n,u_{n+1})=1$.  Hence one of $n,n+1$ lies in $S$ for every
late $n\equiv-1\pmod{6\ell}$, and Lemma~\ref{long243:res:periodicobstruction} with
$L=2$, $h=6\ell$ gives $\underline d(S)\ge1/(12\ell)$.  This contradiction
leaves $c=\pm1$.
\end{proof}

\begin{lemma}[the multiplier supply, and the reducible case]\label{long243:res:reduciblecase}
On the primitive tail, $\gcd(a_n,v_n)=1$, the multipliers at distinct indices
are pairwise coprime, infinitely many of them exceed $1$, and the cubic
$Q_{m,c}$ of \eqref{long243:eq:Qmc} is irreducible over $\Q$.
\end{lemma}
\leannote{Lean: \lproof{ErdosProblems/Erdos243/PaperCompleteR11/PrimitiveMultiplierSupply.lean}{244}{primitive_zero_density_paper_multiplier_lemma}.}

\begin{proof}
A prime dividing $a_n$ and $v_n$ would divide $u_{n+1}$ and $v_{n+1}$, against
$\gcd(u_{n+1},v_{n+1})=1$; and every earlier multiplier divides every later
$v_n$, so distinct multipliers are coprime.  If $a_n=1$ for all large $n$ then
$v_n$ is eventually a positive constant and $u_{n+1}=u_n-v_n$ decreases without
bound, against positivity.  So infinitely many distinct primes divide late
multipliers.

Suppose $Q_{m,c}$ had a rational root.  Choose a prime $\ell\nmid6m$ dividing a
late multiplier $a_N$, large enough that the root reduces modulo $\ell$.  Then
$Q_{m,c}$ vanishes modulo $\ell$ on a full residue class, while $\ell\mid v_n$
and therefore $\ell\nmid u_n$ for every $n>N$.  Agreement at any late index of
that class is impossible, and
Lemma~\ref{long243:res:periodicobstruction} with $L=1$, $h=\ell$ gives
$\underline d(S)\ge1/\ell$.
\end{proof}

\subsection{A square condition from three consecutive numerators}

At a prime dividing a middle numerator, the two-step recurrence forces the
negative product of its neighbours to be a square. For the proposed
cubic this becomes a square condition on $r^2-1$ at every root of a
polynomial modulo almost every prime. The next lemma turns that condition
into a square in the cubic number field.

\begin{lemma}[square specialisation]\label{long243:res:squarespec}
Let $f\in\Q[T]$ be irreducible with root $\alpha$, and let $H\in\Q[T]$ satisfy
$H(\alpha)\ne0$.  If for all but finitely many primes $\ell$ every root
$r\in\Fp_\ell$ of the reduction of $f$ has $H(r)$ a square in $\Fp_\ell^\times$,
then $H(\alpha)$ is a square in $\Q(\alpha)^\times$.
\end{lemma}
\leannote{Lean: \lproof{ErdosProblems/Erdos243/PaperCompleteR21/GaloisSignFlipClosure.lean}{211}{squareSpecialisation_of_chebotarev}. Conditional on Chebotarev's density theorem; see the \hyperref[long243:sec:coverage]{coverage section}.}

\begin{proof}
Suppose not.  Put $K=\Q(\alpha)$ and choose $\beta$ with $\beta^2=H(\alpha)$,
so $K(\beta)/K$ is quadratic.  Take a finite Galois extension $L/\Q$ containing
$K(\beta)$ and $\sigma\in\Gal(L/K)$ restricting nontrivially to $K(\beta)$, so
that $\sigma\alpha=\alpha$ and $\sigma\beta=-\beta$.  By the Chebotarev density
theorem in the form stated by Stevenhagen and Lenstra
\cite[\S3, p.~15]{stevenhagenlenstra1996},
there are infinitely many unramified rational primes whose Frobenius
class in $\Gal(L/\Q)$ is that of $\sigma$; only the infinitude is used.
Discard the finitely many primes dividing denominators or the discriminant,
the prime $2$, and those at which $\beta$ reduces to zero.  The Frobenius elements at the primes of $L$ above a fixed $\ell$
form that conjugacy class, so one of those primes has Frobenius exactly
$\sigma$. At this prime the residue field satisfies
$\overline\alpha^{\,\ell}=\overline\alpha$ and
$\overline\beta^{\,\ell}=-\overline\beta$, so $\overline\alpha\in\Fp_\ell$ while
$\overline\beta\notin\Fp_\ell$.  Then $H(\overline\alpha)=\overline\beta^{\,2}$
is a nonsquare in $\Fp_\ell$ at a root of the reduction of $f$, against the
hypothesis.
\end{proof}

The hypothesis is needed at every root and at all but finitely many primes.  A
favourable finite set of tested primes would not suffice.

\begin{proposition}[the square forced by the numerator recurrence]\label{long243:res:transportsquare}
Write $\kappa=m/6$ and $\eta=6c/m$, so that $Q_{m,c}(n)=\kappa f(n+1)$ for
$f(T)=T^3-T+\eta$.  Then $\alpha^2-1$ is a square in $\Q(\alpha)^\times$ for a
root $\alpha$ of $f$.
\end{proposition}
\leannote{Lean: \lproof{ErdosProblems/Erdos243/PaperCompleteR21/GaloisSignFlipClosure.lean}{240}{transport_square_of_chebotarev}. Conditional on Chebotarev's density theorem; see the \hyperref[long243:sec:coverage]{coverage section}.}

\begin{proof}
Eliminating $v_n$ from \eqref{long243:eq:primitivetail} gives
\begin{equation}\label{long243:eq:secondorderu}
 u_{n+2}=(a_n+a_{n+1})u_{n+1}-a_n^2u_n .
\end{equation}
Let $\ell$ be a prime and $k$ an index with $\ell\mid u_k$.  Reading
\eqref{long243:eq:secondorderu} at $n=k-1$ modulo $\ell$ gives
$u_{k+1}\equiv-a_{k-1}^2u_{k-1}$, hence
\begin{equation}\label{long243:eq:forcedsquare}
 -u_{k-1}u_{k+1}\equiv(a_{k-1}u_{k-1})^2 \pmod\ell,
\end{equation}
so $-u_{k-1}u_{k+1}$ is a square modulo $\ell$.

Let $\ell\nmid6m$ and let $r\in\Fp_\ell$ be a root of $f$.  Then $r\ne0$ and
$r\ne\pm1$, since $f(0)=\eta\ne0$ and $f(\pm1)=\eta$.  Using $\eta=r-r^3$,
\[
 f(r-1)=-3r(r-1),\qquad f(r+1)=3r(r+1),
\]
so at an index $k\equiv r-1\pmod\ell$ we may reduce the polynomial values in
$\Fp_\ell$. All equalities in the following calculation are in that field:
$Q(k)=0$, $Q(k-1)=-3\kappa r(r-1)$, $Q(k+1)=3\kappa r(r+1)$, and
\begin{equation}\label{long243:eq:threevalue}
 -Q(k-1)Q(k+1)=9\kappa^2r^2(r^2-1).
\end{equation}
If $r^2-1$ were a nonsquare modulo $\ell$, then so would be the right side of
\eqref{long243:eq:threevalue}, because $9\kappa^2r^2$ is a nonzero square; agreement at
all three of $k-1,k,k+1$ would then contradict \eqref{long243:eq:forcedsquare}.  That
prohibition recurs at every index of the class $r-1$ modulo $\ell$, and
Lemma~\ref{long243:res:periodicobstruction} with $L=3$, $h=\ell$ would give
$\underline d(S)\ge1/(3\ell)$.  So $r^2-1$ is a square in $\Fp_\ell^\times$ at
every root of every good reduction, and Lemma~\ref{long243:res:squarespec} applies with
$H(T)=T^2-1$.
\end{proof}

\subsection{The square condition forces \texorpdfstring{$m=12$}{m=12}}

\begin{lemma}[classification of the scales]\label{long243:res:scaletwelve}
Let $m$ be a positive integer, let $c\in\{-1,1\}$ and $\eta=6c/m$, and suppose $T^3-T+\eta$ is
irreducible over $\Q$ with a root $\alpha$ such that $\alpha^2-1$ is a square
in $\Q(\alpha)^\times$.  Then $m=12$.
\end{lemma}
\leannote{Lean: \lproof{ErdosProblems/Erdos243/PaperCompleteR20/ScaleTwelve.lean}{36}{scale_twelve_of_square_in_rootField}.}

\begin{proof}
Choose $\beta\in\Q(\alpha)$ with $\beta^2=\alpha^2-1$. The substitution
$z=\alpha+\beta$ is useful because its inverse is already in the same
field and expresses $\alpha$ symmetrically in $z,z^{-1}$.
Then $(\alpha+\beta)(\alpha-\beta)=1$, so
\begin{equation}\label{long243:eq:zalpha}
 z^{-1}=\alpha-\beta,\qquad \alpha=\tfrac12(z+z^{-1}),
\end{equation}
and $z$ generates $\Q(\alpha)$. Write its minimal polynomial as
$Z^3+bZ^2+dZ+w$ with $b,d,w\in\Q$ and $w\ne0$. We use the first two
traces of $\alpha$ to restrict these coefficients, then the third trace
to impose $\eta=6c/m$. All traces below are from $\Q(\alpha)$ to $\Q$.
The polynomial $T^3-T+\eta$ gives
\begin{equation}\label{long243:eq:alphatraces}
 \Tr\alpha=0,\qquad \Tr\alpha^2=2,\qquad \Tr\alpha^3=-3\eta .
\end{equation}
Newton's identities give $\Tr z=-b$ and $\Tr z^{-1}=-d/w$, so the first
equation in \eqref{long243:eq:alphatraces} and \eqref{long243:eq:zalpha} give
\begin{equation}\label{long243:eq:dbw}
 d=-bw .
\end{equation}
For the second trace, Newton's identities and $d=-bw$ give
\[
 \Tr z^2=b^2+2bw,\qquad \Tr z^{-2}=b^2-2b/w.
\]
Since the trace of $1$ is $3$, squaring
\eqref{long243:eq:zalpha} yields
$4\Tr\alpha^2=\Tr z^2+6+\Tr z^{-2}$. Substituting
$\Tr\alpha^2=2$ gives $b^2+b(w-w^{-1})=1$.
Multiplication by $w$ factors this equation:
\begin{equation}\label{long243:eq:factored}
 (bw-1)(b+w)=0 .
\end{equation}
For the third trace, the same identities give
\[
 \begin{aligned}
 \Tr z^3&=-b^3-3b^2w-3w,\\
 \Tr z^{-3}&=b^3-3b^2/w-3/w.
 \end{aligned}
\]
The cross terms in $(z+z^{-1})^3$ have trace
$3(\Tr z+\Tr z^{-1})=0$ by \eqref{long243:eq:dbw}.
Thus $8\Tr\alpha^3=\Tr z^3+\Tr z^{-3}$.
Using $\Tr\alpha^3=-3\eta$ now gives
\begin{equation}\label{long243:eq:etatrace}
 \eta=\frac{(b^2+1)(w+w^{-1})}8 .
\end{equation}
By \eqref{long243:eq:factored}, either $b=-w$ or $b=w^{-1}$, and \eqref{long243:eq:etatrace}
becomes
\[
 \eta=\frac{(w^2+1)^2}{8w}
 \qquad\text{or}\qquad
 \eta=\frac{(w^2+1)^2}{8w^3}.
\]
Write $w=r/s$ with $r\in\Z\mathbin{\backslash}\{0\}$, $s\in\Z_{>0}$ and $\gcd(r,s)=1$.
Since $\eta=6c/m$,
\begin{equation}\label{long243:eq:mformula}
 m=\frac{48c\,rs^3}{(r^2+s^2)^2}
 \qquad\text{or}\qquad
 m=\frac{48c\,r^3s}{(r^2+s^2)^2}.
\end{equation}
Now $\gcd(r^2+s^2,rs)=1$, so integrality of $m$ and $|c|=1$ force
$(r^2+s^2)^2\mid48$, hence $r^2+s^2\in\{1,2,4\}$.  With $r\ne0$ and $s\ge1$ the
value $1$ is too small, and $4$ is not a sum of two nonzero squares, so
$r^2+s^2=2$ and $|r|=s=1$.  Then \eqref{long243:eq:mformula} gives $m=12cr$, and
$m>0$ gives $m=12$.
\end{proof}

The square condition leaves the two possible cubics
\begin{equation}\label{long243:eq:survivors}
 Q_{12,\pm1}(n)=2n(n+1)(n+2)\pm1 .
\end{equation}
The square condition uses three consecutive numerators. The last step
also uses the preceding denominator update, so it tests four consecutive
numerators instead.

\subsection{The two remaining cubics fail modulo seven}\label{long243:sec:modseven}

\begin{lemma}[the two forbidden words]\label{long243:res:modseven}
For $Q_{12,1}$, every sufficiently late four-index window beginning at
$n\equiv0\pmod7$ contains an exceptional index. For $Q_{12,-1}$, the same
holds for windows beginning at $n\equiv1\pmod7$. In either case,
$\underline d(S)\ge1/7$.
No phase-free four-index obstruction is asserted.
\end{lemma}
\leannote{Lean: \lproof{ErdosProblems/Erdos243/PaperCompleteR20/TwoForbiddenWords.lean}{18}{plus_one_forbidden_word}, \lproof{ErdosProblems/Erdos243/PaperCompleteR20/TwoForbiddenWords.lean}{38}{minus_one_forbidden_word}.}

\begin{proof}
Reduce modulo $7$.  For $c=1$ the values of $Q_{12,1}$ at $n=0,1,2,3$ are
$1,13,49,121$, that is
\begin{equation}\label{long243:eq:wordplus}
 (1,6,0,2),
\end{equation}
and for $c=-1$ the values at $n=1,2,3,4$ are $11,47,119,239$, that is
\begin{equation}\label{long243:eq:wordminus}
 (4,5,0,1).
\end{equation}
Each four-term pattern recurs with period $7$.

Index a proposed occurrence locally by $0,1,2,3$, and use $a_i$ for the
multiplier at local index $i$. Let $(c_0,c_1,0,c_3)$ be the residues of
$u$ and $(d_0,d_1,d_2,d_3)$ those of $v$. Thus
$c_{i+1}=a_ic_i-d_i$ and $d_{i+1}=a_id_i$ modulo $7$.  The middle zero gives
$a_1c_1=d_1$, then $d_2=a_1d_1$ and $c_3=-d_2$, so
\begin{equation}\label{long243:eq:d1square}
 d_1^2=c_1\,a_1d_1=c_1d_2=-c_1c_3 .
\end{equation}
For \eqref{long243:eq:wordplus} this is $-6\cdot2=2$, and for \eqref{long243:eq:wordminus} it
is $-5\cdot1=2$.  The square roots of $2$ modulo $7$ are $3$ and $4$, so
$d_1\in\{3,4\}$ in both cases.

The first step gives $a_0=(c_1+d_0)/c_0$ and hence
$d_1=d_0(d_0+c_1)/c_0$, which is legitimate because $c_0\in\{1,4\}$ is
invertible.  For \eqref{long243:eq:wordplus} this is $d_1=d_0(d_0+6)$, with values
\[
 (0,0,2,6,5,6,2) \qquad\text{at } d_0=0,1,\ldots,6,
\]
and for \eqref{long243:eq:wordminus} it is $d_1=2d_0(d_0+5)$, with values
\[
 (0,5,0,6,2,2,6).
\]
Both images are $\{0,2,5,6\}$, which misses $\{3,4\}$.  All seven initial
denominator residues have been displayed, so the calculation is exhaustive and
rests on no unreported search.

Hence every sufficiently late window $[7k,7k+3]$ for the plus profile,
and every sufficiently late window $[1+7k,4+7k]$ for the minus profile,
contains an element of $S$. Windows within each family are pairwise disjoint.
There are $X/7+O(1)$ such windows contained in $[0,X]$, so choosing one
exceptional index in each gives $\#(S\cap[0,X])\ge X/7-O(1)$ and therefore
$\underline d(S)\ge1/7$. This stronger local bound is not promoted to a
uniform lower bound for arbitrary cubic profiles.
\end{proof}

\begin{proof}[Proof of Theorem~\ref{long243:res:cubicexclusion}]
Assume $\underline d(S)=0$.  Lemmas~\ref{long243:res:gcdshape}
and~\ref{long243:res:reduciblecase} put the primitive tail in the shape \eqref{long243:eq:Qmc}
with $c=\pm1$ and $Q_{m,c}$ irreducible.  Each excluded branch already
contradicts the assumption by producing a positive, branch-dependent lower
density, namely $1/(12\ell)$ or $1/\ell$ for the relevant prime $\ell$.
Proposition~\ref{long243:res:transportsquare} and Lemma~\ref{long243:res:scaletwelve} force
$m=12$, and Lemma~\ref{long243:res:modseven} then gives
$\underline d(S)\ge1/7$, the final contradiction.
\end{proof}

\subsection{From an asymptotic ratio to an exact polynomial}

We now connect the arithmetic exclusion to the rate in the original
sequence. First compare $C_n$ with a sequence whose consecutive ratio is
exactly $1+\lambda/n$. The remainder will have every positive-order
forward difference tending to zero. A sufficiently high difference of
$C_n$ then tends to zero as well; because it is an integer, it must
vanish eventually. This is the step that turns the asymptotic estimate
into a polynomial identity.

\begin{lemma}[integer extraction at a regular rate]\label{long243:res:extraction}
Let $\lambda>1$ and let $C_n$ be positive integers with
\begin{equation}\label{long243:eq:regularrate}
 \frac{C_{n+1}}{C_n}=1+\frac\lambda n+o(n^{-\lambda}).
\end{equation}
Then $\lambda$ is an integer $d\ge2$, and there are $A\in\Q_{>0}$ and
$B\in\Q$ with $C_n=A\,n(n+1)\cdots(n+d-1)+B$ for all large $n$.
\end{lemma}
\leannote{Lean: \lproof{ErdosProblems/Erdos243/PaperCompleteR21/RegularRateExtraction.lean}{874}{regular_rate_extraction}, \lproof{ErdosProblems/Erdos243/PaperCompleteR21/RegularRateExtraction.lean}{1151}{regular_rate_extraction_cubic}.}

\begin{proof}
Put $F_n=\Gamma(n+\lambda)/\Gamma(n)$, so that $F_{n+1}/F_n=1+\lambda/n$
exactly and $F_n\sim n^\lambda$ by the gamma-ratio asymptotic
\cite[5.11.12]{dlmf_gamma}. Write $\varepsilon_n=o(n^{-\lambda})$ for the
error in \eqref{long243:eq:regularrate} and $z_n=C_n/F_n$, so that
$z_{n+1}/z_n=1+\varepsilon_n/(1+\lambda/n)$.  Since $\lambda>1$ the errors are
absolutely summable, so the product converges and $z_n\to K$ for some $K>0$,
with $z_n-K=o(n^{1-\lambda})$. To see the stated error, put
$\eta_n=\sup_{k\ge n}k^\lambda|\varepsilon_k|\to0$; the tail of the
logarithmic product has absolute value at most
$O(\eta_n\sum_{k\ge n}k^{-\lambda})=o(n^{1-\lambda})$.
The factors are positive eventually, so the limiting product is nonzero.  Hence
\begin{equation}\label{long243:eq:dsmall}
 \delta_n:=C_n-KF_n=F_n(z_n-K)=o(n).
\end{equation}
The recurrence gives
$\delta_{n+1}-\delta_n=(\lambda/n)\delta_n+\varepsilon_nC_n$, in which the
first term is $o(1)$ by \eqref{long243:eq:dsmall} and the second is $o(1)$ because
$C_n=O(n^\lambda)$.  So $\Delta\delta_n\to0$, and therefore
$\Delta^j\delta_n\to0$ for every $j\ge1$.

For the comparison sequence, the Gamma recurrence gives
\[
 \Delta F_n
 =\frac{\Gamma(n+1+\lambda)}{\Gamma(n+1)}
  -\frac{\Gamma(n+\lambda)}{\Gamma(n)}
 =\lambda\,\frac{\Gamma(n+\lambda)}{\Gamma(n+1)}.
\]
Iterating this identity gives
\[
 \Delta^jF_n=\lambda(\lambda-1)\cdots(\lambda-j+1)\,
 \frac{\Gamma(n+\lambda)}{\Gamma(n+j)},
\]
whose right side is $O(n^{\lambda-j})$.  Choose an integer $j>\lambda$; then
$\Delta^jF_n\to0$, so $\Delta^jC_n\to0$.  These are integers, so
$\Delta^jC_n=0$ for all large $n$. Choose $N$ beyond this threshold.
The Newton forward-difference formula gives
\[
 C_{N+t}=\sum_{i=0}^{j-1}\binom ti\Delta^iC_N\qquad(t\ge0),
\]
so $C_n$ agrees eventually with a polynomial with rational coefficients.
Its growth $C_n\asymp n^\lambda$ forces
the degree to be $\lambda$, so $\lambda=d$ is an integer, and $d\ge2$ because
$\lambda>1$.  For that integer $F_n=n(n+1)\cdots(n+d-1)$ exactly, so
\eqref{long243:eq:dsmall} says that the difference of two polynomials of degree $d$ is
$o(n)$, hence constant. The leading coefficient of the rational polynomial
is $K$, so $K\in\Q_{>0}$; the constant difference is rational as well.
This gives the stated shape with $A=K$ and $B\in\Q$.
\end{proof}

The precision in \eqref{long243:eq:regularrate} carries weight.  The positive integers
$C_n=n(n+1)(n+2)+(-1)^n$ satisfy $C_{n+1}/C_n=1+3/n+O(n^{-3})$ and are not
eventually polynomial, so the little-$o$ remainder cannot in general be replaced by a big-$O$
remainder at the same exponent.
That countermodel is a scalar sequence rather than an exact orbit, and it bears
on Lemma~\ref{long243:res:extraction} alone.

\begin{lemma}[the tail ratio reads the growth defect]\label{long243:res:tailratio}
Let $a_n$ be strictly increasing positive integers with $a_{n+1}/a_n^2\to1$ and
$\sum_n1/a_n$ rational, and let $(C_n,D_n)$ be the integer tail.
Write $\gamma_n=a_n^2/a_{n+1}-1$.  Then
\[
 \frac{C_{n+1}}{C_n}=1+\gamma_n+O(a_n^{-1}),
 \qquad a_n\ge\exp(c\,2^n)\ \text{eventually for some }c>0 .
\]
\end{lemma}
\leannote{Lean: \lproof{ErdosProblems/Erdos243/PaperCompleteR7/QuantitativeTail.lean}{204}{canonical_tail_ratio_quantitative}.}

\begin{proof}
Since $a_{n+1}/a_n^2\to1$ and $a_n\to\infty$, there is an $N$ with
$a_{n+1}\ge a_n^2/2\ge2a_n$ for $n\ge N$.  Iterating $a_{n+1}\ge a_n^2/2$ from
an index where $a_n>2$ gives $a_n\ge\exp(c2^n)$ eventually.  The same doubling
gives, for $m>N$,
\[
 \frac1{a_m}\ \le\ x_m\ =\ \sum_{k\ge m}\frac1{a_k}\ \le\ \frac2{a_m},
\]
so $x_{m}=(1+\rho_{m})/a_{m}$ with
$0\le\rho_{m}=a_mx_{m+1}\le 2a_m/a_{m+1}\le4/a_m$, using
$a_{m+1}\ge a_m^2/2$.  Hence
\[
 \frac{C_{n+1}}{C_n}=\frac{a_nD_nx_{n+1}}{D_nx_n}
 =a_n\cdot\frac{(1+\rho_{n+1})/a_{n+1}}{(1+\rho_n)/a_n}
 =\frac{a_n^2}{a_{n+1}}\cdot\frac{1+\rho_{n+1}}{1+\rho_n}.
\]
The last factor is $1+O(1/a_n)$ and $a_n^2/a_{n+1}$ is bounded, so the product
is $a_n^2/a_{n+1}+O(a_n^{-1})$.
\end{proof}

The exact identity in Section~\ref{long243:sec:lcmrecords} gives the same
estimate: $C_{n+1}/C_n=1-E_n/C_n$, while
$0<\gamma_n+E_n/C_n<3/a_n$ eventually. The identity follows from
Theorem~\ref{long243:res:defect}; its bound also uses vanishing relative
error and near-quadratic growth. Equation~\eqref{long243:eq:weightedgrowth}
is the resulting weighted-sum application, not the identity itself.

\begin{proof}[Proof of Theorem~\ref{long243:res:cubicrate}]
The hypothesis $a_n^2/a_{n+1}=1+3/n+o(n^{-3})$ gives $a_{n+1}/a_n^2\to1$.
Suppose $\sum_n1/a_n$ were rational and pass to the integer tail, so
that $(a,C,D)$ satisfies \eqref{long243:eq:cubicorbit} with every $C_n$ and $D_n$ a
positive integer.  By Lemma~\ref{long243:res:tailratio} and $a_n\ge\exp(c2^n)$, the
error term $O(a_n^{-1})$ is $o(n^{-3})$, so
\[
 \frac{C_{n+1}}{C_n}=1+\frac3n+o(n^{-3}).
\]
Lemma~\ref{long243:res:extraction} at $\lambda=3$ gives
$C_n=A\,n(n+1)(n+2)+B$ for all large $n$, with $A\in\Q_{>0}$ and $B\in\Q$.
That contradicts Theorem~\ref{long243:res:cubicexclusion}.
\end{proof}

No restriction on the rational normalisation was assumed: the proof
derives the permitted primitive leading coefficient. The two residue
patterns \eqref{long243:eq:wordplus} and \eqref{long243:eq:wordminus},
the image $\{0,2,5,6\}$ modulo $7$ and the condition
$(r^2+s^2)^2\mid48$ have been calculated above. The public checkpoint
also contains formalised polynomial extraction, primitive cubic
normalisation and finite transport results, including
\href{https://github.com/wcook04/plectis-erdos/blob/3d6d938d696fed0fb71dd55115a18a73738ff223/lean/ErdosProblems/Erdos243/PaperCompleteR11/CubicZeroDensityShape.lean\#L112}{the primitive-shape theorem}. These are components, not an assembled proof of
Theorem~\ref{long243:res:cubicrate}. No declaration of that complete
irrationality implication has been identified in the supplied source index.
Lemma~\ref{long243:res:squarespec} uses the classical finite-Galois
Chebotarev theorem; its use must be justified by the number-field argument,
not by finite residue computations.

Theorem~\ref{long243:res:cubicrate} assumes neither rationality nor an
integer orbit: these enter only under the supposition that the reciprocal
sum is rational. Its rate implies $a_{n+1}/a_n^2\to1$, so it treats a
subclass of the growth sequences in Problem~\ref{long243:res:problem}.
It proves irrationality on that subclass, rather than constructing a
Sylvester recurrence. Sylvester tails have the much smaller deviation
$O(1/a_n)$ and do not belong to it.

Writing $P_n=\prod_{k<n}a_k$, at the cubic rate $P_n/a_n$ grows like
$n^3$ and its increment like $n^2$.
Thus neither the bounded-increment criterion nor Koizumi's nonpositive
upper-limit criterion applies. The rate also fails his $1+o(1/n)$
condition, and Duverney's signed defect series diverges. We do not decide
whether the LCM-weighted criteria apply; that requires information about
the repeated factors in the prefix product. Conversely,
Theorem~\ref{long243:res:cubicexclusion} assumes only positive integer
$a_n,C_n,D_n$ and the recurrences~\eqref{long243:eq:cubicorbit}, not the
near-quadratic growth limit. The two arguments therefore retain distinct
hypotheses.

In familiar approximation notation,
$C_n=q\bigl(\prod_{k<n}a_k\cdot\sum_m1/a_m-\sum_{k<n}\prod_{j<n,\,j\ne k}a_j\bigr)$
is a linear form in the reciprocal sum with integer coefficients.
Here these integer remainders grow like $n^3$ instead of tending to zero.
The precision $o(n^{-3})$ lets finite differences identify their exact
polynomial form; the recurrence then excludes it. The integer example
$n(n+1)(n+2)+(-1)^n$ above explains why an $O(n^{-3})$ error alone
does not justify that polynomial conclusion.

\section{Prior work and formalisation}\label{long243:sec:priorwork}

\paragraph{Classical criteria and pseudo-greedy expansions.}
The rational-tail method predates these coordinates. Erd\H{o}s and Straus
\cite[Theorem~2.1, pp.~85--86]{erdosstraus1974} use eventual integer
remainders under a small-numerator hypothesis. In polynomial Cantor
series, Han\v{c}l and Tijdeman split the numerator into finitely many
shifted products. Rationality is then characterised by the vanishing of
the resulting polynomial sum
\cite[Theorem~2.2 and its proof, pp.~39--40]{hancltijdeman2008}.
Their theorem provides methodological context, not a result for arbitrary
arrays of signed integer coefficients: the rearrangement of the infinite
sum needs boundary control, supplied there by the polynomial hypothesis.
Neither result supplies the first-crossing argument for a lower error
bound proved below.

Several classical results give sufficient conditions directly on the
sequence $(a_n)$. Their hypotheses use different quantities, so the
comparisons must be made separately. Koizumi's product criterion is
implied by eventual nonnegativity of the integer error $E_n$, under the
identification of coordinates given below. It therefore already covers
the descent argument in Section~\ref{long243:sec:descent}
(see~\cite[Cor.~4(1) and Prop.~1(1), pp.~14--15]{koizumi2025}).
The original Erd\H{o}s--Straus criterion instead uses a least common multiple.
Erd\H{o}s and Straus proved that if $\lim a_n/a_{n-1}^{2}=1$, the reciprocal
sum is rational, and $(a_n)$ has no Sylvester tail, then
\[
 \limsup_{n\to\infty}\ \frac{[a_1,\ldots,a_n]}{a_{n+1}}
 \left(\frac{a_{n+1}^{2}}{a_{n+2}}-1\right)>0,
\]
where $[a_1,\ldots,a_n]$ is the least common
multiple~\cite[Theorem~3, p.~132]{erdosstraus1964}.  This is the indexing in
the original theorem. With $A_r=\operatorname{lcm}(a_1,\ldots,a_{r-1})$
and $r=n+1$, its expression is exactly
$(A_r/a_r)(a_r^2/a_{r+1}-1)$, the LCM quantity in the short note.
The index change is not an additional difference between those two
criteria. Erd\H{o}s's 1988 survey, followed by the supplied
catalogue summary, instead prints $a_n^2/a_{n+1}-1$ in the second factor while
retaining the same prefix-LCM quotient; that displayed summary is off by one
and is not followed here~\cite[p.~105]{erdos1988}\cite{erdosproblems}.
Koizumi records the convenient sufficient rate
$a_n^2/a_{n+1}=1+o(1/n)$ under which the Erd\H{o}s--Straus criterion settles
the problem~\cite[Remark~3, p.~16]{koizumi2025}.  Tijdeman and Yuan give a
criterion of the same kind for Ahmes series with positive integer numerators,
weighted by the least common multiple of the earlier denominators
\cite{tijdemanyuan2002}.

Duverney's signed criterion gives a further comparison. Let $(a_n)$
be positive integers tending to infinity and $\epsilon_n\in\{-1,1\}$.
Under the sufficient hypothesis
$\sum_{n\ge0}|a_{n+1}/a_n^{2}-1|<\infty$, the sum
$\sum_{n\ge0}\epsilon_n/a_n$ is rational if and only if
\[
 a_{n+1}=a_n^2-(\epsilon_{n+1}/\epsilon_n)a_n
             +\epsilon_{n+2}/\epsilon_{n+1}
\]
for all large $n$~\cite[Corollary~3.2, p.~287]{duverney2001}.
The all-positive specialisation is the form relevant here.

There is a qualification concerning the convergence assumption.
Display~(3.6) in that corollary has no absolute values. In the proof on
pp.~299--300, the reduced auxiliary fractions $p'_n/q'_n$ satisfy
$p'_{n+1}\mid q'_n$ and
\[
 p'_n\le p'_N\prod_{k=N}^{n-1}\frac{q'_k}{p'_k}.
\]
The required boundedness follows if
$\prod_{k\ge N}(p'_k/q'_k)$ has a positive limit. Absolute convergence
of the growth-defect series ensures this, using the summable error in
Duverney's estimate~(3.3). Signed convergence alone does not justify
that product step, even for positive rational factors. For each integer
$m\ge2$, take the pair $1+1/m$, $1-1/m$ successively $m$ times.
The deviations cancel after every pair, and the intervening partial sums
are $1/m\to0$, so their series converges. But the product through the
block $m=N$ is
\[
 \prod_{m=2}^{N}\left(1-\frac1{m^2}\right)^m
 =\frac{(N+1)^N}{2N^{N+1}}\sim\frac e{2N}\longrightarrow0;
\]
the equality follows by cancelling powers of consecutive integers.
This example does not impose $p'_{n+1}\mid q'_n$. It refutes only the
general product inference, not Duverney's arithmetic criterion.
We therefore use only the absolute-convergence form; the interpretation
with merely signed convergence is not needed in any proof here.

In the all-positive case $\epsilon_n=1$, absolute convergence also gives a
direct comparison with the classical product criterion. With the
one-based notation $P_n=\prod_{1\le j<n}a_j$ used above,
\[
 \frac{P_{n+1}/a_{n+1}}{P_n/a_n}=\frac{a_n^2}{a_{n+1}}.
\]
The ratios on the right have an absolutely convergent sum of deviations
from $1$, because the same is true of their reciprocals and those
reciprocals tend to $1$. Thus $P_n/a_n$ has a positive finite limit,
and its increments tend to zero. This case already satisfies the
classical nonpositive-upper-limit condition; the bounded-increment
result allows a larger class of sequences.

A different quantitative question is treated by Duverney, Kurosawa and
Shiokawa \cite[Theorem~1, author-version p.~2; proof in Section~3]{duverneykurosawashiokawa2020}.
For rational $x_n>1$, their result assumes eventual $x_{n+1}\ge x_n^2$
and control of the accumulated denominators of $x_{n+1}/x_n^2$; it computes
the irrationality exponent of a signed reciprocal series. We use it only as
a restricted comparison. In particular, a Sylvester tail and the cubic rate
considered here approach quadratic growth from the other side.

Badea's positive-term criterion is adjacent but different. For a convergent
series $\sum_n b_n/a_n$ with $a_n,b_n$ positive integers, eventual strict
inequality
\[
 a_{n+1}>\frac{b_{n+1}}{b_n}a_n^2-\frac{b_{n+1}}{b_n}a_n+1
\]
forces irrationality, while rationality under the corresponding non-strict
inequality forces eventual equality~\cite[Theorem~A, p.~313; Cor.~2.2,
p.~316]{badea1993}. For $b_n=1$, its hypothesis is
$a_{n+1}\ge a_n^2-a_n+1$, an inequality between consecutive denominators,
not the one-step sign condition $E_n\ge0$ used in integer descent.
Under the standing rational-tail and growth hypotheses, however, either
inequality imposed eventually forces a Sylvester tail and hence the
other. Their eventual forms are equivalent in this setting; that fact
does not supply either condition for a general signed error.
The general positive-coefficient criterion is not identified as a checked
theorem in the supplied evidence.

Koizumi's pseudo-greedy expansion~\cite{koizumi2025} chooses $a_n$ by
rounding $x_n^{-1}+1$ to the nearest integer, with a half-integer rounded
upwards. Here $x_n$ is the remaining sum before subtracting $1/a_n$.
Thus $a_n=\lfloor x_n^{-1}+3/2\rfloor$, and the gap
$\varepsilon_n=x_n^{-1}+1-a_n$ is the signed rounding error.
Write the rational initial sum as $r=p/q$ with positive integers $p,q$. His Lemma~4
~\cite[pp.~11--12]{koizumi2025} produces integers $c_n>0$,
$d_n$ and $e_n$ with $x_n=c_n/d_n$ and $\varepsilon_n=e_n/c_n$, satisfying
\[
 a_n=\frac{d_n-e_n}{c_n}+1,\qquad
 c_{n+1}=c_n-e_n,\qquad
 d_{n+1}=a_nd_n ,
\]
and the proof of that lemma also gives $c_{n+1}=a_nc_n-d_n$.
These are the recurrences of Section~\ref{long243:sec:state}, after
matching the starting index and the initial normalization. For a sequence
satisfying the hypotheses of Problem~\ref{long243:res:problem}, the rounding rule is guaranteed
only after a finite restart~\cite[Corollary~3, p.~9]{koizumi2025}. Restarting there
with the tail in lowest terms gives $(c_n,d_n,e_n)$; the original
$(C_n,D_n,E_n)$ on that tail is a fixed positive integer multiple of
this triple, with the indices matched. Thus $E_n/C_n=\varepsilon_n$
is unchanged by the restart. The integer $E_n$ is not necessarily a
numerator in lowest terms, and the rounding range is not asserted
before the restart. The first identity
above gives $E_n=D_n-(a_n-1)C_n$; the remaining identities are exactly
the numerator and denominator recurrences, including
Proposition~\ref{long243:res:update}. On an all-negative tail we later
write $e_n=-E_n>0$ for the magnitude. That local use of lower-case $e_n$
has the opposite sign to Koizumi's signed integer $e_n$.

His Theorem~3~\cite[pp.~12--13]{koizumi2025} proves the equivalence of
two assertions. Conjecture~1~\cite[p.~4]{koizumi2025} asks whether, for
a positive rational $r$, the condition $\varepsilon_n\to0$ forces
$\varepsilon_n=0$ eventually. Question~1~\cite[p.~3]{koizumi2025} is the
question of Erd\H{o}s and Graham stated as
Problem~\ref{long243:res:problem} above.
Two implications used below already appear in these coordinates: his Lemma~3, that $\varepsilon_n=0$ forces $\varepsilon_{n+1}=0$,
is the absorption of Theorem~\ref{long243:res:absorb}, and his Proposition~1(2), that
$\varepsilon_n\ge0$ for all large $n$ forces $\varepsilon_n=0$ for all large
$n$, is the descent of Theorem~\ref{long243:res:descent}.  Both statements are given in~\cite[Lemma~3, p.~10; Prop.~1(2), p.~14]{koizumi2025}.
Koizumi attributes Proposition~1(2) to Badea.  Its contrapositive says that a
counterexample must have negative errors infinitely often; it assumes the sign
is eventually nonnegative, while Theorem~\ref{long243:res:bounded} below assumes only
that the negative part is eventually bounded. It permits negative errors
between $-B$ and $0$ and imposes no independent upper bound on positive
errors; the relative-error limit is still required.

The growth hypothesis in Problem~\ref{long243:res:problem} is calibrated by two facts
about Sylvester's sequence. With $a_1=2$, it satisfies
$a_n\sim c_0^{2^{n}}$ with $c_0=1.2640847\ldots$. Deleting initial terms
and reindexing produces sequences with $a_n\sim C^{2^{n}}$ for arbitrarily
large $C$ whose reciprocals still sum
to a rational number~\cite[arXiv v4, p.~2]{kovactao2024}.  The classical sufficient
condition for irrationality, $\lim_n a_n^{1/2^{n}}=\infty$, is therefore
sharp.  Kova\v{c} and Tao identify that condition as folklore
~\cite[arXiv v4, p.~2]{kovactao2024}; they attribute the sharpness observation to
Erd\H{o}s (1975).  A sequence with
$a_n/a_{n-1}^{2}\to1$ has $a_n^{1/2^{n}}$
convergent, so that criterion says nothing about the sequences of
Problem~\ref{long243:res:problem}, and rationality is genuinely possible there.
Sylvester's
sequence is \texttt{A000058} in the OEIS; \texttt{A129871} is the variant
$1,2,3,7,43,\ldots$ with an initial $1$ prepended. For
the recent literature on irrationality of Ahmes series we refer to Kova\v{c}
and Tao~\cite{kovactao2024}, who resolve several problems of Erd\H{o}s and
Graham drawn from the same two sources cited above, and whose introduction
gives a sample of the intermediate work, including S\'andor (1984) and Badea
(1987).  They do not treat Problem~\#243; the rigidity conclusion asked for
there is not among their results.

Crmari\'{c} and Kova\v{c} \cite[Theorems~1--2]{crmarickovac2025}
address the neighbouring Problem~\#270, not Problem~\#243. Allowing integer
$f(n)\to\infty$ in $\sum_n(\prod_{j=1}^{f(n)}(n+j))^{-1}$ gives every
positive real value; imposing nondecreasing $f$ gives a measure-zero value
set. The latter assertion does not rule out individual rational values.
Their extension of Kakeya's subsum argument explains why decay alone need
not force irrationality when the summands remain freely selectable. The
present exact denominator recurrence imposes additional compatibility, so
neither direction is an implication between their theorem and ours.

The \emph{Formal Conjectures} collection contains a mathematically equivalent
unproved declaration, up to its zero-based indexing
~\cite{formalconjectures243}.  Its summand is $\Q$-valued, so Lean's
\texttt{Summable} hypothesis asserts the existence of a sum in $\Q$; the
finite indexing shift changes that sum only by a rational prefix.  The
declaration therefore does encode the rationality premise, but its proof is
\texttt{sorry}.  Its role here is statement-level prior art; it supplies no
proof authority, and the development in this note is independent of it.
The checked results in this project include the state implications and the
specified theorems about the reciprocal sequence; they are not confined to
the state system. Earlier formal work on the remainder method should also be
distinguished from a solution of this problem: Koutsoukou-Argyraki and Li's
Archive of Formal Proofs entry \cite{kouli2020} records an Isabelle/HOL
formalisation of Erd\H{o}s--Straus (1974), Theorem~2.1, Corollary~2.10 and
Theorem~3.1. That is formal prior art for those classical criteria, not an
existing formal proof of Erd\H{o}s~\#243.

\paragraph{Results and proof status.}
The cubic-rate theorem (Theorem~\ref{long243:res:cubicrate}) and the three
results its proof passes through are checked in the development, with the
input of Lemma~\ref{long243:res:squarespec} as their only assumption;
Section~\ref{long243:sec:coverage} states that assumption in full. The
bounded-negative theorem
(Theorem~\ref{long243:res:bounded}) and the scalar summability theorem
(Theorem~\ref{long243:res:mass}) have checked declarations. The checked
theorems include the estimates for rational tails, proved from a product bound
and from the least common multiple of the denominators.
The complete weighted-record statement (Theorem~\ref{long243:res:weightedrecord})
is kernel-checked in the main development, alongside its earlier proof source
in the separate release.
For the further record criteria, each evidence paragraph identifies the
ordinary argument and the checked lemma it uses. In particular,
Lemma~\ref{long243:res:oddpowersupply} below proves that the reduced
denominator supplies arbitrarily large odd prime powers, and that proof is
kernel-checked in the main development. The static-model conclusions in
Section~\ref{long243:sec:open} are proved by ordinary arguments and are
checked in the development as well.

The additional bound or convergence assumption is not derived from the
problem's hypotheses. An equivalence between convergence and the desired
recurrence does not prove either of them. In the stable-gcd case the record-increment theorem
gives a positive lower bound for the log-log coefficient, not necessarily an
infinite coefficient. All conclusions about nonzero tails retain failure of eventual
Sylvester behaviour as a hypothesis.

The lower-error-bound proof is independent of the cubic argument: it
uses exact updates, absorption, gcd stabilisation and a CRT first
crossing. The scalar proof uses a product bound and integer descent.
Section~\ref{long243:sec:open} states the resulting necessary conditions
on a counterexample; it does not classify mixed-sign tails.

\section{Integer numerators, denominators and errors}\label{long243:sec:state}

Write a rational reciprocal tail as $C_n/D_n$, without requiring the
fraction to be reduced. Removing $1/a_n$ and clearing its denominator
gives
\[
 D_{n+1}=a_nD_n,\qquad C_{n+1}=a_nC_n-D_n.
\]
We measure the difference from a Sylvester tail by
$E_n=D_n-(a_n-1)C_n$. These are Koizumi's recurrences
\cite[Lemma~4, pp.~11--12]{koizumi2025}; their formal definitions are the
\lword{Erdos243/ReciprocalTailRigidity.lean}{41}{nextDenState}{denominator
update}, the
\lword{Erdos243/ReciprocalTailRigidity.lean}{45}{nextTailState}{numerator update},
and the
\lword{Erdos243/ReciprocalTailRigidity.lean}{49}{centeredState}{error}.

We also study these identities for arbitrary integer sequences, before
constructing any reciprocal series. An \emph{exact orbit} means sequences
$a_n,C_n,D_n$ satisfying the two displayed recurrences at every index;
$E_n$ always denotes the error just defined. Here and in the following
recurrence sections the indices may start at $0$. The terms $a_n$ are the
multipliers, and $C_n,D_n$ are the numerator and denominator, not
necessarily coprime. An exact orbit need not obey the pseudo-greedy
rounding rule. That rule gives $-C_n/2\le E_n<C_n/2$, whereas the
absorption argument below needs only $|E_n|<C_n$. At a negative error we
write $e_n=-E_n>0$. For an isolated step we use $a,D,C\in\Z$ without
subscripts.

\begin{proposition}[update law]\label{long243:res:update}
For all $a,D,C\in\Z$,
\[
 aC-D=C-\bigl(D-(a-1)C\bigr).
\]
Consequently, every exact orbit satisfies $C_{n+1}=C_n-E_n$.
\end{proposition}
\leannote{Lean: \lproof{ErdosProblems/Erdos243/ReciprocalTailRigidity.lean}{57}{nextTailState_eq_sub_centered}.}

\begin{proof}
$C-\bigl(D-(a-1)C\bigr)=aC-D$.
\end{proof}

\noindent The identity is formalised as the
\lword{Erdos243/ReciprocalTailRigidity.lean}{57}{nextTailState_eq_sub_centered}{update
law}. Thus $C_n$ decreases when $E_n>0$, increases when $E_n<0$, and
stays unchanged when $E_n=0$. This is why bounds on the negative error
become bounds on upward increments.

\begin{example}[the Sylvester orbit]\label{long243:ex:sylvester}
Take $a_n=2,3,7,43,1807,\ldots$ with $a_{n+1}=a_n^{2}-a_n+1$, and start the
state at $D_0=C_0=1$.  The two updates give

\begin{center}
\begin{tabular}{@{}rrrrr@{}}
\toprule
$n$ & $a_n$ & $D_n$ & $C_n$ & $E_n$\\
\midrule
$0$ & $2$ & $1$ & $1$ & $0$\\
$1$ & $3$ & $2$ & $1$ & $0$\\
$2$ & $7$ & $6$ & $1$ & $0$\\
$3$ & $43$ & $42$ & $1$ & $0$\\
$4$ & $1807$ & $1806$ & $1$ & $0$\\
\bottomrule
\end{tabular}
\end{center}

\noindent and $D_n=a_n-1$ with $C_n=1$ at every index: if $D_n=a_n-1$ and
$C_n=1$ then $C_{n+1}=a_n-(a_n-1)=1$ and $D_{n+1}=a_n(a_n-1)=a_{n+1}-1$.
Hence $E_n=D_n-(a_n-1)C_n=0$ throughout and the numerator never moves, which
is Proposition~\ref{long243:res:update} in the stationary case.
\end{example}

\paragraph{Construction from the reciprocal sum.}
Let $T_n=\sum_{k\ge n}1/a_k$ and suppose $T_0\in\Q$.  Put $D_0$ equal to a
common denominator and $D_n=D_0a_0\cdots a_{n-1}$, and set $C_n=D_nT_n$.  Then
$C_n\in\Z$ for every $n$, and from $T_n=1/a_n+T_{n+1}$ one gets
$C_{n+1}=a_nC_n-D_n$, which is the tail update.  On Sylvester's sequence
$T_n=1/(a_n-1)$ exactly, so $D_n=(a_n-1)C_n$ and $E_n=0$: the error
measures deviation from the Sylvester tail identity, and it vanishes
identically on the Sylvester orbit.

The elementary identities in \texttt{ReciprocalTailRigidity.lean} are stated
for an abstract integer system. Their reciprocal-tail interpretation
is supplied separately by the checked \texttt{PaperCompleteR7} modules,
including \href{https://github.com/wcook04/plectis-erdos/blob/3d6d938d696fed0fb71dd55115a18a73738ff223/lean/ErdosProblems/Erdos243/PaperCompleteR7/ProductDefect.lean\#L211}{the product condition for a rational reciprocal sum}. The natural-state identities used in this passage include the
\lword{Erdos243/ReciprocalTailRigidity.lean}{1869}{natTail_eq_nextTailState}{tail realisation},
the \lword{Erdos243/ReciprocalTailRigidity.lean}{1881}{natDen_eq_nextDenState}{denominator realisation},
and the \lword{Erdos243/ReciprocalTailRigidity.lean}{1978}{natTail_eq_sub_centeredState}{signed update}.
Thus the abstract integer identities and their application to a rational
reciprocal sum have separate formal statements.

Conversely, the recurrences identify a reciprocal sum when the ratio
$C_n/D_n$ tends to zero. Regard that ratio as a real number, the
\lword{Erdos243/ReciprocalTailRigidity.lean}{1901}{tailRatio}{tail ratio}.
Assume $a_n>0$, $D_n>0$, $C_{n+1}+D_n=a_nC_n$ and
$D_{n+1}=a_nD_n$. Dividing the numerator update by $a_nD_n$ gives
$C_n/D_n=1/a_n+C_{n+1}/D_{n+1}$, the
\lword{Erdos243/ReciprocalTailRigidity.lean}{1905}{tailRatio_eq_reciprocal_add_next}{one-step
reciprocal identity}. Iterating gives
\[
 \frac{C_0}{D_0}=\sum_{n<N}\frac1{a_n}+\frac{C_N}{D_N},
\]
as in the \lword{Erdos243/ReciprocalTailRigidity.lean}{1929}{tailRatio_eq_partialReciprocalSum_add}{finite
telescoping identity}. If $C_N/D_N\to0$, taking the limit identifies
$\sum_n1/a_n=C_0/D_0$, the
\lword{Erdos243/ReciprocalTailRigidity.lean}{1955}{reciprocalSeries_hasSum_of_tailRatio_tendsto_zero}{series
realisation}.  No growth, error bound or separate rationality hypothesis is used in
these three identities. Positivity and the stated limit of $C_n/D_n$
are still required. They prove the direction from recurrences to a series;
the construction from a rational reciprocal sum was given above and is
also formalised in separate declarations.

Under the relative-error hypothesis used below, the required limit is
automatic. Suppose $a_n>1$, $C_n>0$, $D_n\ge0$ and $E_n/C_n\to0$ on an
exact integer orbit. If $D_0=0$, then $D_n=0$ and
$E_n/C_n=1-a_n\le-1$ at every index, a contradiction. Thus $D_0\ge1$ and
$D_n\ge D_0 2^n$. For each $\varepsilon>0$, the update gives
$C_{n+1}\le(1+\varepsilon)C_n$ eventually. Taking $0<\varepsilon<1$
shows that $C_n/D_n\to0$, so the telescoping identity realises the
reciprocal sum as $C_0/D_0$.

The same assumptions also imply the quadratic growth limit. Put
$\theta_n=E_n/C_n$. Then $a_n=D_n/C_n+1-\theta_n\to\infty$, and the
exact recurrences give
\[
 \frac{a_{n+1}}{a_n^2}
 =\frac1{1-\theta_n}-\frac1{a_n}
   +\frac{1-\theta_{n+1}}{a_n^2}\longrightarrow1.
\]
In particular, the multipliers are strictly increasing eventually.
Thus a global orbit with these positivity and relative-error assumptions
already gives a sequence of the type in Problem~\ref{long243:res:problem}
after deletion of a finite prefix. No separate tail-limit estimate is
needed for such an orbit; without the relative-error assumption, that
estimate remains a condition of the general telescoping argument.

\begin{proposition}[scaling the numerator and denominator]\label{long243:res:scale}
For every $s,a,D,C\in\Z$,
\[
 a(sD)=s\,(aD),\qquad
 a(sC)-sD=s\,(aC-D),
\]
\[
 sD-(a-1)sC=s\,[D-(a-1)C].
\]
\end{proposition}
\leannote{Lean: \lproof{ErdosProblems/Erdos243/PaperCompleteR7/Arithmetic.lean}{31}{state_scale}.}

\noindent The three identities are formalised for the
\lword{Erdos243/ReciprocalTailRigidity.lean}{65}{nextDenState_scale}{denominator},
\lword{Erdos243/ReciprocalTailRigidity.lean}{70}{nextTailState_scale}{numerator}
and
\lword{Erdos243/ReciprocalTailRigidity.lean}{75}{centeredState_scale}{error}.
Multiplying the numerator and denominator by the same factor therefore
preserves the recurrence and scales the error by that factor. When the
factor divides all entries, we can divide it out instead. This is used
in the finite calculation of Appendix~\ref{long243:app:residue} and in
the induction in Theorem~\ref{long243:res:periodic}.

\section{When a zero error forces the Sylvester recurrence}\label{long243:sec:defect}

The converse to the Sylvester example follows by eliminating $D_n$
from two successive updates. Define
\[
 \Delta_n=a_{n+1}-(a_n^2-a_n+1),
\]
the
\lword{Erdos243/ReciprocalTailRigidity.lean}{53}{sylvesterDefect}{Sylvester
defect}. The identity below relates $\Delta_n$ to two consecutive errors.
It will show both that eventual zero error forces the recurrence and that,
under $|E_n|<C_n$, a single zero error forces the next one to vanish.

\begin{theorem}[defect identity]\label{long243:res:defect}
For all $a,a',D,C\in\Z$,
\[
 \begin{aligned}
 \relax[a'-(a^2-a+1)]\,(aC-D)
 &=a^2[D-(a-1)C]\\
 &\quad-\bigl[aD-(a'-1)(aC-D)\bigr],
 \end{aligned}
\]
that is, $\Delta_n\,C_{n+1}=a_n^{2}E_n-E_{n+1}$.
\end{theorem}
\leannote{Lean: \lproof{ErdosProblems/Erdos243/ReciprocalTailRigidity.lean}{1775}{sylvesterDefect_mul_nextTailState}.}

\begin{proof}
A direct expansion: both sides equal
$a'aC-a'D-a^{3}C+a^{2}D+a^{2}C-aD-aC+D$.
\end{proof}

\noindent The identity is formalised as the
\lword{Erdos243/ReciprocalTailRigidity.lean}{1775}{sylvesterDefect_mul_nextTailState}{defect
identity}. No sign or growth hypothesis is needed for the identity itself.
Such hypotheses enter only in its consequences below.

\begin{example}[one defect and the error it creates]\label{long243:ex:defect}
Continue Example~\ref{long243:ex:sylvester} but replace $a_3=43$ by $a_3=44$.  The
states $D_3=42$ and $C_3=1$ are unchanged, since they depend only on
$a_0,a_1,a_2$, and $E_2=0$ still.  The defect is
$\Delta_2=a_3-(a_2^2-a_2+1)=44-43=1$, and
$E_3=D_3-(a_3-1)C_3=42-43=-1$, so the identity reads
\[
 \Delta_2C_3=1\cdot1=1=49\cdot0-(-1)=a_2^{2}E_2-E_3 .
\]
By Proposition~\ref{long243:res:update} the numerator then rises, $C_4=C_3-E_3=2$.  A
single unit of defect at one index has produced a negative error of magnitude
$1$ at the next, increasing the numerator from $1$ to $2$.
\end{example}

\begin{theorem}[two vanishing errors force the step]\label{long243:res:step}
Let $a,a',D,C\in\Z$ with $aC-D\ne0$. If
\[
 D-(a-1)C=0,\qquad aD-(a'-1)(aC-D)=0,
\]
then $a'=a^2-a+1$.
\end{theorem}
\leannote{Lean: \lproof{ErdosProblems/Erdos243/ReciprocalTailRigidity.lean}{1787}{sylvesterNext_eq_of_centered_zero}.}

\begin{proof}
Theorem~\ref{long243:res:defect} gives
$[a'-(a^2-a+1)]\,(aC-D)=0$. The second factor is nonzero, so the first
factor must vanish.
\end{proof}

\noindent The conclusion is formalised as the
\lword{Erdos243/ReciprocalTailRigidity.lean}{1787}{sylvesterNext_eq_of_centered_zero}{local
rigidity step}.  The hypothesis $C_{n+1}\ne0$ is not removable: at $C_{n+1}=0$
the identity gives no information about $a'$.

\subsection{Relations among three consecutive numerators}

Eliminating the denominator from two successive steps gives a relation
among three reduced numerators.  Write $u,u_1,u_2$ for the
reduced numerator coordinates, $v,v_1$ for the corresponding denominator
coordinates, and $h,h_1$ for the two common factors removed in reduction.  If
the two steps have multipliers $a,a_1$, then the exact hypotheses are
\[
 hu_1+v=au,\qquad hv_1=av,\qquad
 h_1u_2+v_1=a_1u_1.
\]

\begin{theorem}[eliminating the denominator from two steps]\label{long243:res:secondorder}
Let $a,a_1,u,u_1,u_2,v,v_1,h,h_1$ be integers satisfying
\[
 hu_1+v=au,\qquad hv_1=av,\qquad h_1u_2+v_1=a_1u_1.
\]
Then $a^2u+hh_1u_2=h(a+a_1)u_1$.
\end{theorem}
\leannote{Lean: \lproof{ErdosProblems/Erdos243/PaperCompleteR7/Arithmetic.lean}{42}{reduced_second_order_int}.}

\begin{proof}
Multiply the first equation by $a$, replace $av$ using the second equation,
and then replace $h_1u_2+v_1$ using the third.  The remaining terms factor as
the displayed right-hand side.
\end{proof}

\noindent This is the
\lword{Erdos243/DynamicCancellation.lean}{270}{reducedStep_secondOrder}{recurrence obtained by eliminating the denominator}.  Both cancellation factors remain in the formula. The identity adds no
assumption to the two exact steps, but it still contains both multipliers
and cannot by itself force $a_1=a^2-a+1$;
Lemma~\ref{long243:res:oldmodulussaturation} specifies what it forgets
modulo a fixed old denominator.

When neither step cancels a common factor, the same recurrence gives
a square identity. Write $p,p_1,p_2$ for three consecutive numerators,
$q=ap-p_1$ for the denominator before the first step, and $a,a_1$ for
the two multipliers. Then
\[
 p_2+a^2p=(a+a_1)p_1.
\]
The resulting square identity is as follows.

\begin{theorem}[a square identity without cancellation]\label{long243:res:curvature}
Let $a,a_1,p,p_1,p_2,q$ be integers with $q=ap-p_1$ and
$p_2+a^2p=(a+a_1)p_1$.  Then
\[
 q^2+\bigl(pp_2-p_1^2\bigr)=(a_1-a)pp_1.
\]
\end{theorem}
\leannote{Lean: \lproof{ErdosProblems/Erdos243/DynamicCancellation.lean}{309}{cancellationFree_curvature_square}.}

\begin{proof}
Substitute $q=ap-p_1$ and $p_2=(a+a_1)p_1-a^2p$:
\[
 \begin{aligned}
 q^2+pp_2-p_1^2
 &=(ap-p_1)^2+p\bigl((a+a_1)p_1-a^2p\bigr)-p_1^2\\
 &=(a_1-a)pp_1.
 \end{aligned}
\]
\end{proof}

The term $q^2$ is nonnegative, while $pp_2-p_1^2$ measures the difference
between the product of the outer numerators and the square of the middle one.  The right side records the change in multiplier, so the
identity gives a necessary algebraic constraint without
asserting a sign or a global monotonicity.  It is checked as
\lword{Erdos243/DynamicCancellation.lean}{309}{cancellationFree_curvature_square}{the
square identity when no factor is cancelled}; no global exclusion of approximate solutions is
claimed.

The reduced error update must also be compatible with the denominator
recurrence. The following calculation gives the exact compatibility
condition.
Here $e,e_1$ are signed errors in reduced coordinates, not the positive
magnitudes denoted by $e_n$ in the all-negative case. Let $\Delta$ be a
proposed value of the Sylvester defect. If
\[
 hu_1=u-e,
 \qquad he_1=a^2e-\Delta(u-e),
\]
then the denominator equality
\[
 h\bigl((a_1-1)u_1+e_1\bigr)
   =a\bigl((a-1)u+e\bigr)
\]
holds precisely when the following difference vanishes:
\[
 \begin{aligned}
 &h\bigl((a_1-1)u_1+e_1\bigr)-a\bigl((a-1)u+e\bigr)\\
 &\qquad=\bigl[a_1-(a^2-a+1)-\Delta\bigr]\,(u-e).
 \end{aligned}
\]
Consequently, when $u-e\ne0$,
\[
 h\bigl((a_1-1)u_1+e_1\bigr)
   =a\bigl((a-1)u+e\bigr)
 \quad\Longleftrightarrow\quad
 a_1=a^2-a+1+\Delta.
\]
The
\lword{Erdos243/FeedbackRealizability.lean}{64}{feedback_denominator_transport_iff}{equivalence with the denominator recurrence}
is formalised using the
\lword{Erdos243/FeedbackRealizability.lean}{15}{feedback_transport_defect}{factored difference between the two recurrences}.  The hypothesis $u-e\ne0$ is essential: if
$u-e=0$, the factorised mismatch cannot identify $a_1$.

The identities hold under the local hypotheses displayed in their
statements. It is their proposed global applications, not the identities,
that would need additional information, such as control of cancellation or
a sign estimate for $pp_2-p_1^2$. No such information is deduced here for
an arbitrary rational near-quadratic tail.

\begin{lemma}[saturation modulo an old denominator]\label{long243:res:oldmodulussaturation}
Let $M\ge2$. Any finite word $r_0,\ldots,r_k$ of units modulo $M$ is
compatible with the cancellation-free recurrences modulo $M$, with all
reduced denominator residues equal to zero. Consequently the eliminated
two-step identity alone imposes no further restriction on such unit words
when the multiplier residues are free.
\end{lemma}
\leannote{Lean: \lproof{ErdosProblems/Erdos243/PaperCompleteR21/ReducedStepLocalArithmetic.lean}{39}{unit_word_saturates_old_modulus}.}
\begin{proof}
Set $v_i\equiv0\pmod M$ and choose
$a_i\equiv r_{i+1}r_i^{-1}\pmod M$. Then
$r_{i+1}+v_i\equiv a_ir_i$ and $v_{i+1}\equiv a_iv_i$.
Eliminating $v_i$ gives the two-step identity automatically.
\end{proof}
This is finite residue compatibility, not existence of a positive integer
orbit, much less a rational tail with near-quadratic growth.
For $M\mid v_T$ on a tail with no further common-factor cancellation,
the actual $u_n$ are indeed units modulo
$M$. Useful further restrictions must therefore retain information discarded
here: for example multiplier size, changing moduli, or primes outside the
old denominator. The square restriction used in the cubic proof concerns
precisely a prime at which a middle numerator vanishes.

\begin{theorem}[sequence form]\label{long243:res:eventual}
Let $a,D,C:\N\to\Z$ satisfy $D_{n+1}=a_nD_n$ and
$C_{n+1}=a_nC_n-D_n$.  If $E_n=0$ for all sufficiently large $n$ and
$C_{n+1}\ne0$ for all sufficiently large $n$, then $a_{n+1}=a_n^{2}-a_n+1$ for
all sufficiently large $n$.
\end{theorem}
\leannote{Lean: \lproof{ErdosProblems/Erdos243/ReciprocalTailRigidity.lean}{1805}{sylvesterNext_eventually_of_centered_zero}.}

\noindent The sequence form is formalised as the
\lword{Erdos243/ReciprocalTailRigidity.lean}{1805}{sylvesterNext_eventually_of_centered_zero}{eventual
Sylvester recurrence}.  Take the larger of the two thresholds and apply
Theorem~\ref{long243:res:step} at each later index. This is the final
step of the criteria below that force $E_n$ to vanish eventually.

The defect identity has a second consequence, which is what makes a single
vanishing error worth having.

\begin{theorem}[zero is absorbing]\label{long243:res:absorb}
Let $a,C,D:\N\to\N$ be an exact orbit of natural numbers, so
$C_{n+1}+D_n=a_nC_n$ and $D_{n+1}=a_nD_n$, and let
$E_n=D_n-(a_n-1)C_n$.  Suppose the centring is strict, $|E_n|<C_n$ for
every $n$.  If $E_n=0$ then $E_{n+1}=0$.
\end{theorem}
\leannote{Lean: \lproof{ErdosProblems/Erdos243/ReciprocalTailRigidity.lean}{2227}{centeredState_zero_absorbing}.}

\begin{proof}
Putting $E_n=0$ in Theorem~\ref{long243:res:defect} gives
$\Delta_n C_{n+1}=-E_{n+1}$, so $C_{n+1}$ divides $E_{n+1}$.  A multiple of
$C_{n+1}$ of absolute value smaller than $C_{n+1}$ is zero, and
$|E_{n+1}|<C_{n+1}$ by hypothesis.
\end{proof}

\noindent The conclusion is formalised as the
\lword{Erdos243/ReciprocalTailRigidity.lean}{2227}{centeredState_zero_absorbing}{absorption
of a vanishing error}, with the companion
\lword{Erdos243/ReciprocalTailRigidity.lean}{2254}{eventually_nonzero_of_zero_absorbing}{contrapositive
along a tail}: if zero is absorbing beyond some index and $E$ does not vanish
eventually, then $E$ is nowhere zero beyond that index.

Beyond an eventual centring threshold, either $E$ reaches zero and
stays there, or it is nowhere zero. A zero before that threshold need not
persist, as Example~\ref{long243:ex:defect} shows: the altered multiplier
creates $E_3=-1$ with $C_3=1$, exactly where strict centring fails.
Sections~\ref{long243:sec:constant} and~\ref{long243:sec:periodic} address
the special all-negative case. The mixed-sign analysis in
Sections~\ref{long243:sec:bounded} and~\ref{long243:sec:mass} separates
arbitrarily late negative errors from an eventually nonnegative tail.

\section{A criterion using new maxima of an LCM numerator}\label{long243:sec:lcmrecords}

Clear each rational tail with a least common multiple rather than a
product. The resulting numerator is a positive integer. We will show that
the Sylvester recurrence is equivalent to convergence of a weighted sum
over steps at which this numerator exceeds all its previous values.
The first-crossing proof permits a fixed amount to be subtracted from
each such increase. The equivalence is not a proof of convergence under
the original hypotheses.

\paragraph{Formalisation.}
The proof below is ordinary mathematics. The first-crossing arithmetic has
formalised components in \texttt{LcmRecordExcess.lean}. In addition, the
separate release contains
\href{https://github.com/wcook04/plectis-erdos-lean/blob/52f29ad173b04e3bac941b3663f2b9aebe5de0bb/ErdosProblems/Erdos243/PaperCompleteR8/CanonicalWeightedRecords.lean\#L268}{the complete criterion for real-valued weights} and its natural-weight version. These release sources contain proof
bodies, but they are outside the supplied successful main-repository build;
a Comparator configuration alone is not a replay receipt. The finite-orbit
data retained with the source are diagnostic examples, not a proof of
termination for all rational seeds.

The sum counts only steps setting new maxima and subtracts the same
fixed amount from each actual jump. Unlike Koizumi's eventual-nonnegative
case in Proposition~1(2), it permits both signs of the error.
The weaker model in Proposition~\ref{long243:res:coprimalitycap} assumes
only nondivisibility by whole moduli. It omits both coprimality to their
prime factors and the denominator recurrence, so its counterexamples
do not refute an argument using those additional properties.

In this section the indices start at $0$. Write the full reciprocal
sum as $p/q$, with positive integers $p,q$, put $x_n=\sum_{k\ge n}1/a_k$, and take
$D_0=q$. Set
\[
 L_n=\operatorname{lcm}(q,a_0,\ldots,a_{n-1}),\quad
 M_n=D_n/L_n,\quad U_n=C_n/M_n,\quad V_n=E_n/M_n.
\]
The rational tail has denominator dividing $L_n$, so $U_n=L_nx_n$
and $V_n$ are integers. This clears the denominator but need not reduce
the fraction: $U_n$ and $L_n$ may still have common factors.
With $\rho_n=\gcd(L_n,a_n)$, the exact updates are
\[
 M_{n+1}=M_n\rho_n,\qquad
 \rho_nU_{n+1}=U_n-V_n,\qquad V_n=L_n-(a_n-1)U_n.
\]
These are Bado's LCM coordinates, with the indexing translated explicitly.
Take his denominator parameter to be $q$ and identify his $a_{n+1}$
with our $a_n$. Then his $M_n$, $\Delta_n$, $K_n$, $u_{n+1}$ and
$g_{n+1}$ are respectively our $L_n$, $M_n$, $U_n$, $V_n$ and $\rho_n$.
His (19) becomes the middle update above
\cite[Prop.~7.1 and (16)--(20), pp.~6--7]{bado2026}.

Write $R_n=\max_{j\le n}U_j$ and
$\mathcal R=\{n:U_{n+1}>R_n\}$.  Strict centring already gives
$U_{n+1}<U_n$ when $\rho_n\ge2$, so every sufficiently late strict rise has
$\rho_n=1$.  The stronger eventual bound $-U_n\le2V_n$, supplied by
$V_n/U_n=E_n/C_n\to0$, gives the quantitative estimate
$U_{n+1}\le3U_n/4$ when $\rho_n\ge2$.
At a sufficiently late LCM record, $\rho_n=1$ and the actual jump is
$d_n=U_{n+1}-U_n=-V_n>0$. At a contracting step with $\rho_n\ge2$ this
identity need not hold.

\begin{theorem}[boundedness and a weighted sum over new maxima]\label{long243:res:arithmeticrecord}
Let $a_n,L_n,U_n$ be positive integers and $V_n$ integers satisfying
\[
 \begin{gathered}
 L_{n+1}=\operatorname{lcm}(L_n,a_n),\qquad \rho_n=\gcd(L_n,a_n),\\
 \rho_nU_{n+1}=U_n-V_n,\qquad
 V_n=L_n-(a_n-1)U_n,\qquad -U_n\le2V_n.
 \end{gathered}
\]
Let $f:[1,\infty)\to[0,\infty)$ be finite and nonincreasing, with
$\int_1^\infty f(t)\,dt=\infty$. For each fixed integer $B\ge0$,
\[
 \sup_n U_n<\infty
 \quad\Longleftrightarrow\quad
 \sum_{n\in\mathcal R}(-V_n-B)_+f(U_n)<\infty.
\]
\end{theorem}
\leannote{Lean: \lproof{ErdosProblems/Erdos243/PaperCompleteR11/ArithmeticWeightedRecord.lean}{264}{arithmetic_weighted_record_dichotomy}.}

\begin{proof}
A bounded integer running maximum increases only finitely many times, so
bounded $U_n$ gives a finite sum. Suppose instead that $U_n$ is unbounded.
There are infinitely many record rises, each with $\rho_n=1$. Its multiplier is
greater than one because $d_n=(a_n-1)U_n-L_n>0$. If $r<s$ are record
indices, then $a_r\mid L_s$ and $\gcd(a_s,L_s)=1$, hence
$\gcd(a_r,a_s)=1$. The record multipliers are therefore distinct and pairwise
coprime. For any fixed $B\ge1$, choose $B$ of them greater than $B$,
and take $T$ after their indices. These earlier multipliers
$m_0,\ldots,m_{B-1}$ all divide $L_T$.

Put $P=\prod_i m_i$ and choose $x$ by the Chinese remainder theorem with
$m_i\mid x+i$. Consider the heights $\tau=x+B+kP>R_T$, $k\in\Z$.
We first show that a crossing requires $d_n>B$, then count how many
heights one jump can cross. A first crossing
$U_n\le R_n<\tau\le U_n+d_n$ is a record step. If $d_n\le B$, then
$U_n\in[\tau-B,\tau)$, so some $m_i$ divides $U_n$. It also divides
$L_n$, hence divides $d_n=(a_n-1)U_n-L_n$, contradicting
$0<d_n\le B<m_i$.

If the step first crosses $h\ge1$ such heights, their spacing gives
$(h-1)P<d_n$. With $r=d_n-B\ge1$ and $P\ge B+1$, we have
$d_n=B+r\le Pr$, whence $h\le r$. Monotonicity of $f$ now gives
\[
 \sum_{\substack{\tau\text{ first crossed}\\\text{at step }n}}f(\tau)
 \le(d_n-B)f(U_n).
\]
Each selected height has one first crossing, and the first selected
height is at most $R_T+P$. Summing and comparing each interval of length
$P$ with its left endpoint gives, for $R_N\ge R_T+P$, the finite bound
\begin{equation}\label{long243:eq:weightedcrossing}
 \sum_{\substack{T\le n<N\\n\in\mathcal R}}(-V_n-B)_+f(U_n)
 \ge\frac1P\int_{R_T+P}^{R_N} f(t)\,dt.
\end{equation}
Since $R_n\to\infty$, the right side diverges for every $B\ge1$.
The case $B=0$ follows by domination.
\end{proof}

Only the lower centring bound was used. In particular, neither vanishing relative error nor a growth estimate for $C_n$ supplies the CRT moduli: the
record multipliers themselves do so. Vanishing relative error enters in the following
application to the original sequence.

\begin{theorem}[a convergent weighted sum over new maxima]\label{long243:res:weightedrecord}
Assume the growth and rationality hypotheses of Problem~\ref{long243:res:problem},
and let $f$ be as in Theorem~\ref{long243:res:arithmeticrecord}. The sequence is
eventually Sylvester if and only if, for some integer $B\ge0$,
\[
 \sum_{n\in\mathcal R}(-V_n-B)_+f(U_n)<\infty.
\]
\end{theorem}
\leannote{Lean: \lproof{ErdosProblems/Erdos243/PaperCompleteR11/CanonicalRecords.lean}{202}{canonical_weighted_record_excess}.}

\begin{proof}
If the sequence is not eventually Sylvester, zero error is never reached on a sufficiently late
tail. Integrality gives $1/U_n\le|V_n|/U_n=|E_n|/C_n\to0$, so $U_n$ is
unbounded. Theorem~\ref{long243:res:arithmeticrecord}, applied after the centring
threshold, forces divergence for every $B$. Once the running maximum of
the retained tail exceeds the omitted prefix maximum, deleting that
prefix no longer changes which steps set new records. Unboundedness
ensures that this crossing occurs. Conversely, a Sylvester tail
telescopes to $x_n=1/(a_n-1)$, so $V_n=0$ eventually.
\end{proof}

The weights $f(t)=1/t$ and $f(t)=1/[t\log(et)]$ are admissible: they are
nonnegative and nonincreasing, but their integrals diverge. The faster-decaying
weight $1/t^2$ is not admissible. The divergence requirement is used at the
last line of the crossing argument; without it, unbounded numerators need
not make the lower bound diverge. For example, $f(t)=1/[t\log(et)]$ gives
the sufficient condition
\[
 \sum_{n\in\mathcal R}\frac{(-V_n-B)_+}{U_n\log(eU_n)}<\infty.
\]
Its finite lower bound in~\eqref{long243:eq:weightedcrossing} is
$P^{-1}\log\bigl(\log(eR_N)/\log(e(R_T+P))\bigr)$.
Further fixed iterated logarithmic factors are allowed whenever the integral
still diverges. These are specialisations of one crossing theorem.

The criterion also has an exact expression in the original growth defect.
Put $\gamma_n=a_n^2/a_{n+1}-1$ and $\theta_n=E_n/C_n$. The defect identity
gives
\[
 \gamma_n+\theta_n=
 \frac{(1-\theta_n)(a_n-1+\theta_{n+1})}{a_{n+1}},
 \qquad 0<\gamma_n+\theta_n<3/a_n
\]
eventually. Thus the two nonnegative summands
$U_nf(U_n)(\gamma_n-B/U_n)_+$ and $(-V_n-B)_+f(U_n)$ differ by at most
$3U_nf(U_n)/a_n$. To see summability directly, put $P_n=\prod_{j<n}a_j$.
The tail estimate $x_n\sim1/a_n$ gives $C_n/a_n\sim qP_n/a_n^2$, and
\[
 \frac{P_{n+1}/a_{n+1}^{2}}{P_n/a_n^{2}}
 =\frac{a_n^3}{a_{n+1}^{2}}\longrightarrow0.
\]
Thus $\sum_n C_n/a_n<\infty$ by the ratio test. Since $U_n\le C_n$ and
$f(U_n)\le f(1)$, the comparison error is summable.
Consequently Theorem~\ref{long243:res:weightedrecord} is equivalent to finiteness of
\begin{equation}\label{long243:eq:weightedgrowth}
 \sum_{n\in\mathcal R}U_nf(U_n)
 \left(\frac{a_n^2}{a_{n+1}}-1-\frac B{U_n}\right)_+
\end{equation}
for some $B$. The original hypotheses do not currently supply this
finiteness. In particular, termwise convergence to zero is insufficient.
The corresponding formal comparison starts from the same positive,
strictly increasing integer sequence, its rational reciprocal sum and
the near-quadratic growth limit. It uses the summand in
\eqref{long243:eq:weightedgrowth} at record indices and zero elsewhere,
with the same class of weights and existential choice of $B$.
The formal estimate uses the same ratio-test argument, with the sufficient
constant $16$ in place of $3$. This transfers the criterion; it does not
supply the required finiteness from the original hypotheses.
The separate-release source
\href{https://github.com/wcook04/plectis-erdos-lean/blob/52f29ad173b04e3bac941b3663f2b9aebe5de0bb/ErdosProblems/Erdos243/PaperCompleteR8/GrowthDebtSummability.lean\#L234}{\texttt{GrowthDebtSummability.lean}}
contains the factored-growth formulation. It is outside the documented
main-repository build. The on-page comparison above is an ordinary proof;
this prose revision does not constitute a new Lean check of that release.

The crossed heights lie above $R_n$, but the divisibility argument uses
the starting numerator $U_n$ and the full jump $U_{n+1}-U_n$. Replacing
this by $R_{n+1}-R_n$ would discard the recovery from an earlier decrease.


\paragraph{Integer coefficients.}
The following extension is not used in either the bounded-increment proof
or the cubic-rate argument. We give the proof here to keep the short note's
comparison separate from its main unit-fraction argument.

The first-crossing argument also works when the summand has an integer
numerator $b_n$. In that case
\[
 V_n=b_nL_n-(a_n-1)U_n,\qquad
 \rho_nU_{n+1}=U_n-V_n.
\]
Assume that $a_n\ge2$, that $L_n,U_n$ are positive integers, and that
$L_{n+1}=\operatorname{lcm}(L_n,a_n)$, with
$\rho_n=\gcd(L_n,a_n)$. Suppose $V_n\ge-B$ eventually, with an integer
$B\ge0$. Then $U_{n+1}\le U_n+B$; if $B=0$, boundedness follows at once.
For $B\ge1$, a record step with $R_n>B$ must have $\rho_n=1$, since
$\rho_n\ge2$ would give
\[
 U_{n+1}\le\frac{U_n+B}{2}\le\frac{R_n+B}{2}<R_n.
\]
Thus an unbounded sequence would supply infinitely many pairwise coprime
record multipliers. Choose $B$ of them larger than $B$, and then a CRT
block of $B$ consecutive integers above the previous maximum, each
divisible by one of the chosen multipliers. The first step past the upper
end of the block is a high record step, so $\rho_n=1$.
Its positive jump is at most $B$, hence the preceding numerator $U_n$
lies in the block. The corresponding multiplier divides both $U_n$ and
$L_n$, hence divides the jump $(a_n-1)U_n-b_nL_n$. A positive jump at most $B$ cannot have a divisor
larger than $B$. This also explains why no growth assumption on
$a_n$, centring, or relative-error limit is needed for boundedness.

With the additional limit $V_n/U_n\to0$, choose an integer $K$ bounding
$U_n$. Eventually $|V_n|<U_n/K\le1$, so the integer $V_n$ is zero.
The recurrence becomes $\rho_nU_{n+1}=U_n$; the positive integer
sequence $U_n$ is then nonincreasing and hence eventually constant. For positive $b_n$, the classical comparisons
are Badea's Corollary~2.2 \cite[p.~316]{badea1993} and the criterion of
Tijdeman and Yuan \cite{tijdemanyuan2002}. Further finite examples
separating boundedness from stationarity are in the
\href{https://github.com/wcook04/plectis-erdos/blob/eccd8afc6db3c02a2265d0601b727d6c5e4467d5/lean/ErdosProblems/Erdos243/CoefficientUniformBoundedHeight.md}
{coefficient proof supplement}.

\section{New maxima of reduced numerators}\label{long243:sec:records}

In Section~\ref{long243:sec:lcmrecords}, a prime divisor of $L_n$
continues to divide every later $L_j$. This section instead writes the
reciprocal tail in lowest terms as $u_n/v_n$. Cancellation can then remove
prime factors from $v_n$, so their persistence needs a separate proof.
The criteria below use the size of new maxima of $u_n$ and the prime powers
that remain in its denominator. They do not identify $u_n$ with the
LCM numerator $U_n$.

\paragraph{Fractions in lowest terms.}
Put
\[
 G_n=\gcd(C_n,D_n),\qquad u_n=C_n/G_n,\qquad v_n=D_n/G_n,
 \qquad \tilde e_n=E_n/G_n.
\]
Thus $\gcd(u_n,v_n)=1$ and $\tilde e_n$ is a signed integer.
The factor $h_n=G_{n+1}/G_n$ records the common factor removed at the next
step. The relation to the preceding section is exact:
\[
 \frac{G_n}{M_n}=\gcd(U_n,L_n),\qquad
 u_n=\frac{U_n}{\gcd(U_n,L_n)},\qquad
 v_n=\frac{L_n}{\gcd(U_n,L_n)}.
\]
Indeed, $C_n=M_nU_n$ and $D_n=M_nL_n$.
Clearing with an LCM and reducing to lowest terms are different operations.
For example, the exact step $(C,D,a)=(3,6,3)$ has $E=0$ and gives
$(C',D')=(3,18)$. With $L=6$, its LCM numerator falls from $3$ to $1$,
while the reduced numerator stays $1$: the LCM factor is $\rho=3$,
but the reduction factor is $h=1$. This step begins a Sylvester tail.

In general, before reduction the next numerator is
$w_n=a_nu_n-v_n=u_n-\tilde e_n$, so
\[
 h_nu_{n+1}=w_n,\qquad h_nv_{n+1}=a_nv_n,\qquad
 \gcd(u_n,u_{n+1})=1.
\]
The negative magnitude $e_n=-E_n$ used earlier is not $\tilde e_n$;
neither is Koizumi's real gap $\varepsilon_n$.

For the maxima and their increments write
\[
 R_n=\max_{k\le n}u_k,\qquad H_n=\max_{j\le n}C_j,
 \qquad s_n=R_{n+1}-R_n.
\]
At a strict rise let $d_n=u_{n+1}-u_n$. Then
$s_n=(d_n-(R_n-u_n))_+$. Thus the actual jump includes recovery of
the earlier decrease $R_n-u_n$, whereas the record increment does not. For example,
if $R_n=10$, $u_n=5$ and $u_{n+1}=12$, then $s_n=2$ but $d_n=7$.
This illustrates the definitions, not a claimed reciprocal-tail orbit.
We will use
\[
 m_n=(-\tilde e_n)_+,\qquad \mathcal A_n=\frac{R_n}{u_n}m_n,
 \qquad \delta_n=\left(\frac{a_n^2}{a_{n+1}}-1\right)_+,
 \qquad \ell(x)=\log_2\log_2\max(4,x).
\]
The factor $R_n/u_n$ in $\mathcal A_n$ measures how far the current
numerator has fallen below its previous maximum. It equals one at a
maximum and can be large after a decrease. The symbol $\mathcal A_n$
is distinct from the prefix product $A_n$ used later.

Unless a statement in this section specifies an abstract recurrence,
we assume the hypotheses of Problem~\ref{long243:res:problem} and use
its integer tails, for which $|E_n|/C_n\to0$ after a finite shift.
We will write out the conclusion $a_{n+1}=a_n^2-a_n+1$ eventually
rather than give it a separate symbol.

\paragraph{Comparison with Duverney's reduced parameters.}
For the positive reciprocal series, the eventually reduced parameters in Duverney's
Theorem~3.1 and Section~5.2 \cite[pp.~285--286, 299--300]{duverney2001}
can be taken as follows. The symbols $p_n,q_n$ in this comparison are
Duverney's auxiliary integers, not the fixed numerator and denominator
of the full reciprocal sum:
\[
 p_n=u_n,\qquad q_n=w_n=a_nu_n-v_n=h_nu_{n+1},\qquad
 E_n=G_n(p_n-q_n).
\]
Indeed, $\gcd(u_n,w_n)=\gcd(u_n,v_n)=1$. Since
$a_nv_n=a_n^2u_n-a_nw_n$, reduction of the next tail gives
\[
 h_n=\gcd(w_n,a_nv_n)=\gcd(w_n,a_n^2),\qquad
 p_{n+1}\mid q_n.
\]
For the gcd equality, subtract the multiple $a_nw_n$ and use the
coprimality of $u_n,w_n$; the divisibility follows from
$q_n=h_np_{n+1}$. Eliminating $v_{n+1}$ then gives Duverney's recurrence:
\[
 a_{n+1}=\frac{p_n}{q_n}a_n^2-a_n+\frac{q_{n+1}}{p_{n+1}}.
\]
Thus $q_n/p_{n+1}=h_n$ is precisely the cancellation factor, and
$q_n/p_n\to1$. This identifies the reduced parameters, not every possible
unreduced choice in the original theorem. An upper bound on $q_n-p_n$ controls the increase before cancellation
in the reduced coordinates. To transfer that bound directly to
$(-E_n)_+=G_n(q_n-p_n)_+$ would also require control of $G_n$.

The source references distinguish complete Lean proofs from written
arguments that use formalised lemmas. The supplied index records a build
of the listed public modules at revision \texttt{6b78209ab63a}; the links
below retain their original revisions. It separately records successful
Comparator checks for specified theorems in the release, not a build of
that release as a whole. A challenge statement containing \texttt{sorry}
is not a proof, and a configuration entry alone does not record a successful
check. Citing a formalised lemma does not discharge the extra hypotheses
or steps of the written argument using it.

\begin{lemma}[the denominator valuation transition]\label{long243:res:valuationtransition}
Let $u,v,a$ be positive integers, $\gcd(u,v)=1$, and $w=au-v>0$.
Put $h=\gcd(w,av)$ and $v'=av/h$. For a prime $p$, write
$r=\nu_p(a)$, $s=\nu_p(v)$ and $t=\nu_p(w)$. Then
\[
 \nu_p(v')=
 \begin{cases}
 \max(r,s),&r\ne s,\\
 \max(0,2s-t),&r=s.
 \end{cases}
\]
In particular $\nu_p(v')\le\max(r,s)$. A strict loss relative to $s$
requires $r=s\ge1$ and $t>s$.
\end{lemma}
\leannote{Lean: \lproof{ErdosProblems/Erdos243/PaperCompleteR21/ReducedStepLocalArithmetic.lean}{244}{reduced_denominator_valuation_transition}, \lproof{ErdosProblems/Erdos243/PaperCompleteR21/ReducedStepLocalArithmetic.lean}{262}{reduced_denominator_valuation_le_max}, \lproof{ErdosProblems/Erdos243/PaperCompleteR21/ReducedStepLocalArithmetic.lean}{279}{reduced_denominator_valuation_strict_loss}.}
\begin{proof}
Always $\nu_p(v')=r+s-\min(t,r+s)$. If $r\ne s$, primitivity implies
$t=\min(r,s)$: when $s>0$, $u$ is a $p$-adic unit, and when $s=0<r$,
$v$ is a unit. If $r=s>0$, both terms in $au-v$ are divisible by $p^s$,
so $t\ge s$; substituting gives the second case. If $r=s=0$, the same
formula gives zero without needing $u$ to be a unit.
\end{proof}
Cancellation is not the same as a fall in the reduced denominator valuation.
For example, $(u,v,a)=(2,15,9)$ gives $w=3$, $h=3$ and
$(u',v')=(1,45)$. The numerator before cancellation exceeds $u$ by only
$w-u=1$, yet the $3$-adic valuation of the reduced denominator rises
from $1$ to $2$.
This is a finite exact step, not an infinite counterexample.

\begin{corollary}[persistence of a prime power]\label{long243:res:powerpersistence}
Suppose $p^k\mid v_s$. If $w_n<p^{k+1}$ at every step from $s$ through
$t-1$, then $p^k\mid v_t$.
\end{corollary}
\leannote{Lean: \lproof{ErdosProblems/Erdos243/PaperCompleteR20/PowerPersistence.lean}{27}{primePower_persists}.}
\begin{proof}
At a first loss of divisibility by $p^k$, the current exponent is some
$j\ge k$ and the valuation lemma forces $\nu_p(w_n)>j$. This would give
$w_n\ge p^{j+1}\ge p^{k+1}$, contrary to the hypothesis.
\end{proof}
These are ordinary local calculations. They explain the persistence
threshold used below without identifying any new Lean declaration.

\subsection{How fast the running maximum must increase}

\begin{theorem}[increments of the running maximum]\label{long243:res:recorddichotomy}
Let $\Theta=\limsup_n(H_{n+1}-H_n)/\ell(H_n)$ on a rational-tail orbit
under the standing hypotheses, and suppose $E_n$ is not eventually zero.
Then either $G_n$ is unbounded and $\Theta$ is infinite, or
$G_n$ stabilises at a value $g$ and $\Theta\ge g\,v_T/\varphi(v_T)$ for every
late $T$, so that $\Theta>g\ge1$.  For any orbit under the standing
hypotheses, therefore, $\Theta=0$ or $\Theta>1$, and $\Theta\le1$ forces
the eventual Sylvester recurrence.
\end{theorem}
\leannote{Lean: \lproof{ErdosProblems/Erdos243/PaperCompleteR11/QuantitativeRecordDichotomy.lean}{340}{canonical_quantitative_record_dichotomy}, \lproof{ErdosProblems/Erdos243/PaperCompleteR11/QuantitativeRecordDichotomy.lean}{308}{canonical_recordTheta_zero_or_gt_one}, \lproof{ErdosProblems/Erdos243/PaperCompleteR11/QuantitativeRecordDichotomy.lean}{321}{canonical_recordTheta_eq_zero_iff}, \lproof{ErdosProblems/Erdos243/PaperCompleteR11/InclusiveLimsup.lean}{54}{canonical_recordTheta_gt_one}.}

The constant in the next condition is measured against a double
logarithm, not against $C_n$. This grows much more slowly than any
positive power of $C_n$. Every fixed bound on $(-E_n)_+$ satisfies
it on a nonzero tail because $C_n\to\infty$; it also permits unbounded
negative parts on the double-logarithmic scale, with the coefficient
specified below. In contrast, an error of size
$\sqrt{C_n}$ has relative size tending to zero but violates the condition.
These are comparisons of bounds, not constructions of reciprocal tails.

\begin{corollary}[the double-logarithmic bound]\label{long243:res:loglogboundary}
If
\[
 \limsup_n\frac{(-E_n)_+}{\ell(C_n)}\le1,
\]
then the sequence is eventually Sylvester. Every counterexample therefore
satisfies
\[
 \limsup_n\frac{(-E_n)_+}{\ell(C_n)}>1.
\]
\end{corollary}
\leannote{Lean: \lproof{ErdosProblems/Erdos243/PaperCompleteR11/InclusiveLimsup.lean}{129}{canonical_inclusive_logLog_criterion}, \lproof{ErdosProblems/Erdos243/PaperCompleteR11/InclusiveLimsup.lean}{108}{canonical_negativeError_limsup_gt_one}.}

\begin{proof}[Proof of Theorem~\ref{long243:res:recorddichotomy} and
Corollary~\ref{long243:res:loglogboundary}]
We first prove the auxiliary lower bound $\Theta\ge1$. This part needs
only an exact positive integer orbit with $a_n>1$, $D_0\ge1$,
$|E_n|/C_n\to0$ and error not eventually zero. In particular, it will also
apply to the abstract orbit in Theorem~\ref{long243:res:slownegative}.
The key input is that, outside a set of indices of density zero,
$a_n$ is coprime to the entire preceding denominator $D_n$. This supplies
almost one new coprime modulus per step. The denominator growth then
places the resulting CRT block at the required double-logarithmic height.

\emph{How often a multiplier is coprime to the preceding denominator.}
We use the LCM factorisation from
Section~\ref{long243:sec:lcmrecords}. In the present zero-based indexing, put
\[
 L_n=\lcm(D_0,a_0,\ldots,a_{n-1}),\qquad
 M_n=D_n/L_n,\qquad \rho_n=\gcd(L_n,a_n).
\]
The finite telescoping identity gives
\[
 \frac{C_n}{M_n}
 =L_n\left(\frac{C_0}{D_0}-\sum_{j<n}\frac1{a_j}\right)
 \in\Z_{>0}.
\]
Thus $M_n\mid C_n$, while $M_0=1$ and
$M_{n+1}=\rho_nM_n$. Since $L_n$ and $D_n$ have the same prime divisors,
\[
 \#\{n<N:\gcd(a_n,D_n)>1\}
 =\#\{n<N:\rho_n>1\}
 \le\log_2M_N\le\log_2C_N=o(N).
\]
Here the last estimate follows by summing
$\log(C_{n+1}/C_n)=\log(1-E_n/C_n)=o(1)$.
This is the density-one coprimality argument of
Bado~\cite[Theorem~9.1, p.~7]{bado2026}, written with the exact-orbit
hypotheses used here. Only the finite telescoping identity is needed;
no estimate on record increments has entered this count.

\emph{Choosing the moduli and controlling their product.}
Absorption and vanishing relative error give $C_n\to\infty$, hence
$H_n\to\infty$. The maximum is still subexponential: for every
$\varepsilon>0$, choose $K_\varepsilon$ with
$C_j\le K_\varepsilon e^{\varepsilon j}$ for every $j$. Then
$H_n\le K_\varepsilon e^{\varepsilon n}$, so $\log H_n=o(n)$.
Since $D_n\ge2^n$ and
$a_n=D_n/C_n+1-E_n/C_n$, eventually
\[
 a_n\ge2^{n/2},\qquad a_n<2D_n,\qquad \ell(D_n)\le n+O(1).
\]
For the last inequality, $D_{n+1}<2D_n^2$ gives
$1+\log_2D_{n+1}<2(1+\log_2D_n)$. Iteration from a fixed late index,
followed by another logarithm, gives $\ell(D_n)\le n+O(1)$.

Suppose $H_{n+1}-H_n\le c\ell(H_n)$ eventually for some $0<c<1$.
For a large integer $B$, take the first $B$ indices at or after
$\lceil3\log_2B\rceil$ for which $\gcd(a_n,D_n)=1$.
Write the corresponding multipliers as $m_0,\ldots,m_{B-1}$ and let
$s$ be the index immediately after the last choice. The density estimate
just proved gives $s=B+o(B)$: deleting $O(\log B)$ initial indices and
$o(s)$ exceptional indices leaves $B$ choices.
Each $m_i>2B$ for large $B$, and the chosen multipliers are pairwise
coprime, because each earlier one divides the denominator preceding a
later one. All divide $D_s$. For $P=\prod_i m_i$ we therefore have
\[
 (2B)^B<P\le D_s,\qquad H_s<P,\qquad B<P,
 \qquad \ell(3P)\le B+o(B).
\]
The bound on $H_s$ uses $\log H_s=o(s)$ and $s=B+o(B)$.
Thus the product is large enough to place the block beyond the previous
maximum, but small enough that the allowed jump at its height is less
than the block length: $c\ell(3P)<B$ for all sufficiently large $B$.
Both comparisons are needed; existence of a distant CRT block alone
would not control a height-dependent jump bound.

\emph{Crossing the block.}
Choose $x\in[P,2P)$ by the Chinese remainder theorem so that
$m_i\mid x+i$ for $0\le i<B$. At the first record $C_t\ge x$
after $s$, the preceding maximum is below $x$ and its increase is
at most $c\ell(2P)<B$. Thus
$x\le C_t<x+B<3P$, so some $m_i$ divides both $C_t$ and $D_t$.
The exact updates preserve this common divisor. At the next record
$C_r$ we would therefore have
\[
 m_i\mid C_r-C_t,\qquad
 0<C_r-C_t\le c\ell(C_t)\le c\ell(3P)<B<m_i,
\]
which is impossible. This proves the auxiliary bound $\Theta\ge1$.

The integer normalisation matters here. Scaling $(C,D,E)$ by a positive
integer $k$ scales the running maximum by $k$ and its limit-superior
coefficient by $k$, since $\ell(kx)/\ell(x)\to1$. Thus the bound with
coefficient $1$ is not invariant under arbitrary clearing of denominators.
The next step keeps track of the common factor rather than discarding it.

\emph{The common gcd and the sharper coefficient.}
For any fixed $N$, divide the tail from $N$ onwards by $G_N$.
The resulting integer orbit satisfies the same hypotheses. Its running
maximum is eventually $H_n/G_N$, and
$\ell(H_n/G_N)/\ell(H_n)\to1$. The auxiliary bound therefore gives
$\Theta\ge G_N$. If $G_N$ is unbounded, then $\Theta=+\infty$.

Otherwise $G_n=g$ eventually. Fix a later index $T$ with $v_T>1$.
The reduced exact tail has pairwise coprime multipliers, each coprime
to $v_T$. For a large integer $L$, exactly
\[
 k_L=\frac{\varphi(v_T)}{v_T}L+O(v_T)
\]
of the offsets $0,\ldots,L-1$ are coprime to $v_T$.
Assign to those offsets the next $k_L$ multipliers, beginning at $T$.
Put $s=T+k_L$ and $Q=v_s=v_T\prod_{T\le j<s}a_j$.
Solve $x\equiv0\pmod{v_T}$ and $x\equiv-j$ modulo the multiplier
assigned to offset $j$, and take the solution in $[Q,2Q)$.
Every integer in $[x,x+L)$ then fails to be coprime to $v_s$.
No later reduced numerator can lie there.

The running maximum $R_s$ of the reduced numerator is below $Q$ for
large $L$, by the growth estimates above. The first crossing of $x$
after $s$ must therefore have a record increment at least $L$.
Its preceding maximum is below $2Q$, and
$\ell(2Q)\le s+O(1)=T+k_L+O(1)$. The corresponding record indices tend
to infinity with $L$, so
\[
 \limsup_n\frac{R_{n+1}-R_n}{\ell(R_n)}
 \ge\lim_{L\to\infty}\frac{L}{T+k_L+O(1)}
 =\frac{v_T}{\varphi(v_T)}.
\]
Since $H_n=gR_n$ eventually, this gives
$\Theta\ge g v_T/\varphi(v_T)>g$.
An eventually Sylvester tail instead has constant $C_n$ and $\Theta=0$.
Finally,
\[
 H_{n+1}-H_n=\bigl(-E_n-(H_n-C_n)\bigr)_+\le(-E_n)_+.
\]
Together with $\ell(C_n)\le\ell(H_n)$, this proves the corollary.
\end{proof}

\paragraph{Formalisation.}
The supplied formal sources contain persistence and pointwise-rise
versions of the arithmetic steps:
\lword{Erdos243/SlowRiseBarrier.lean}{111}{commonDivisor_persists}{persistence of
a common divisor}, the
\lword{Erdos243/SlowRiseBarrier.lean}{51}{exists_consecutiveMultiples_between}{bounded
CRT block}, the
\lword{Erdos243/SlowRiseBarrier.lean}{194}{slowRise_landing}{slow-rise landing
lemma}, and the
\lword{Erdos243/SlowRiseBarrier.lean}{280}{no_slowNegative_of_coprimeBlock}{coprime-block
exclusion}.  The growth estimates, the count of usable moduli and the running-maximum
argument are proved above in ordinary mathematics. The cited declarations
verify their own arithmetic statements, not this assembled theorem.
In particular, their landing bound controls increments of $C_n$;
the argument above controls increments of $H_n$ and is a written
extension, not a direct invocation of that declaration.

No upper bound on $\Theta$ is proved. If $G_n$ is unbounded, the theorem
forces $\Theta=+\infty$. If $G_n=g$ eventually, it proves only
$\Theta\ge g\,v_T/\varphi(v_T)$ for every late $T$. Since $v_T\mid v_{T+1}$,
the ratios $\varphi(v_T)/v_T$ are nonincreasing and have a limit
$\sigma\ge0$. If $\sigma=0$, the lower bound again forces
$\Theta=+\infty$; if $\sigma>0$, it gives $\Theta\ge g/\sigma$ but does
not determine whether $\Theta$ is finite. Replacing $H_n$ by the running
maximum of $u_n$ divides the coefficient by $g$, since
$\ell(gx)/\ell(x)\to1$.

\subsection{Bounds that allow for cancellation and earlier decreases}

\begin{theorem}[bounds allowing for previous decreases]\label{long243:res:recordamplified}
Under the standing hypotheses, the following are equivalent: eventual Sylvester behaviour; $\limsup_n\mathcal A_n<\infty$;
$\limsup_nR_n\delta_n<\infty$.  Each of those two limits superior is $0$ or
$+\infty$.
\end{theorem}
\leannote{Lean: \lproof{ErdosProblems/Erdos243/PaperCompleteR21/AmplifiedRecordEquivalence.lean}{975}{recordAmplified}, \lproof{ErdosProblems/Erdos243/PaperCompleteR21/AmplifiedRecordEquivalence.lean}{931}{amp_bddAbove_iff_sylvester}, \lproof{ErdosProblems/Erdos243/PaperCompleteR21/AmplifiedRecordEquivalence.lean}{948}{Rdelta_bddAbove_iff_amp}, \lproof{ErdosProblems/Erdos243/PaperCompleteR21/AmplifiedRecordEquivalence.lean}{587}{eventuallySylvester_of_amp_le}, and 6 further declarations in the \hyperref[long243:sec:coverage]{coverage section}.}

\begin{corollary}[the critical rate]\label{long243:res:criticalrate}
Under the standing hypotheses and $\delta_n=O(1/n)$, with no convergence of $n\delta_n$ assumed,
eventual Sylvester behaviour is equivalent to $u_n=O(n)$ and to $(-\tilde e_n)_+=O(1)$.  A
counterexample at the critical rate therefore has $\limsup_nu_n/n=\infty$ and
$\limsup_n(-\tilde e_n)_+=\infty$.
\end{corollary}
\leannote{Lean: \lproof{ErdosProblems/Erdos243/PaperCompleteR21/AmplifiedRecordEquivalence.lean}{1033}{criticalRate}, \lproof{ErdosProblems/Erdos243/PaperCompleteR21/AmplifiedRecordEquivalence.lean}{1197}{criticalRate_counterexample}, \lproof{ErdosProblems/Erdos243/PaperCompleteR21/AmplifiedRecordEquivalence.lean}{1016}{R_le_of_u_le}.}

\begin{proof}[Proof of Theorem~\ref{long243:res:recordamplified}]
Suppose the orbit is not eventually Sylvester and
$\mathcal A_n\le K$ eventually, for an integer $K\ge1$.
Absorption and $|\tilde e_n|/u_n\to0$ give $u_n\to\infty$.
Negative reduced errors must occur arbitrarily late, since otherwise
$h_nu_{n+1}=u_n-\tilde e_n$ would make $u_n$ eventually nonincreasing.

The bound controls cancellation as well as upward motion. Fix a late
index $s$, and let $t>s$ be the first subsequent negative-error index.
At the intervening steps the reduced numerator is nonincreasing, so
$u_t\le u_{s+1}$. Since $R_t\ge u_s$ and
$m_t=|\tilde e_t|\ge1$,
\[
 \mathcal A_t=\frac{R_tm_t}{u_t}\ge\frac{u_s}{u_{s+1}}
 =\frac{h_s}{1-\tilde e_s/u_s}.
\]
Thus $h_s\le3K/2<2K$ after the relative-error threshold.
Also $m_n\le\mathcal A_n\le K$, since $R_n\ge u_n$, and hence
$u_{n+1}\le u_n+K$.

We can now choose primes too large to be removed by cancellation.
The density-one argument in the proof of
Theorem~\ref{long243:res:recorddichotomy} supplies infinitely many
late multipliers $a_n$ coprime to $D_n$. These multipliers are pairwise
coprime. At each such step, $a_n$ is coprime to $v_n$ and
$\gcd(u_n,v_n)=1$, so $a_nu_n-v_n$ is coprime to both $a_n$ and $v_n$.
Hence $h_n=1$ and $a_n\mid v_{n+1}$.
Only finitely many of the chosen multipliers can have all their prime
divisors at most $2K$, since each uses a different prime from that finite
set. Choose $K$ distinct primes $p_0,\ldots,p_{K-1}>2K$ at these
steps. Once a chosen prime divides a reduced denominator, it divides
every later one: $h_nv_{n+1}=a_nv_n$ and $h_n<2K<p_i$ prevent its
removal. All the primes therefore divide $v_T$ at a common later index $T$.

Consequently every $u_n$ with $n\ge T$
is coprime to all the $p_i$. Choose a CRT block of $K$ consecutive
integers, one divisible by each $p_i$, above $R_T$.
Divergence forces a first crossing, while
$u_{n+1}\le u_n+K$ forces it to land inside the block, a contradiction.
This proves that finite $\limsup\mathcal A_n$ forces a Sylvester tail.

The comparison
$|\delta_n-m_n/u_n|\le3/a_n$ gives
$|R_n\delta_n-\mathcal A_n|\le3R_n/a_n\to0$.
Here $R_n\le H_n$, $\log H_n=o(n)$ and $a_n$ grows doubly
exponentially. Thus the two boundedness conditions are equivalent.
On a Sylvester tail, $u_n=1$ and $m_n=0$ eventually, so
$\mathcal A_n\to0$ and $R_n\delta_n\to0$. Since both quantities are
nonnegative, their limits superior can only be $0$ or $+\infty$.
\end{proof}

The estimate at the first later negative error is the lemma in the release,
\href{https://github.com/wcook04/plectis-erdos-lean/blob/52f29ad173b04e3bac941b3663f2b9aebe5de0bb/ErdosProblems/Erdos243/SaturatedSquareTransport.lean\#L378}{negative error scaled by the previous maximum after cancellation}: at the first later negative error,
$u_t\le u_{s+1}$ and $u_su_t\le R_t|\tilde e_t|u_{s+1}$.
The supplied index also records this lemma in a built public module.
That evidence concerns the lemma, not a formal verification of the complete
proof above. The exact example
$15/134=1/10+1/85+1/5695$, with steps in reduced fractions
$(15,134)\to(4,335)\to(1,5695)$ and $\mathcal A_1=15/4$,
attains equality in the displayed bound.

For Corollary~\ref{long243:res:criticalrate}, the comparison
$|\delta_n-m_n/u_n|\le3/a_n$ and the rate $\delta_n=O(1/n)$ give
$m_n/u_n=O(1/n)$. If $u_n=O(n)$, then $m_n=O(1)$. Conversely,
$m_n=O(1)$ gives $u_{n+1}\le u_n+m_n$, hence $u_n=O(n)$ and
$R_n=O(n)$. In either case
$\mathcal A_n=R_nm_n/u_n=O(1)$, so
Theorem~\ref{long243:res:recordamplified} gives eventual Sylvester
behaviour. A Sylvester tail has $u_n=1$ and $m_n=0$ eventually,
which proves the converse implications.

The condition bounds the negative reduced error after multiplying it by
$R_n/u_n$. At a current maximum this is just $(-\tilde e_n)_+$; after a
decrease the multiplier is larger. A bound on $(-\tilde e_n)_+$ does
not directly control this additional factor. Under the critical-rate
hypothesis, the preceding corollary supplies the required control.
A Sylvester tail satisfies the condition with eventual value zero.
We do not derive it for every rational tail satisfying the growth hypothesis.

\begin{lemma}[large odd prime powers in the reduced denominator]\label{long243:res:oddpowersupply}
Under the standing hypotheses, for every fixed $A>0$ and every sufficiently large $n$, the
reduced denominator $v_n$ has an odd prime-power divisor $Q=p^k$ with
\[
 Q>(H_n+2)^A,\qquad H_n=\max_{j\le n}C_j.
\]
The prime $p$ may depend on $n$; no stable-gcd or prime-arrival assumption
is imposed.
\end{lemma}
\leannote{Lean: \lproof{ErdosProblems/Erdos243/PaperCompleteR21/ReducedDenominatorPrimePowers.lean}{790}{oddPrimePower_supply}, \lproof{ErdosProblems/Erdos243/PaperCompleteR21/ReducedDenominatorPrimePowers.lean}{718}{oddPrimePower_supply_nat}.}
\begin{proof}
The denominator grows too fast for all its odd prime-power factors to
remain small compared with $H_n$. Indeed, the estimate for the rational
tail gives $C_{n+1}/C_n\to1$, hence $\log H_n=o(n)$. Since $x_n=u_n/v_n\le2/a_n$ and $u_n\ge1$,
$v_n\ge a_n/2$. Also $v_n\mid L_n=\lcm(q,a_0,\ldots,a_{n-1})$,
so the $2$-primary part of $v_n$ is at most $\max(q,a_{n-1})=a_{n-1}$
for all large $n$. Its odd part $W_n$ therefore satisfies
\[
 W_n\ge\frac{a_n}{2a_{n-1}}\ge\frac{a_{n-1}}4,
 \qquad \log W_n\ge c2^{n-1}-O(1)
\]
for some $c>0$.
Put $B_n=(H_n+2)^A=\exp(o(n))$. If every exact odd prime-power factor of
$W_n$ were at most $B_n$, then
$W_n\mid\lcm(1,\ldots,\lfloor B_n\rfloor)$, giving
\[
 \log W_n\le B_n\log B_n=\exp(o(n)).
\]
This contradicts the preceding exponential lower bound for all sufficiently
large $n$. An offending exact prime-power factor is the required $Q$.
\end{proof}
The argument compares the logarithm of the odd denominator part with the
logarithm of a factorial bound. It needs neither the prime number theorem
nor a lower density of new primes, and makes no such claim.

\begin{theorem}[unit record increments]\label{long243:res:unitrecord}
Under the standing hypotheses, if $R_{n+1}-R_n\le1$ for all large $n$, then $a_{n+1}=a_n^2-a_n+1$
for all large $n$.  Hence the sequence is eventually Sylvester if and only if
$\#\{n:R_{n+1}-R_n\ge2\}$ is finite.  No hypothesis is placed on drawdowns, on
record-setting jumps, or on the cancellation factors $h_n$.
\end{theorem}
\leannote{Lean: \lproof{ErdosProblems/Erdos243/PaperCompleteR21/ReducedDenominatorPrimePowers.lean}{846}{unitRecordIncrement_sylvesterNext}, \lproof{ErdosProblems/Erdos243/PaperCompleteR21/ReducedDenominatorPrimePowers.lean}{950}{unitRecordIncrement_criterion}.}

\begin{proof}
Suppose the orbit is not eventually Sylvester. Absorption gives
$\tilde e_n\ne0$ on a late tail, so vanishing relative error and integrality give
$u_n\to\infty$. Choose $s$ beyond the unit-increment and centring
thresholds and the threshold in Lemma~\ref{long243:res:oddpowersupply}
with $A=3$.
Since $H_s\ge R_s\ge1$, it supplies an odd $Q=p^k\mid v_s$ with
$Q>(H_s+2)^3>4(R_s+2)$ and $Q\ge16$.
Let $Y$ be the least multiple of $p$ above $R_s$. Then
$Y\le R_s+p< Q/4+p\le5Q/4$.
At the first $t>s$ with $u_t\ge Y$, the integral unit-record increment
forces $u_t=Y$. Before $t$, one has
$w_n<3u_n/2<15Q/8<pQ=p^{k+1}$.
Corollary~\ref{long243:res:powerpersistence} therefore gives $p^k\mid v_t$.
But $p\mid u_t=Y$, contradicting $\gcd(u_t,v_t)=1$.
The converse follows because a Sylvester tail has $u_n=1$ eventually.
\end{proof}

The separate release contains the
\href{https://github.com/wcook04/plectis-erdos-lean/blob/52f29ad173b04e3bac941b3663f2b9aebe5de0bb/ErdosProblems/Erdos243/RecordIncrementBarrier.lean\#L329}{unit-increment landing lemma} and the
\href{https://github.com/wcook04/plectis-erdos-lean/blob/52f29ad173b04e3bac941b3663f2b9aebe5de0bb/ErdosProblems/Erdos243/RecordIncrementBarrier.lean\#L425}{one-unit record-increment implication}.
The supplied index also records these statements in a built public
module. Their assumption about the existence of a prime power is proved
here by the preceding written lemma. The complete theorem about the
rational tail is not thereby an assembled Lean theorem.


The two-unit criterion discussed below bounds the actual jump at each
record step, rather than the increase in the running maximum. As numerical
bounds on individual steps, neither condition implies the other: the first
allows $s_n=2$, while $s_n\le1$ permits large jumps after an earlier
decrease. This is a comparison of the bounds, not a claim that they remain
logically independent under all the standing tail hypotheses; there both
criteria force eventual Sylvester behaviour.
The finite example
$(8,177)\to(7,4071)\to(10,2373393)$ with multipliers $23$ and $583$ has records
$8$ and $10$, increment $2$, jump $3$ and $\gcd(8,10)=2$, which is why the
parity argument for crossing an odd level does not transfer from actual
jumps to increments of the running maximum.

\subsection{Counting jumps before a prime power can be lost}

\begin{theorem}[counting crossings before a prime power is lost]\label{long243:res:epochenergy}
Let the orbit satisfy the reduced recurrences of this section, with
$2|\tilde e_n|<u_n$ from an
index $s$.  Let $p\ge3$ be prime, $Q=p^{\ell}$ divide $v_s$ with $Q\ge16$, put
$L=pQ/2$, and assume $R_s<L/2$. Assume that $u_t\ge L$ for some $t>s$,
and let $\tau$ be the first such index, let $J$ be the set of steps in $[s,\tau)$ that first cross at least one odd
multiple of $p$ in $(L/2,L]$, and put $X=\sum_{n\in J}(d_n-2)$.  Then every
$n\in J$ is a record step with $h_n=1$ and $d_n\ge3$, and
$pQ\le(8p+8)\lvert J\rvert+4X+8p$.
\end{theorem}
\leannote{Lean: \lproof{ErdosProblems/Erdos243/PaperCompleteR21/ProtectedEpochBarrierCount.lean}{149}{epoch_energy_named_crossing_set}, \lproof{ErdosProblems/Erdos243/PaperCompleteR21/ProtectedEpochBarrierCount.lean}{34}{mem_barrierIdx_iff}.}

\begin{proof}[Proof of Theorem~\ref{long243:res:epochenergy}]
For $s\le n<\tau$, one has $u_n<L$ and
$w_n<3u_n/2<3pQ/4<p^{\ell+1}$, since $Q=p^{\ell}$.
Corollary~\ref{long243:res:powerpersistence} gives $Q\mid v_t$ throughout
$[s,\tau]$, so $p\nmid u_t$ there. A first crossing of a level above $R_s$
is a record step. If $h_n\ge2$ then $u_{n+1}=w_n/h_n<3u_n/4$, so every
such record step has $h_n=1$.
An odd multiple of $p$ cannot be landed on. A jump of size $1$ crossing it
would land on it; a jump of size $2$ avoiding the landing would have two
even endpoints, contrary to $\gcd(u_n,u_{n+1})=1$. Thus every $n\in J$
has $d_n\ge3$.
It remains to count these levels. The interval $(L/2,L]$ has length
$pQ/4$, and odd multiples of $p$ are spaced by $2p$, so it contains
at least $Q/8-1$ of them. The levels first crossed by a step of size
$d_n$ have the same spacing, so
there are at most $1+d_n/(2p)$ of them. Summing over $J$ gives
\[
 Q/8-1\le |J|+\frac{2|J|+X}{2p},
\]
which rearranges to the asserted inequality.
\end{proof}

The next series ignores record jumps of size at most two, apart from
finitely many initial terms, and assigns a smaller cost to a large jump
when it begins at a large numerator. A Sylvester tail has no late records,
so both series converge. The proof shows that a non-Sylvester rational
tail contributes a fixed positive amount on arbitrarily late finite
intervals. Thus convergence is a genuine additional requirement, not a
restatement of the pointwise limit $|E_n|/C_n\to0$.

\begin{theorem}[two convergence criteria for large record jumps]\label{long243:res:energycriterion}
Under the standing hypotheses, the sum
$\mathcal E=\sum_{n\ \mathrm{record}}\bigl(\mathbf 1_{d_n\ge3}u_n^{-1/2}
+(d_n-2)_+u_n^{-1}\bigr)$ is finite if and only if the sequence is eventually Sylvester; and
$\sum_{n\ \mathrm{record}}(d_n-2)_+u_n^{-1/2}$ is finite if and only if the sequence is eventually Sylvester.
\end{theorem}
\leannote{Lean: \lproof{ErdosProblems/Erdos243/PaperCompleteR21/RecordJumpEnergySeries.lean}{468}{energy_criterion}, \lproof{ErdosProblems/Erdos243/PaperCompleteR21/RecordJumpEnergySeries.lean}{414}{energy_summable_iff}, \lproof{ErdosProblems/Erdos243/PaperCompleteR21/RecordJumpEnergySeries.lean}{455}{energySqrt_summable_iff}, \lproof{ErdosProblems/Erdos243/PaperCompleteR21/RecordJumpEnergySeries.lean}{161}{energy_le_two_energySqrt}, and 2 further declarations in the \hyperref[long243:sec:coverage]{coverage section}.}

\begin{proof}
Assume first that the recurrence is not eventually Sylvester. Then $u_n\to\infty$.
For every sufficiently late $s$, Lemma~\ref{long243:res:oddpowersupply}
with $A=3$ gives an odd $Q=p^k\mid v_s$ with
$Q>(H_s+2)^3$, hence $Q\ge16$ and $pQ>4R_s$.
Set $L=pQ/2$ and take its first crossing. The preceding theorem applies.
Because $Q\ge p$ and $p\ge3$,
\[
 \frac{8p}{L}=\frac{16}{Q}\le1,
 \qquad \frac{8p+8}{\sqrt L}
 \le8\sqrt2(1+1/p)<16.
\]
Dividing the preceding counting inequality by $L$ therefore gives
\[
 1\le16\frac{|J|}{\sqrt L}+4\frac{X}{L}
 \le16\sum_{n\in J}\left(\frac1{\sqrt{u_n}}+
                  \frac{d_n-2}{u_n}\right),
\]
since $u_n<L$ for $n<\tau$ and $d_n\ge3$ on $J$.
Thus an arbitrarily late finite window contributes at least $1/16$ to
$\mathcal E$, contradicting convergence. No disjointness of the windows is
needed: every window lies in a tail of the nonnegative series.
On a Sylvester tail $u_n=1$ eventually, so there are no late records and
$\mathcal E$ is finite.
Finally, term by term at a record,
\[
 \mathbf1_{d_n\ge3}u_n^{-1/2}+(d_n-2)_+u_n^{-1}
 \le2(d_n-2)_+u_n^{-1/2}.
\]
Hence finiteness of the second series implies finiteness of the first,
and convergence in the converse direction again follows from eventual Sylvester behaviour.
\end{proof}

\paragraph{Formalisation.}
The separate release contains
\href{https://github.com/wcook04/plectis-erdos-lean/blob/52f29ad173b04e3bac941b3663f2b9aebe5de0bb/ErdosProblems/Erdos243/ProtectedEpochEnergy.lean\#L185}{the integer inequality counting crossings before a prime power is lost}
and its counting lemmas.
The supplied index also records the integer inequality in a built
public module. The existence of the required prime power, the real-valued
bound and the complete equivalence above have written proofs; the Lean inequality is the checked
integer form, not a checked transcription of the whole assembled theorem.
The checked main-repository declaration
\lword{Erdos243/PrimitiveRecordBarrier.lean}{435}{recordRiseTwo_sylvesterNext_eventually}{the two-unit record-rise implication}
retains a supply premise: arbitrarily late $s$ admit an odd
$p^k\mid v_s$ with $3R_s<p^k$.
Lemma~\ref{long243:res:oddpowersupply} supplies this premise in ordinary
mathematics. The declaration itself has not thereby become a complete
checked theorem about rational tails. This distinction is about formal coverage, not
an unfilled mathematical supply argument in the text.

\subsection{Bounds on the error and on the original sequence}

\begin{theorem}[slow negative part]\label{long243:res:slownegative}
Let $(a,C,D)$ be an exact orbit of natural numbers with $a_n>1$, $C_n>0$, $D_0\ge1$, under
vanishing relative error.  Suppose that for some $\delta\in(0,1)$ and all large $n$
with $E_n<0$ one has $-E_n\le(1-\delta)\ell(C_n)$.  Then $E_n=0$ for all large
$n$, and $a_{n+1}=a_n^2-a_n+1$ for all large $n$.
\end{theorem}
\leannote{Lean: \lproof{ErdosProblems/Erdos243/PaperCompleteR21/ExactOrbitRecordDichotomy.lean}{559}{slowNegative_eventually_zero_and_sylvesterNext_unconditional}, \lproof{ErdosProblems/Erdos243/PaperCompleteR21/ExactOrbitRecordDichotomy.lean}{363}{exactOrbit_recordTheta_gt_one}, \lproof{ErdosProblems/Erdos243/PaperCompleteR21/ExactOrbitRecordDichotomy.lean}{522}{exactOrbit_one_le_recordTheta}, \lproof{ErdosProblems/Erdos243/PaperCompleteR21/ExactOrbitRecordDichotomy.lean}{288}{exactOrbit_unbounded_of_error_not_eventually_zero}, and 3 further declarations in the \hyperref[long243:sec:coverage]{coverage section}.}

\noindent Theorem~\ref{long243:res:bounded} is the case of a constant bound, since a
constant is eventually below $(1-\delta)\ell(C_n)$ on a nonzero tail,
where $C_n\to\infty$.

\begin{proof}
Suppose the error is not eventually zero. The auxiliary lower bound in
the proof of Theorem~\ref{long243:res:recorddichotomy} applies to this
exact integer orbit: it used only $a_n>1$, $C_n>0$, $D_0\ge1$ and
vanishing relative error. It gives
$\limsup_n(H_{n+1}-H_n)/\ell(H_n)\ge1$.
On the other hand, the hypothesis and $C_n\le H_n$ give
\[
 H_{n+1}-H_n\le(-E_n)_+
 \le(1-\delta)\ell(C_n)\le(1-\delta)\ell(H_n)
\]
for all large $n$, a contradiction. Thus $E_n=0$ eventually, and
Theorem~\ref{long243:res:eventual} gives the Sylvester recurrence.
\end{proof}

\paragraph{Formalisation.}
The corresponding pointwise-rise argument has formalised arithmetic
lemmas:
\lword{Erdos243/SlowRiseBarrier.lean}{111}{commonDivisor_persists}{persistence of
a common divisor}, the
\lword{Erdos243/SlowRiseBarrier.lean}{162}{multiplierOverlap_dvd_laterCenteredState}{overlap
transport}, the
\lword{Erdos243/SlowRiseBarrier.lean}{194}{slowRise_landing}{landing lemma} and
the
\lword{Erdos243/SlowRiseBarrier.lean}{280}{no_slowNegative_of_coprimeBlock}{coprime-block
exclusion}.  The count of indices and the simultaneous choice of constants are
written out in the earlier proof; no end-to-end Lean verification of
this consequence is claimed. Corollary~\ref{long243:res:loglogboundary} includes the
boundary coefficient $1$, while the preceding theorem on the running
maximum also discounts rises that recover an earlier decrease. The present
statement is retained for its direct bound on $E_n$.

To compare the bounded hypothesis with the classical criteria, number
the original sequence from $1$ in this subsection and put
$P_n=\prod_{1\le j<n}a_j$. Then
\[
 \frac{P_n}{a_n}\left(\frac{a_n^2}{a_{n+1}}-1\right)
 =\frac{P_{n+1}}{a_{n+1}}-\frac{P_n}{a_n}.
\]
Thus the bound concerns upward increments, not boundedness of $P_n/a_n$.
The quadratic growth limit controls only the ratio of consecutive terms.
Sylvester tails satisfy the bound, as do exact-square sequences
$a_n=b^{2^{n-1}}$ with $b\ge2$, for which the difference is zero.
The latter never satisfies the Sylvester recurrence, so the bounded
criterion below recovers irrationality of its reciprocal sum.
The rounded examples after Corollary~\ref{long243:res:onethreshold}
show more generally how a term $c/n$ in the ratio produces increments
of order $n^{c-1}$. The theorem allows a finite positive upper limit
where the classical product criterion requires a nonpositive one;
neither bound is derived from the original problem alone.

\begin{theorem}[bounded or slowly growing increments of the product ratio]\label{long243:res:strausbounded}
Let $a_1<a_2<\cdots$ be positive integers with $a_{n+1}/a_n^2\to1$
and $\sum_{n\ge1}1/a_n=p/q$, where $p,q$ are positive integers. Put
\[
 Q_n=\frac{a_1a_2\cdots a_{n-1}}{a_n}
     \left(\frac{a_n^2}{a_{n+1}}-1\right).
\]
If $\limsup_nQ_n<\infty$, the sequence is eventually Sylvester.
The same conclusion holds if, for some $\delta>0$ and all large $n$,
\[
 Q_n\le\frac{1-\delta}{q}\,
 \ell(a_1\cdots a_{n-1}/a_n).
\]
\end{theorem}
\leannote{Lean: \lproof{ErdosProblems/Erdos243/PaperCompleteR7/ProductDefect.lean}{211}{original_coordinate_bounded_defect}, \lproof{ErdosProblems/Erdos243/PaperCompleteR21/SlowGrowthProductIncrements.lean}{155}{original_coordinate_slow_growth_defect}, \lproof{ErdosProblems/Erdos243/PaperCompleteR21/SlowGrowthProductIncrements.lean}{61}{recordTheta_le_of_slow_negative}, \lproof{ErdosProblems/Erdos243/PaperCompleteR21/SlowGrowthProductIncrements.lean}{111}{prefix_ratio_le_canonicalNumerator}.}

For the bounded clause, use the stated denominator $q$, put
$P_n=\prod_{j<n}a_j$, $x_n=\sum_{k\ge n}1/a_k$, and form the integer tail
$D_n=qP_n$, $C_n=D_nx_n$, $E_n=D_n-(a_n-1)C_n$. Exact cancellation gives
\[
 E_n+qQ_n=qP_n\left(\frac1{a_{n+1}}
 -(a_n-1)\sum_{k\ge n+2}\frac1{a_k}\right).
\]
The expression in parentheses is positive eventually, since
$\sum_{k\ge n+2}1/a_k\le4/a_{n+1}^2$ and
$a_{n+1}>4(a_n-1)$ eventually. Thus $E_n\ge-qQ_n$,
and $\limsup Q_n<\infty$ supplies the eventual lower bound required
by Theorem~\ref{long243:res:bounded}. This proves the bounded clause
without first estimating the absolute size of the comparison error.

For the slow-growth clause and the later summability comparisons,
the same exact formula gives
\[
 0<E_n+qQ_n\le qP_n/a_{n+1}=O(P_n/a_n^2)
\]
eventually. The ratio test in Section~\ref{long243:sec:lcmrecords}
shows that these errors have a convergent sum, hence tend to zero.
The factor $1/q$ in the slow-growth clause cancels the factor $q$ in
$E_n+qQ_n=o(1)$. Suppose the sequence is not eventually Sylvester. Absorption
and vanishing relative error give $C_n\to\infty$, and the tail estimate
$C_n\sim qP_n/a_n$ implies
\[
 \frac{\ell(C_n)}{\ell(P_n/a_n)}\longrightarrow1.
\]
For $0<\delta<1$, the assumed bound and $E_n+qQ_n=o(1)$ consequently give
$(-E_n)_+\le(1-\delta/2)\ell(C_n)$ eventually. This contradicts
Theorem~\ref{long243:res:slownegative}. If $\delta\ge1$, the hypothesis gives
$Q_n\le0$ eventually and the bounded clause already applies. The growth
argument used in this deduction is a written argument, not an additional
assembled Lean theorem.

The same comparison permits any positive real clearing factor
$F_n\le qP_n$: multiplying $E_n+qQ_n=o(1)$ by $F_n/(qP_n)$ gives
\[
 F_n-(a_n-1)F_nx_n+
 \frac{F_n}{a_n}\left(\frac{a_n^2}{a_{n+1}}-1\right)=o(1).
\]
Neither $F_nx_n$ nor $F_n-(a_n-1)F_nx_n$ need be integral. The LCM
choice $F_n=L_n$ makes both integers, but its update includes the
factor $\rho_n$ of Section~\ref{long243:sec:lcmrecords}. The short note
uses only this LCM specialisation; no general clearing factor is
needed in its bounded-increment proof.

\paragraph{Attribution.}
Erd\H{o}s and Straus require $\limsup\le0$ in the corresponding criterion, with
the least common multiple in place of the product and the growth factor one
index later, so their quantity is
$[a_1,\ldots,a_n]a_{n+1}^{-1}(a_{n+1}^{2}/a_{n+2}-1)$~\cite[Theorem~3,
p.~132]{erdosstraus1964}; Koizumi's Corollary~4(1) uses the product
form~\cite[Cor.~4(1), pp.~14--15]{koizumi2025}.  The bounded clause follows from
Theorem~\ref{long243:res:bounded}. For Koizumi's product expression it
replaces a nonpositive upper limit by an arbitrary finite upper limit.
The material distinction here is product versus least common multiple.
Reindexing $r=n+1$ makes the classical LCM expression
$[a_1,\ldots,a_{r-1}]a_r^{-1}(a_r^2/a_{r+1}-1)$, exactly the expression
in the short note's LCM corollary.

The comparison $E_n+qQ_n=o(1)$ has its classical predecessor in
Koizumi's proof of Corollary~4(1) \cite[(11)--(12), p.~15]{koizumi2025}.
The bounded clause uses the checked rational-tail estimates; the
slow-growth deduction above is an ordinary written argument.

\begin{corollary}[the $1/n$ threshold]\label{long243:res:onethreshold}
Let $a_1<a_2<\cdots$ be positive integers with $a_{n+1}/a_n^2\to1$
and $\sum_n1/a_n\in\Q$. If
\[
 \limsup_{n\to\infty} n\,(a_n^2/a_{n+1}-1)_+<1,
\]
then $a_{n+1}=a_n^2-a_n+1$ for all large $n$. The same conclusion holds if, for some $K\ge0$ and $\varepsilon>0$,
\[
 \frac{a_n^2}{a_{n+1}}-1\le\frac1n+\frac{K}{n^{1+\varepsilon}}
 \quad\hbox{eventually}.
\]
In particular the one-sided bound by $1/n$ is included.
\end{corollary}
\leannote{Lean: \lproof{ErdosProblems/Erdos243/PaperCompleteR21/ProductDefectThresholds.lean}{92}{original_coordinate_strict_one}, \lproof{ErdosProblems/Erdos243/PaperCompleteR20/InclusiveOne.lean}{207}{original_coordinate_inclusive_one}, \lproof{ErdosProblems/Erdos243/PaperCompleteR20/InclusiveOne.lean}{225}{original_coordinate_inclusive_one_pointwise}.}

\begin{proof}
Put $\gamma_n=a_n^2/a_{n+1}-1$ and $t_n=\bigl(\prod_{j<n}a_j\bigr)/a_n$, so
that $t_{n+1}/t_n=1+\gamma_n$.  Choose $r\in(0,1)$ and $N$ with
$\gamma_n^{+}\le r/n$ for all $n\ge N$.  Then
\[
 t_n\le t_N\prod_{k=N}^{n-1}\Bigl(1+\frac rk\Bigr)=O(n^{r}),
\]
so $(t_n\gamma_n)_+=t_n\gamma_n^{+}=O(n^{r-1})$ tends to zero and
$\limsup_nt_n\gamma_n\le0$.  Koizumi's Corollary~4(1)
\cite[pp.~14--15]{koizumi2025} gives the first conclusion.
For the second, positivity of $1+\gamma_n$ and
\[
 1+\gamma_n\le(1+1/n)(1+K/n^{1+\varepsilon})
\]
give $t_n=O(n)$, since the second product converges. Consequently
$t_n(\gamma_n)_+=O(1)$, and the criterion bounding the product expression applies.
\end{proof}

The inclusive clause requires more than $\limsup n(\gamma_n)_+\le1$.
That weaker condition does not give the product estimate: for example,
the scalar sequence $\gamma_n=(1+1/\log n)/n$ has $t_n\asymp n\log n$
and $t_n\gamma_n\asymp\log n$. This is a limitation of the estimate, not a
counterexample to the original problem.

More generally, if $c>0$, $\varepsilon>0$ and
\[
 \frac{a_n^2}{a_{n+1}}=1+\frac cn+O(n^{-1-\varepsilon}),
\]
then the consecutive-ratio identity gives $t_n\sim Kn^c$ for some
$K>0$ and $t_{n+1}-t_n\sim Kc n^{c-1}$. To see the constant, take
logarithms and subtract $c\log(1+1/n)$; the remainder is absolutely summable.
Hence the increments are bounded for $c\le1$ and unbounded for $c>1$.
These rates are realised by increasing integer sequences: set
$a_{n+1}=\lceil n a_n^2/(n+c)\rceil$ and choose the seed large enough
that $a_{n+1}\ge a_n^2/(1+c)\ge2a_n$ throughout. The rounding remainder
in $[0,1)$ gives an error in the ratio bounded by
$(1+c)^2/a_n^2$, smaller than every fixed inverse power of $n$.
This calculation supplies the family comparison behind the two rounded
examples retained in the short note. It concerns their growth and does
not assume rationality.

The second clause has a concrete application outside Koizumi's
nonpositive-upper-limit product condition. The sequence
$a_1=4$, $a_{n+1}=\lceil n a_n^2/(n+1)\rceil$ begins
$4,8,43,1387,\ldots$ and satisfies
\[
 a_n\ge2\cdot2^{2^{n-1}},\qquad
 0\le1+\frac1n-\frac{a_n^2}{a_{n+1}}<\frac4{a_n^2}.
\]
The lower bound follows from $a_{n+1}\ge a_n^2/2$; the upper error bound
follows by writing the rounding remainder in $[0,1)$.
Consequently $t_n\sim Kn$ and $t_n\gamma_n\to K>0$ by the product
comparison above. The corollary proves that the reciprocal sum is
irrational, since a Sylvester tail would have
$\gamma_n=O(1/a_n)$ rather than $\gamma_n\sim1/n$.

The strict clause includes Koizumi's sufficient
rate $1+o(1/n)$ \cite[Remark~3, p.~16]{koizumi2025}; the inclusive clause
allows a bounded positive increment rather than requiring a nonpositive
upper limit.

\subsection{Signs of the error and of the growth ratio}

We keep the preceding one-based indexing. The growth defect need not
have the opposite sign to $E_n$: the next identity displays the
correction, including its sign on a Sylvester tail.

\begin{proposition}[comparison of the two signs]\label{long243:res:classicalhalfspace}
The factor $M_n=D_n/L_n$ divides $G_n=\gcd(C_n,D_n)$. Thus $G_n/M_n$
is a positive integer and $E_n/M_n=(G_n/M_n)\tilde e_n$ is an integer
with the same sign as $E_n$. The sign of the growth ratio minus one also
depends on a correction term. The exact recurrences give, whenever
$a_{n+1}C_nC_{n+1}\ne0$,
\begin{equation}\label{long243:eq:shiftedsign}
 \frac{a_n^2}{a_{n+1}}-1=-\frac{E_n}{C_n}+\Lambda_n,
 \qquad
 \Lambda_n=\frac{\bigl(1-E_n/C_n\bigr)\bigl(a_n-1+E_{n+1}/C_{n+1}\bigr)}{a_{n+1}},
\end{equation}
and, under the standing positive rational-tail hypotheses,
$0<\Lambda_n<3/a_n$ for all sufficiently large $n$. The
Erd\H{o}s--Straus quantity of Theorem~3 is
\[
 Z_n^{\mathrm{ES}}=
 \frac{[a_1,\ldots,a_n]}{a_{n+1}}
 \left(\frac{a_{n+1}^2}{a_{n+2}}-1\right).
\]
Its least common multiple includes $a_n$ but not the clearing
denominator $q$. It is not $Q_n=(P_n/a_n)\gamma_n$ from the preceding subsection,
nor $L_n\gamma_{n+1}/a_{n+1}$ under our convention
$L_n=\lcm(q,a_1,\ldots,a_{n-1})$.
Being a positive multiple of the next growth defect, $Z_n^{\mathrm{ES}}$
has the sign of $\Lambda_{n+1}-E_{n+1}/C_{n+1}$. For all sufficiently
large $n$, it is positive when $E_{n+1}\le0$; for $E_{n+1}>0$, it is
negative precisely when
$E_{n+1}/C_{n+1}>\Lambda_{n+1}$. On a Sylvester tail $E_n=0$ and
$\Lambda_n=(a_n-1)/a_{n+1}>0$.
\end{proposition}
\leannote{Lean: \lproof{ErdosProblems/Erdos243/PaperCompleteR21/ClassicalHalfspaceSigns.lean}{206}{overlapDebt_dvd_gcd}, \lproof{ErdosProblems/Erdos243/PaperCompleteR21/ClassicalHalfspaceSigns.lean}{225}{canonicalError_div_overlapDebt}, \lproof{ErdosProblems/Erdos243/PaperCompleteR21/ClassicalHalfspaceSigns.lean}{325}{growthDefect_eq_neg_relativeError_add_shiftedCorrection}, \lproof{ErdosProblems/Erdos243/PaperCompleteR21/ClassicalHalfspaceSigns.lean}{415}{canonical_growthDefect_identity}, and 8 further declarations in the \hyperref[long243:sec:coverage]{coverage section}.}

\begin{proof}
Since both $C_n=D_nx_n$ and $L_nx_n$ are integers, $M_n=D_n/L_n$
divides $C_n$ as well as $D_n$. It therefore divides $G_n$.  For \eqref{long243:eq:shiftedsign}, write
$\theta_n=E_n/C_n$; Proposition~\ref{long243:res:update} gives $C_{n+1}=C_n(1-\theta_n)$
and $D_n/C_n=a_n-1+\theta_n$, so
$a_{n+1}-1+\theta_{n+1}=a_nD_n/C_{n+1}=a_n(a_n-1+\theta_n)/(1-\theta_n)$.
Multiplying the asserted identity by $a_{n+1}$ and substituting reduces it to
$a_{n+1}+a_n-1+\theta_{n+1}=a_n^2/(1-\theta_n)$, which is the displayed
relation with $a_n$ added to both sides.  The bound follows from
$\theta_n\to0$ and $a_{n+1}\ge a_n^2/2$.  On Sylvester's sequence
$\theta_n=0$ and $a_{n+1}=a_n^2-a_n+1$, so
$a_n^2/a_{n+1}-1=(a_n-1)/a_{n+1}=\Lambda_n>0$.
\end{proof}


A nonpositive upper limit admits positive values tending to zero; it is
not an eventual pointwise sign condition. For the \emph{product} quantity
$Q_n=(P_n/a_n)\gamma_n$, the identity $E_n+qQ_n=o(1)$ shows that eventual
$E_n\ge0$ implies $\limsup Q_n\le0$, as in Koizumi's Corollary~4(1)
\cite[pp.~14--15]{koizumi2025}.
For the classical LCM expression $Z_n^{\mathrm{ES}}$, the corresponding
integer error is also taken one index later. Its published hypotheses are
those of Erd\H{o}s--Straus Theorem~3
\cite[p.~132]{erdosstraus1964}. Reindexing relates it to an LCM-cleared
error, not to $E_n$ with the unchanged product factor $qQ_n$.

\paragraph{A comparison with $\mathcal B$-free integers.}
For a fixed family $\mathcal B=\{m_i\}$, the integers divisible by none
of the $m_i$ are called $\mathcal B$-free integers. With pairwise coprime
moduli and $\sum_i1/m_i<\infty$, this is the standard setting of
\cite[\S1.2]{elabdalaoui2015}. The next proposition is an elementary
interval count in that setting. It tests what can follow from avoiding
whole multiples alone. Unlike the reciprocal-tail argument, it imposes
all the moduli from the outset and has no denominator recurrence or
cancellation factors.

\begin{proposition}[an elementary interval bound]\label{long243:res:coprimalitycap}
Let $m_0<m_1<\cdots$ be pairwise coprime integers at least $2$ with
$\theta=\sum_i1/m_i<1$. For all integers $x\ge1$ and $L\ge1$ satisfying
$L>k/(1-\theta)$, where $k=\#\{i:m_i\le x+L\}$, the interval $[x,x+L)$ contains
an integer divisible by no $m_i$. If also $\ell(m_i)=i+O(1)$, then for
every $\epsilon>0$ there are an index $T$ and a strictly increasing
sequence of positive integers $(u_n)$ such that
\[
 m_i\nmid u_n\quad\text{for every }i\ge T\text{ and every }n,
 \qquad
 u_{n+1}-u_n\le(1+\epsilon)\ell(u_n)\quad\text{eventually}.
\]
\end{proposition}
\leannote{Lean: \lproof{ErdosProblems/Erdos243/PaperCompleteR21/WindowAvoidance.lean}{201}{exists_avoiding_in_window}, \lproof{ErdosProblems/Erdos243/PaperCompleteR21/WindowAvoidance.lean}{423}{exists_slow_rise_avoiding_sequence}.}

\begin{proof}
The interval $[x,x+L)$ contains exactly $L$ integers.  Each lies in
$[1,\infty)$ and is smaller than $x+L$, so a
modulus exceeding $x+L$ divides none of them. Each of the $k$ remaining
moduli divides at most $L/m_i+1$ of them, so the covered count is at most
$L\theta+k<L$. This interval count itself does not use pairwise
coprimality; that condition is needed for the exact CRT proportions
in the later gap theorem.

For the second assertion, choose a tail of the family and reindex it so that
$\theta<\epsilon/(1+\epsilon)$.  If
$k(z)=\#\{i:m_i\le z\}$, the hypothesis $\ell(m_i)=i+O(1)$ gives
$k(z)\le\ell(z)+C$ for all large $z$ and some constant $C$.  Choose
\[
 (1-\theta)^{-1}<\rho<1+\epsilon,
 \qquad L(y)=\left\lceil\rho\ell(y)\right\rceil .
\]
Then $L(y)=O(\ell(y))=o(y)$, and the definition of $\ell$ gives
$\ell(y+1+L(y))=\ell(y)+o(1)$.  Consequently, for all large integers $y$,
\[
 \frac{k(y+1+L(y))}{1-\theta}
 \le \frac{\ell(y)+C+o(1)}{1-\theta}
 < \rho\ell(y)\le L(y),
\]
while $L(y)\le(1+\epsilon)\ell(y)$.  The first part also supplies an admissible $u_0$ beyond this threshold.
Apply it successively with $x=u_n+1$ and length $L(u_n)$, and choose
$u_{n+1}$ in the resulting window.  The sequence is strictly increasing,
avoids every retained modulus, and has the required rise bound.
\end{proof}

The small-increment sequence avoids only the retained moduli $m_i$,
$i\ge T$. Discarding the prefix makes their reciprocal sum small.
For the full family, Theorem~\ref{long243:res:gapconstant} instead gives
the coefficient $\prod_i(1-1/m_i)^{-1}>1$ for the increasing enumeration
of integers avoiding all the moduli.


Proposition~\ref{long243:res:coprimalitycap} concerns avoidance of whole
multiples, not coprimality to a composite modulus. The distinction matters:
no integer in $[2,5)$ is coprime to $30$, although none is divisible by $30$.
Thus this proposition alone does not disprove a coprime-walk extension of
Theorem~\ref{long243:res:barrier}. Proposition~\ref{long243:res:variablerise}
below gives a separate counterexample using sparse \emph{prime} moduli.
Under the additional scale $\ell(m_j)=j+O(1)$,
Theorem~\ref{long243:res:gapconstant} determines the exact maximal-gap
coefficient for these $\mathcal B$-free integers. Its proof uses a finite
sieve estimate and the Chinese remainder theorem, not an orbit
construction. Neither static construction supplies an exact
reciprocal-tail orbit.

\paragraph{Conditions on a possible counterexample.}
Together, these criteria give necessary conditions on a counterexample.
For every fixed $\delta\in(0,1)$, it must have
\[
 (1-\delta)\ell(C_n)\ \le\ -E_n\ =\ o(C_n)
 \qquad\text{at infinitely many }n.
\]
It must also have a divergent sum of relative increases, unbounded
negative error scaled by the previous maximum, and infinitely many record
jumps of the reduced numerator of size at least three, with $h_n=1$.
These are necessary conditions, not a construction of a counterexample. The residue-saturation lemma
shows precisely one limitation of the fixed-old-modulus method; it is not
an impossibility theorem about every future approach. A useful next target
is interaction between new prime support, three consecutive numerators and
the global bound on numerator growth.

\section{Descent when the error is nonnegative}\label{long243:sec:descent}

We return to zero-based indexing for the exact recurrences.
If $E_n\ge0$, the update $C_{n+1}=C_n-E_n$ makes $C_n$
nonincreasing. A nonincreasing sequence of natural numbers stabilises.
The next theorem records this elementary case before we turn to
negative errors.

\begin{theorem}[descent]\label{long243:res:descent}
Let $C,E:\N\to\N$ satisfy $C_{n+1}+E_n=C_n$ for every $n$.  Then $E_n=0$ for
all sufficiently large $n$.
\end{theorem}
\leannote{Lean: \lproof{ErdosProblems/Erdos243/ReciprocalTailRigidity.lean}{1843}{centeredState_eventually_zero}.}

\begin{proof}
The recurrence gives $C_{n+1}\le C_n$. Choose an index at which the
natural-number sequence reaches its minimum. Every later value equals
that minimum, and $C_{n+1}+E_n=C_n$ then gives $E_n=0$.
\end{proof}

\noindent The result is formalised as the
\lword{Erdos243/ReciprocalTailRigidity.lean}{1843}{centeredState_eventually_zero}{stabilisation
of the nonnegative error}, using the
\lword{Erdos243/ReciprocalTailRigidity.lean}{1825}{antitone_nat_eventually_constant}{eventual
constancy of a nonincreasing sequence}.

The hypothesis is exactly one-sidedness: $C$ and $E$ take values in $\N$.  If
$E_n<0$ infinitely often, the numerator rises at those indices and this
descent argument no longer applies. The next two sections first treat
constant and periodic negative errors.

\section{When the negative error is constant}\label{long243:sec:constant}

The exclusions in this section and the next concern infinite constant or
periodic negative magnitudes, and they do not bound the length of a finite
transient.  Arbitrarily long constant transients occur on genuine rational
orbits, by the flat-transient theorem of the working report on
Problem~\#243 \cite[``Flat-transient theorem'']{erdosproblemaday243}.

Suppose the error is negative with a constant magnitude, $E_n=-m$ for a fixed
$m>0$ and every $n$.  By
Proposition~\ref{long243:res:update} the numerator increases by $m$ at each step,
so $C_n=c+nm$ with $c=C_0$, and the definition of the error becomes a
\emph{shape equation}
\[
 D_n+m=(a_n-1)\,(c+nm) .
\tag{5.1}\label{long243:eq:shape}
\]
Together with $D_{n+1}=a_nD_n$ this is a closed system in $(a,D)$.
It has no infinite natural-number solution with every $a_n\ge2$.  The following instance shows how the failure happens; the
theorem then shows that it cannot be avoided.

\begin{example}[the shape equation running until it fails]\label{long243:ex:shape}
Take $m=c=1$, so that $C_n=1+n$ and the shape equation~\eqref{long243:eq:shape} reads
$D_n+1=(a_n-1)(1+n)$, and take $D_0=1$.  Each step is now determined: at $n=0$
the equation reads $2=(a_0-1)\cdot1$, so $a_0=3$, and $D_1=a_0D_0=3$; at $n=1$
it reads $4=(a_1-1)\cdot2$, so $a_1=3$, and $D_2=a_1D_1=9$; at $n=2$ it reads
$10=(a_2-1)\cdot3$, which has no integer solution, and the orbit stops.
Longer prefixes occur for other starting values. For $2\le a_0<5000$,
the exact residue search in Appendix~\ref{long243:app:residue} gives a
maximum of $17$ successful updates, hence $18$ multiplier values
including $a_0$.
\end{example}

\begin{theorem}[no constant negative magnitude]\label{long243:res:constant}
For any $m,c\in\N$ with $m>0$, there is no pair of sequences
$a,D:\N\to\N$ with $a_n\ge2$ for all $n$ satisfying
$D_{n+1}=a_nD_n$ and~\eqref{long243:eq:shape}. The same holds if the
shape equation only begins at some index.
\end{theorem}
\leannote{Lean: \lproof{ErdosProblems/Erdos243/PaperCompleteR21/NegativeMagnitudeExclusions.lean}{38}{no_constantNegative_shapeEquation}, \lproof{ErdosProblems/Erdos243/PaperCompleteR21/NegativeMagnitudeExclusions.lean}{52}{no_eventuallyConstantNegative_shapeEquation}.}

\begin{proof}[Proof when $c$ and $m$ are coprime]
Assume first $\gcd(c,m)=1$. Every multiplier must share a prime with
$m$, but each such prime can occur in at most one multiplier.

\emph{Every $a_j$ shares a prime with $m$.}  Suppose $\gcd(a_j,m)=1$.  Then
$m$ is invertible modulo $a_j$, so some $n$ has $c+nm\equiv0$, that is
$a_j\mid C_n$; and $a_j\mid D_n$ for every $n>j$, since $D$ is multiplicative
with $a_j$ among its factors.  Choosing such an $n$ beyond $j$
and reading~\eqref{long243:eq:shape} modulo $a_j$ gives $a_j\mid m$, contradicting
$\gcd(a_j,m)=1$ and $a_j\ge2$.

\emph{A prime divisor of $m$ occurs in at most one $a_j$.}  Suppose $p\mid m$
and $p\mid a_j$.  For $n>j$ we have $p\mid D_n$, so~\eqref{long243:eq:shape} gives
$(a_n-1)C_n\equiv0 \pmod p$.  Now $C_n=c+nm\equiv c$, and $p\nmid c$ because
$p\mid m$ and $\gcd(c,m)=1$; hence $p\mid a_n-1$, so $p\nmid a_n$.  Thus $p$
cannot divide any later multiplier.

This injects infinitely many multipliers into the finite set of prime
divisors of $m$, a contradiction.
\end{proof}

\noindent The case $\gcd(c,m)=1$ is the
\lword{Erdos243/ReciprocalTailRigidity.lean}{147}{no_scalePrimitiveConstantNegative_orbit}{exclusion when $c$ and $m$ are coprime}.  The special case $m=c=1$, worked through in
Example~\ref{long243:ex:shape}, where $m$ has no prime divisors at all and the first
fact is immediately contradictory, is recorded separately as the
\lword{Erdos243/ReciprocalTailRigidity.lean}{100}{no_normalizedConstantNegative_orbit}{case $c=m=1$}.

\begin{proof}[Removing the scale]
For general $c$, put $g=\gcd(c,m)$.  Equation~\eqref{long243:eq:shape} shows
$g\mid D_n$ for every $n$, so dividing $D$, $c$ and $m$ by $g$ leaves a system
of the same shape with coprime data, which the previous case excludes.
\end{proof}

\noindent The formal statements are the
\lword{Erdos243/ReciprocalTailRigidity.lean}{286}{no_constantNegative_orbit}{exclusion
at every scale} and the
\lword{Erdos243/ReciprocalTailRigidity.lean}{350}{no_eventuallyConstantNegative_orbit}{eventual
exclusion}. The latter follows by shifting the orbit to the first index
of constant error.

Theorem~\ref{long243:res:constant} excludes an error that is eventually constant
and negative, at every magnitude and every scale. Constancy fixes both
the possible prime divisors and the congruence used at later indices.
For varying magnitudes, even when their prime divisors belong to a fixed
finite set, this argument does not supply the same later congruence.
It therefore does not settle that case.

\section{When the negative error is periodic}\label{long243:sec:periodic}

The next case allows the magnitude to vary, provided it repeats.  Write
$e_n=-E_n>0$ for the magnitude, as in Section~\ref{long243:sec:state}, so that
$C_{n+1}=C_n+e_n$; suppose $e$ has period $h$, meaning $e_{n+h}=e_n$ for every
$n$, and that the numerator gains a fixed \emph{drift} $M>0$ over one period,
meaning $C_{n+h}=C_n+M$ for every $n$. In fact the update forces
$M=\sum_{j=0}^{h-1}e_j$: periodicity makes the sum over any $h$
consecutive indices the same. Thus $M>0$ follows from the positive
magnitudes; it is not an independent growth assumption.

At $h=1$ this says that the magnitude is constant and that $M$ is its value,
which is the situation of Section~\ref{long243:sec:constant};
Theorem~\ref{long243:res:constant} already excludes it without using
the pointwise bound $e_n<a_n$ in the proof below. The new case is $h\ge2$.

For a constant error, the proof used the finitely many prime divisors of
its magnitude. For a periodic error we use the prime divisors of the
increase $M$ over one period. The required divisibility facts are as
follows. Every multiplier $a_j$ divides each later denominator, by
\lword{Erdos243/ReciprocalTailRigidity.lean}{382}{base_dvd_denState_later}{persistence
of a multiplier}. If a prime divides both $D_n$ and $a_n$, the numerator
update makes it divide $C_{n+1}$, the
\lword{Erdos243/ReciprocalTailRigidity.lean}{402}{divisorLocksTailSucc}{one-step
divisibility implication}. Once it divides both numerator and denominator,
it divides every later numerator, denominator and error, by
\lword{Erdos243/ReciprocalTailRigidity.lean}{429}{divisorLock_persists}{persistence
of a common divisor}.

Periodicity carries this divisibility back to every phase. Suppose
$d\mid M$ and $d\mid a_i,a_j$ with $i<j$. Then $d\mid D_j$, and
$C_{j+1}=a_jC_j-D_j$ gives $d\mid C_{j+1}$. The common divisor
persists in both sequences and in every later error. For any phase $r$,
choose $k$ with $r+kh\ge j+1$; periodicity gives
$d\mid e_{r+kh}=e_r$. Taking $kh\ge j+1$ also gives
$d\mid C_{kh}=C_0+kM$, hence $d\mid C_0$. This is
\lword{Erdos243/ReciprocalTailRigidity.lean}{494}{repeatedDivisor_forces_periodicCommonScale}{the repeated-divisor implication}; it uses the
\lword{Erdos243/ReciprocalTailRigidity.lean}{461}{periodic_add_mul}{period relation}
and the
\lword{Erdos243/ReciprocalTailRigidity.lean}{476}{phaseDrift_add_mul}{increase over successive periods}, each iterated through whole periods.
A repeated prime can therefore be divided out of the entire system.
This explains the induction on $M$ in the proof.

\begin{theorem}[no periodic negative magnitude]\label{long243:res:periodic}
Let $a,D,C,e:\N\to\N$ with $a_n\ge2$, $e_n>0$ and $e_n<a_n$ for every $n$,
satisfying
\[
 D_{n+1}=a_nD_n,\qquad C_{n+1}=C_n+e_n,\qquad D_n+e_n=(a_n-1)C_n ,
\]
and suppose $e_{n+h}=e_n$ and $C_{n+h}=C_n+M$ for some $h>0$ and $M>0$.  This
is impossible.
\end{theorem}
\leannote{Lean: \lproof{ErdosProblems/Erdos243/PaperCompleteR21/NegativeMagnitudeExclusions.lean}{68}{no_periodicNegative_shapeEquation}.}

\begin{proof}
We use strong induction on the drift $M$. A prime divisor of $M$ that
recurs among the multipliers need not give an immediate contradiction.
It instead divides the whole orbit, allowing us to reduce $M$ by that
prime and apply the induction hypothesis.

Suppose first that no prime divisor of $M$ is a common divisor of $C_0$ and of
every magnitude.  The repeated-divisor implication applies in the
contrapositive: a prime divisor of $M$ occurring in two of the multipliers
would be exactly such a common divisor, so each prime divisor of $M$ occurs in
at most one multiplier.  On the other hand, every multiplier shares a prime with $M$.
Indeed, if $\gcd(a_j,M)=1$, choose $k\ge1$ so that
$a_j\mid C_j+kM=C_{j+kh}$. The same multiplier divides $D_{j+kh}$,
and periodicity gives $e_{j+kh}=e_j$. The shape equation at $j+kh$
therefore implies $a_j\mid e_j$, contrary to $0<e_j<a_j$.
Thus infinitely many multipliers would require distinct prime divisors
of the fixed integer $M$, a contradiction.

Otherwise some prime $p$ divides $M$, divides $C_0$, and divides every
magnitude.  Then $p\mid C_n$ for every $n$, since $C_{n+1}=C_n+e_n$ and both
summands on the right are divisible by $p$, and the shape equation
$D_n+e_n=(a_n-1)C_n$ then gives $p\mid D_n$ for every $n$.  Divide $D$, $C$ and
$e$ by $p$.  The multipliers are untouched, so $a_n\ge2$ and $e_n<a_n$ persist
and the magnitudes stay positive; the three recurrences are homogeneous in
$(D,C,e)$ and so survive; the period is still $h$; and the drift becomes
$M/p<M$.  The inductive hypothesis applies.
\end{proof}

\noindent The first of the two cases above is the
\lword{Erdos243/ReciprocalTailRigidity.lean}{548}{no_phasePrimitivePeriodicNegative_orbit}{exclusion when no prime of $M$ divides all phases}, the induction is the
\lword{Erdos243/ReciprocalTailRigidity.lean}{690}{no_periodicNegative_orbit}{exclusion
at every scale}, and the eventual form is the
\lword{Erdos243/ReciprocalTailRigidity.lean}{802}{no_eventuallyPeriodicNegative_orbit}{eventual
exclusion}.

The phase-by-phase proof uses $e_n<a_n$ to prevent a multiplier from
dividing its own nonzero phase magnitude. This is not a rescaling of the
orbit. Nevertheless, for an infinite exact orbit with periodic positive
magnitudes, the bound follows eventually from the other assumptions.
The shape equation gives $C_n>0$, and periodicity gives
\[
 C_n=\frac{M}{h}n+O(1),\qquad \sup_ne_n<\infty,
\]
and $D_0$ cannot be zero: otherwise all $D_n$ vanish and the shape
equation gives $e_n=(a_n-1)C_n\ge C_n$, contradicting these bounds.
Thus $D_n\ge2^nD_0$, and
\[
 a_n-1=\frac{D_n+e_n}{C_n}\longrightarrow\infty.
\]
Consequently $e_n<a_n$ eventually. Applying the eventual exclusion
therefore rules out periodic positive magnitudes without an independent
smallness assumption. The numbered theorem retains the pointwise bound
used by its checked phase-by-phase proof. The later bounded-negative
theorem also excludes this case.

\begin{remark}
On the rational-tail orbit of~\cite{koizumi2025}, the absolute error
satisfies $|E_n|<a_n$ eventually. Indeed, under the dictionary in
Section~\ref{long243:sec:priorwork}, Lemma~4(2) gives
$-C_n/2\le E_n<C_n/2$, hence $|E_n|\le C_n/2$. The proof of Lemma~4(3) gives
$C_{n+1}\le\tfrac32C_n$~\cite[pp.~11--12]{koizumi2025}, hence
$C_n\le C_N(3/2)^{n-N}$ beyond a pseudo-greedy starting index $N$.
On the other hand, convergence of $\sum1/a_n$ gives $a_n\to\infty$,
and $a_{n+1}/a_n^2\to1$ then gives $a_{n+1}\ge2a_n$ eventually.
Thus $|E_n|/a_n\to0$. On a negative-error tail this gives $e_n<a_n$.

For an abstract exact orbit, the same conclusion follows when
$a_n>1$, $C_n>0$ and $E_n/C_n\to0$ hold: the calculation in
Section~\ref{long243:sec:state} already derives the reciprocal series and
its growth from these hypotheses. Summability is not an independent
assumption in that setting. The bound $e_n<a_n$ is still only eventual,
whereas Theorem~\ref{long243:res:periodic} states it at every index;
its eventual form, cited above, is therefore the applicable one.
For a periodic magnitude, boundedness of $e_n$ and divergence of $a_n$
already suffice.
\end{remark}

\section{Bounded increases and coprimality to earlier moduli}\label{long243:sec:barrier}

Periodicity is still a strong assumption.  Removing it needs a different kind
of obstruction, and the one used here is a counting argument about where an
integer sequence with bounded upward steps must land.

\begin{lemma}[shifted blocks of consecutive multiples]\label{long243:res:crt}
Let $m_0,\ldots,m_{B-1}$ be pairwise coprime and at least $2$.  For every bound
there is a $t$ beyond it with $m_i\mid t+i$ for each $i<B$.
\end{lemma}
\leannote{Lean: \lproof{ErdosProblems/Erdos243/PaperCompleteR21/ForbiddenBlockCrossing.lean}{26}{exists_shiftedBlock_consecutiveMultiples}.}

\begin{proof}
Standard Chinese remaindering: solve $t\equiv-i \pmod{m_i}$ simultaneously,
then add multiples of $\prod_i m_i$ to pass the bound.
\end{proof}

\noindent The lemma is formalised as the
\lword{Erdos243/ReciprocalTailRigidity.lean}{839}{exists_shifted_consecutiveMultiples}{shifted
consecutive multiples}.

\noindent Lemma~\ref{long243:res:crt} matches each modulus $m_i$ to its own
multiple $t+i$ inside a window of $B$ consecutive integers. Matchings of a set
of integers to distinct multiples in an interval are the subject of
Erd\H{o}s Problem~\#650, solved by van Doorn, Li and Tang: for every set $S$
of $k$ positive integers, every open interval of length $2\max S$ contains
distinct multiples of at least $\min(k,\lceil 2\sqrt{k}\,\rceil)$ elements of
$S$, and this count is optimal \cite[Theorem~2.1]{vandoornlitang2026}. Their
extremal construction uses the same Chinese remaindering for moduli that need
not be pairwise coprime, with their Claim~3.2 supplying the divisibility of
residue differences by greatest common divisors that the generalised theorem
requires \cite[pp.~4--5]{vandoornlitang2026}. The first-crossing argument
below needs only the pairwise coprime case.

\begin{example}[a block of three consecutive multiples]\label{long243:ex:crt}
Take $B=3$ and $(m_0,m_1,m_2)=(3,4,5)$.  The congruences $t\equiv0\pmod3$,
$t\equiv3\pmod4$ and $t\equiv3\pmod5$ hold exactly when $t\equiv3\pmod{60}$.
At $t=3$ the block is $3,4,5$ with $3\mid3$, $4\mid4$ and $5\mid5$; at $t=63$
it is $63,64,65$ with $3\mid63$, $4\mid64$ and $5\mid65$.  A sequence of
natural numbers that starts at $0$, tends to infinity and rises by at most $3$
at each step cannot step over the block $\{3,4,5\}$: it has a first value at
least $3$, that value is at most $5$, and every member of $\{3,4,5\}$ is
divisible by one of the three moduli.
\end{example}

\begin{theorem}[bounded increases and coprimality to earlier moduli]\label{long243:res:barrier}
Let $u:\N\to\N$ tend to infinity with $u_{n+1}\le u_n+B$ for a fixed integer $B\ge1$.
Then $u$ cannot remain coprime to infinitely many fresh pairwise coprime
moduli: there is no family of pairwise coprime $m_i\ge2$, one for each index,
such that $\gcd(m_i,u_t)=1$ whenever $i<t$.
\end{theorem}
\leannote{Lean: \lproof{ErdosProblems/Erdos243/PaperCompleteR21/ForbiddenBlockCrossing.lean}{44}{no_boundedRise_coprimeToEarlierModuli}.}

The CRT block must lie above the initial segment where some moduli
are not yet available. No monotonicity of $u$ is assumed.

\begin{proof}
Choose $m_0,\ldots,m_{B-1}$ and use Lemma~\ref{long243:res:crt} to find
$t>\max(u_0,\ldots,u_B)$ with $m_i\mid t+i$ for $0\le i<B$.  Since
$u_n\to\infty$, there is a first $n>B$ with $u_n\ge t$.  Minimality and the
rise bound give
\[
 t\le u_n\le u_{n-1}+B<t+B.
\]
Hence $u_n=t+i$ for some $0\le i<B<n$, so $m_i\mid u_n$.  The inequality
$i<n$ now permits the avoidance hypothesis, which says
$\gcd(m_i,u_n)=1$; this contradicts $m_i\ge2$.  The argument uses the first
crossing of the CRT block, not monotonicity of $u$.
\end{proof}

Only unboundedness above is used to obtain this first crossing.
Divergence is stated because it holds for the reduced rational tail in
the application. This distinction also explains why the weighted proof
in Section~\ref{long243:sec:lcmrecords} works with an unbounded running
maximum, without assuming that its numerator tends to infinity.

\noindent The argument is formalised as
\lword{Erdos243/ReciprocalTailRigidity.lean}{903}{no_boundedRise_of_tailAvoidance}{the Chinese-remainder first-crossing consequence}.
Divergence forces $u$ to reach the forbidden block, and the uniform rise
bound prevents it from jumping over the block. A bound only along a
subsequence would leave the intervening steps unrestricted. Without a scale restriction on the moduli,
Proposition~\ref{long243:res:variablerise} gives an unbounded coprime walk
with rises $o(\log\log u_n)$.  The same forbidden-block crossing,
under a two-sided bound on the error, appears in the proof of Bado's
Theorem~5.1 \cite[pp.~4--5]{bado2026}, posted in September 2026.

In the application, $m_i$ is the $i$th multiplier. Its coprimality is
known only for later numerators, which explains the condition $i<t$.

The barrier applies once reduction removes no further common factor.
Call $(u,v)$ a \emph{reduced exact tail} with multipliers $a$ when
$\gcd(u_n,v_n)=1$ and
\[
 u_{n+1}+v_n=a_nu_n,\qquad v_{n+1}=a_nv_n .
\]
Here $u$ is the numerator and $v$ the denominator: the two recurrences are
those of Section~\ref{long243:sec:state}, with $u$ in the role of the numerator $C$
and $v$ in that of the denominator state $D$, and with coprimality imposed at
every index.

\begin{proposition}[reduced tails are pairwise coprime]\label{long243:res:reduced}
In a reduced exact tail, $a_n$ is coprime to $v_n$; the multipliers at distinct
indices are pairwise coprime; and every earlier multiplier is coprime to every
later numerator.
\end{proposition}
\leannote{Lean: \lproof{ErdosProblems/Erdos243/PaperCompleteR7/Reduction.lean}{16}{persistent_coprimality}.}

\begin{proof}
A common prime divisor of $a_n$ and $v_n$ would divide both
$u_{n+1}=a_nu_n-v_n$ and $v_{n+1}=a_nv_n$, contradicting their
coprimality. Hence $\gcd(a_n,v_n)=1$.
For $i<t$, the denominator recurrence gives $a_i\mid v_t$.
Combining this with $\gcd(a_t,v_t)=1$ proves that $a_i$ and $a_t$
are coprime; combining it with $\gcd(u_t,v_t)=1$ proves that $a_i$
and $u_t$ are coprime.
\end{proof}

\noindent These are the
\lword{Erdos243/ReciprocalTailRigidity.lean}{989}{reducedStep_coprime_currentFactor}{step
coprimality}, the
\lword{Erdos243/ReciprocalTailRigidity.lean}{1016}{reducedTail_pairwiseCoprime}{pairwise
coprimality}, and the
\lword{Erdos243/ReciprocalTailRigidity.lean}{1030}{reducedTail_wholeModulusAvoidance}{coprimality to each earlier multiplier}.  Feeding them to Theorem~\ref{long243:res:barrier} gives the form used
later: there is no reduced exact tail whose numerator tends to infinity with a
uniformly bounded upward increment, the
\lword{Erdos243/ReciprocalTailRigidity.lean}{1041}{no_boundedRise_reducedTail}{reduced
exclusion under bounded upward increments}, together with its eventual form, the
\lword{Erdos243/ReciprocalTailRigidity.lean}{1066}{no_eventuallyBoundedRise_reducedTail}{version with an eventual increment bound}.

\begin{example}[Sylvester's sequence as a reduced exact tail]\label{long243:ex:reduced}
Put $u_n=1$ and $v_n=D_n=1,2,6,42,1806,\ldots$ as in
Example~\ref{long243:ex:sylvester}, with multipliers $a_n=v_n+1=2,3,7,43,1807,\ldots$.
Then $\gcd(u_n,v_n)=1$, $u_{n+1}+v_n=1+v_n=a_nu_n$ and $v_{n+1}=a_nv_n$, so
this is a reduced exact tail, and Proposition~\ref{long243:res:reduced} returns the
classical fact that Sylvester's numbers are pairwise coprime.  Its numerator is
constant, so it satisfies every hypothesis of the exclusion just stated except
divergence; since the tail exists, that hypothesis cannot be dropped.
\end{example}

An arbitrary orbit of Section~\ref{long243:sec:state} need not be reduced.
To use the preceding argument, we must show that the common factor stops
changing. One uniform bound on the negative magnitudes at arbitrarily
late indices suffices: each gcd divides every later gcd, and at a negative
index it also divides the nonzero magnitude. Thus a bound at arbitrarily
late negative indices bounds the entire gcd sequence. The next proposition
justifies dividing by its eventual value.

\begin{proposition}[the tail gcd stabilises]\label{long243:res:gcdstab}
Let $(a,D,C)$ be an exact orbit of natural numbers, that is, a triple of $\N$-valued
sequences with $C_{n+1}+D_n=a_nC_n$ and $D_{n+1}=a_nD_n$, whose error is
$E_n=D_n-(a_n-1)C_n$. Suppose some fixed integer $B\ge1$ satisfies
$-B\le E_n<0$ at infinitely many indices. Then
$\gcd(C_n,D_n)$ is eventually constant, and beyond that index the orbit divided
by the stable gcd is a reduced exact tail.
\end{proposition}
\leannote{Lean: \lproof{ErdosProblems/Erdos243/PaperCompleteR7/Reduction.lean}{103}{gcd_stabilises_and_reduces}.}

\begin{proof}
First, $C_n>0$ at every index. If $C_n=0$, the nonnegative recurrence
forces $C_{n+1}=D_n=0$, and both sequences then vanish at all later
indices. This contradicts the existence of arbitrarily late negative
errors.

The two recurrences give $G_n\mid G_{n+1}$ for
$G_n=\gcd(C_n,D_n)$, as in
\lword{Erdos243/ReciprocalTailRigidity.lean}{1430}{tailGcd_dvd_succ}{the gcd divisibility relation}.
Also $G_n=\gcd(C_n,|E_n|)$, since $D_n=E_n+(a_n-1)C_n$;
at a negative index this is
\lword{Erdos243/ReciprocalTailRigidity.lean}{1587}{tailGcd_eq_gcd_negativeMagnitude}{the gcd identity with the negative magnitude}.
For any $n$, choose $t\ge n$ with $-B\le E_t<0$. Then
\[
 G_n\le G_t\le |E_t|\le B.
\]
The entire positive divisibility chain is therefore bounded and eventually
constant, by
\lword{Erdos243/ReciprocalTailRigidity.lean}{1605}{dvdChain_eventuallyConstant_of_cofinally_bounded}{stabilisation of a bounded divisibility chain}.
This is the supplied
\lword{Erdos243/ReciprocalTailRigidity.lean}{1636}{tailGcd_eventuallyConstant_of_cofinally_boundedNegative}{gcd stabilisation result}.
Dividing both sequences by the eventual value preserves the recurrences
and gives coprime numerator and denominator at every later index.
\end{proof}

Other negative errors may exceed $B$ in magnitude. Gcd stabilisation
uses the bounded witnesses, whereas the CRT application needs a bound
on \emph{every} upward increment. The supplied declaration proves
stabilisation; division by the stable gcd gives the reduced tail.

Vanishing relative error also gives sparsity of the strict gcd increases.
The next proposition asserts long finite intervals of constancy, not
constancy on an infinite tail.
For an exact orbit of natural numbers with $C_n>0$, put
\[
 G_n=\gcd(C_n,D_n),\qquad
 \Gamma(N)=\#\{0\le j<N:G_j<G_{j+1}\}.
\]

\begin{proposition}[vanishing relative error makes strict gcd changes sparse]
\label{long243:res:gcdsparse}
Let $a,C,D:\N\to\N$ satisfy $C_{n+1}+D_n=a_nC_n$ and $D_{n+1}=a_nD_n$ with
$C_n>0$, and put $E_n=D_n-(a_n-1)C_n$, $G_n=\gcd(C_n,D_n)$ and
$\Gamma(N)=\#\{0\le j<N:G_j<G_{j+1}\}$.  If $|E_n|/C_n\to0$, then
$\Gamma(N)=o(N)$.  Moreover, for every starting bound $B$ and block length
$L$, some $n\ge B$ satisfies
\[
  G_n=G_{n+1}=\cdots=G_{n+L}.
\]
\end{proposition}
\leannote{Lean: \lproof{ErdosProblems/Erdos243/PaperCompleteR7/Limits.lean}{97}{sparse_gcd_changes}.}

\noindent The first assertion is the
\lword{Erdos243/ReciprocalTailRigidity.lean}{2186}{tailGcd_strictGrowthCount_sublinear_of_normalizedVanishes}{sublinear
strict-growth theorem}; the second is the
\lword{Erdos243/ReciprocalTailRigidity.lean}{2202}{tailGcd_exists_arbitrarilyLate_constBlock_of_normalizedVanishes}{arbitrarily
late constant-block theorem}. The proof first converts vanishing relative
error into subexponential growth of $C_n$. Each strict divisibility increase
of $G_n$ contributes a factor of at least $2$, so
$2^{\Gamma(N)}G_0\le G_N\le C_N$. Taking logarithms gives
$\Gamma(N)=o(N)$. For $L=0$ the second assertion is immediate.
For $L\ge1$, if every block of $L$ successive transitions beyond $B$
contained a strict increase, disjoint such blocks would give
$\Gamma(N)\ge(N-B)/L-O(1)$, a contradiction.

The quantifiers matter.  Proposition~\ref{long243:res:gcdstab} uses a bound on negative
magnitudes at arbitrarily late indices and yields eventual constancy.  Proposition~\ref{long243:res:gcdsparse}
uses only vanishing relative error and yields arbitrarily late constant blocks of
each prescribed finite length. This proof does not establish eventual
constancy under that hypothesis, nor does it bound the negative magnitudes.

\section{A lower bound on the error forces eventual zero}\label{long243:sec:bounded}

The lower bound $E_n\ge-B$ has two uses in the proof: it bounds
negative errors at arbitrarily late indices to stabilise the gcd, and it
bounds every upward increment to apply Theorem~\ref{long243:res:barrier}.

The denominator recurrence is essential. For a scalar sequence,
$C_n=c+bn$ and $E_n=-b$, with positive integers $b,c$, satisfy both the
lower bound and vanishing relative error but never stabilise.
Theorem~\ref{long243:res:constant} excludes this as an exact orbit.
The different, summability-based argument of the next section needs
only the scalar update. On exact positive integer orbits with vanishing
relative error, both sufficient conditions are equivalent to eventual
zero error; neither is established from the unrestricted hypotheses.

\begin{theorem}[bounded negative part]\label{long243:res:bounded}
Let $a,C,D:\N\to\N$ and $E:\N\to\Z$ satisfy
\begin{enumerate}
\item $a_n>1$ and $C_n>0$ for every $n$;
\item the exact dynamics $C_{n+1}+D_n=a_nC_n$ and $D_{n+1}=a_nD_n$;
\item $E_n=D_n-(a_n-1)C_n$ for every $n$;
\item \emph{eventual strict centring}: $|E_n|<C_n$ for all large $n$;
\item \emph{eventually bounded negative part}: $-B\le E_n$ for all large $n$,
for some integer $B\ge0$;
\item \emph{vanishing relative error}: for every integer $K\ge1$ there is an $N$ with
$K\,|E_n|<C_n$ for all $n\ge N$.
\end{enumerate}
Then $E_n=0$ for all sufficiently large $n$.
\end{theorem}
\leannote{Lean: \lproof{ErdosProblems/Erdos243/ReciprocalTailRigidity.lean}{2360}{eventuallyBoundedNegativePart_eventually_zero}.}

\begin{proof}
Suppose the error is not eventually zero. Absorption implies that
$|E_n|\ge1$ after a sufficiently late centring threshold.
Condition~(6) therefore gives $C_n\to\infty$.

If negative indices were finite, Theorem~\ref{long243:res:descent}
would give eventual zero. Thus negative indices occur infinitely often,
and their magnitudes are bounded by $B$; in particular $B\ge1$.
Proposition~\ref{long243:res:gcdstab} makes $G_n=\gcd(C_n,D_n)$
stabilise to an integer $g>0$. The divided sequences
$u_n=C_n/g$, $v_n=D_n/g$ form a reduced exact tail, with
$u_n\to\infty$ and $u_{n+1}-u_n=-E_n/g\le B$.
Their coprimality properties from Proposition~\ref{long243:res:reduced}
now contradict Theorem~\ref{long243:res:barrier}.
\end{proof}

\noindent The theorem is
\lword{Erdos243/ReciprocalTailRigidity.lean}{2360}{eventuallyBoundedNegativePart_eventually_zero}{formalised with hypotheses~(4) and~(5) holding eventually}.
The Lean proof shifts past both thresholds and applies the
\lword{Erdos243/ReciprocalTailRigidity.lean}{2271}{boundedNegativePart_eventually_zero}{version with centring and the lower bound at every index}.
It uses
\lword{Erdos243/ReciprocalTailRigidity.lean}{1739}{tailState_tendsto_atTop_of_nonzero_normalizedVanishes}{divergence from vanishing relative error}
and the
\lword{Erdos243/ReciprocalTailRigidity.lean}{1754}{no_cofinallyBoundedNegative_of_normalizedVanishes}{exclusion for bounded upward increments and arbitrarily late bounded negative magnitudes}, via
\lword{Erdos243/ReciprocalTailRigidity.lean}{1661}{no_cofinallyBoundedNegative_of_tailTendsto}{its formulation with divergence as a hypothesis}.
Both exclusion statements require a uniform bound on all upward
increments. Bounded negative magnitudes at arbitrarily late indices
suffice for gcd stabilisation, but not for the jump bound in the CRT
argument.

\begin{corollary}\label{long243:res:cor}
Under the hypotheses of Theorem~\ref{long243:res:bounded}, together with
$C_{n+1}\ne0$ for all large $n$, the multipliers satisfy
$a_{n+1}=a_n^{2}-a_n+1$ for all sufficiently large $n$.
\end{corollary}
\leannote{Lean: \lproof{ErdosProblems/Erdos243/ReciprocalTailRigidity.lean}{2412}{boundedNegativePart_sylvesterNext_eventually}.}

\begin{proof}
Theorem~\ref{long243:res:bounded} gives eventual zero error.
Theorem~\ref{long243:res:eventual} then gives the Sylvester recurrence.
\end{proof}

Hypotheses~(1)--(3) specify the positive integer recurrences.
Condition~(6) says $|E_n|/C_n\to0$ and implies~(4) by taking $K=1$.
The abstract theorem lists the latter separately to expose its use in
propagating a zero error. For a rational reciprocal sum, the analytic
construction supplies these conditions. The only additional assumption
is~(5), the lower bound on the error. The constant and periodic exclusions
proved earlier apply in their own stated regimes; they do not supply
that bound for a general sequence.

Here is the precise comparison with~\cite{koizumi2025}, which also
shows how Theorem~\ref{long243:res:bounded} applies to
Problem~\ref{long243:res:problem}. Let
$(a_n)$ satisfy the hypotheses of Problem~\ref{long243:res:problem}.  After deleting a
finite prefix, Corollary~3 of~\cite[p.~9]{koizumi2025} makes the sequence the
pseudo-greedy expansion of its own reciprocal sum, and Lemma~4 there supplies
the integers $c_n,d_n,e_n$ of Section~\ref{long243:sec:priorwork}.  Hypothesis~(1) then
holds: $c_n$ is a positive integer by Lemma~4(1), and $a_n>1$ after a further
finite shift, because summability of $\sum1/a_n$ forces $a_n\to\infty$.
Hypothesis~(2) is the pair of recurrences $c_{n+1}=a_nc_n-d_n$ and
$d_{n+1}=a_nd_n$ of Lemma~4(2), and hypothesis~(3) is that lemma's
$a_n=(d_n-e_n)/c_n+1$ read as $e_n=d_n-(a_n-1)c_n$, which is the error $E_n$.
The centring range $-c_n/2\le e_n<c_n/2$ in the same lemma gives
hypothesis~(4) at every index, which is stronger than the eventual form used
here.  Hypothesis~(6) is the vanishing
of the gap sequence in Corollary~3, since $\varepsilon_n=e_n/c_n$.  The sole
hypothesis not supplied is~(5), the eventual bound on the negative part.
Neither the cited construction~\cite{koizumi2025} nor the argument here
derives that bound from growth and rationality alone. This proves the
conditional result, not Problem~\#243.

\section{Finite total relative increase}\label{long243:sec:mass}

Theorem~\ref{long243:res:bounded} bounds individual upward increments of
$C_n$. Here we instead assume that the sum of those increments divided by
$C_n$ converges. A Sylvester tail satisfies this because its numerator is
eventually constant. The integer sequences $C_n=n+1$ and $C_n=(n+1)^2$
do not satisfy it: their relative increases have a divergent harmonic sum,
although each tends to zero. Unlike the preceding argument, this proof
needs no denominator or growth hypothesis.

\begin{samepage}
\begin{theorem}[finite sum of relative increases]\label{long243:res:mass}
Let \(C_n\in\N_{>0}\) and \(E_n\in\Z\) satisfy \(C_{n+1}=C_n-E_n\).
If
\[
 \sum_{n=0}^{\infty}\frac{(-E_n)_+}{C_n}<\infty,
\]
then \(E_n=0\) eventually. For an exact reciprocal-tail orbit this implies
\(a_{n+1}=a_n^2-a_n+1\) eventually. No hypothesis $|E_n|/C_n\to0$ is required.
\end{theorem}
\leannote{Lean: \lproof{ErdosProblems/Erdos243/PaperCompleteR7/Frontier.lean}{84}{finite_negative_mass_paper}.}
\end{samepage}

\begin{proof}
Put \(\delta_n=(-E_n)_+/C_n\). The update gives
\(C_{n+1}\le C_n(1+\delta_n)\), so
\[
 C_N\le C_0\prod_{n<N}(1+\delta_n)
 \le C_0\exp\!\left(\sum_n\delta_n\right).
\]
Choose an integer upper bound \(K\) for \(C_n\). Each strict rise of the
integer sequence contributes at least \(1/K\) to \(\sum_n\delta_n\), so
there are only finitely many rises. The remaining positive integer sequence
is nonincreasing and therefore stabilises, forcing \(E_n=0\). The recurrence
follows from Theorem~\ref{long243:res:eventual}, since \(C_n>0\).
\end{proof}

The same proof allows any positive real $C_0$, provided $E_n\in\Z$:
after the finitely many rises, the bounded positive values lie in the
finite set $(C_0+\Z)\cap(0,K]$, so they stabilise. The linked formal
statement retains integer-valued $C_n$.

The two indispensable features of this argument are positivity and
discrete increments. With $C_n=1+1/(n+1)$ and
$E_n=1/((n+1)(n+2))$, the numerator decreases forever with zero
relative-increase sum, but the errors are not integers. With $C_n=-n-1$
and $E_n=1$, the increments are integers but positivity fails.

\noindent The scalar conclusion is
\lword{Erdos243/SparseResetRecovery.lean}{612}{eventually_zero_of_summable_negativeRelativeMass_scalar}{the eventual vanishing for a positive integer sequence};
its exact-orbit consequence is
\lword{Erdos243/SparseResetRecovery.lean}{690}{sylvesterNext_eventually_of_summable_negativeRelativeMass_scalar}{the scalar recurrence from a convergent sum}.
The separate release contains a
\href{https://github.com/wcook04/plectis-erdos-lean/blob/52f29ad173b04e3bac941b3663f2b9aebe5de0bb/ErdosProblems/Erdos243/PaperCompleteR8/CanonicalNegativeMass.lean\#L19}{formulation for the rational-tail numerators using a convergent sum}, outside the main-repository build identified here. Summability remains
an assumption; the elementary scalar implication needs neither growth nor
vanishing relative error.

For the gap sequence of the pseudo-greedy expansion, the same criterion, with
the same product bound and integer descent, appears in the Erd\H{o}s Problem a
Day working report on Problem~\#243, dated 12 August 2026
\cite[``A global termination criterion'']{erdosproblemaday243}. Bado's
Theorem~11.1 also accounts for the factors removed in LCM clearing
\cite[Thm.~11.1 and Remark~11.3, p.~9]{bado2026}. In the notation of
Section~\ref{long243:sec:lcmrecords}, the update gives
\[
 \log U_N\le\log U_0+
 \sum_{n<N}\frac{(-E_n)_+}{C_n}-\sum_{n<N}\log\rho_n.
\]
His sufficient condition is an upper bound on the difference of these
two sums. The subtractive term records the decrease caused by a repeated
factor in the denominator. It may offset some relative increases; the
hypothesis does not require either sum to converge separately. Unlike the
scalar criterion above, this argument also uses the pseudo-greedy limit:
bounded $U_n$ and $V_n/U_n=E_n/C_n\to0$ make the integer $V_n$ eventually
zero, after which $\rho_nU_{n+1}=U_n$ gives eventual constancy.

Corollary~3~\cite[p.~9]{koizumi2025} and Lemma~4
~\cite[pp.~11--12]{koizumi2025} supply the same positive integer update
after a finite restart. Theorem~\ref{long243:res:mass} therefore applies
if the relative-increase sum converges. This is the sole additional
hypothesis, asked for in Problem~\ref{long243:res:masshyp}; the cited
construction does not establish it.

\begin{example}[the product bound at a geometric rate]\label{long243:ex:mass}
Suppose $C_0=100$ and $(-E_n)_+/C_n\le2^{-n-1}$ for every $n$, so that the
sum of relative increases is at most $1$; the constraint at $n=0$ still permits
$E_0$ as negative as $-50$.  Since $1+x\le e^{x}$,
\[
 C_N\le100\prod_{n<N}\bigl(1+2^{-n-1}\bigr)\le100e<272
\]
for every $N$. For all sufficiently large $n$, we have
$2^{-n-1}<1/272$. A strict rise would then force
$(-E_n)_+/C_n\ge1/C_n>1/272$, a contradiction. Thus $C_n$ is eventually
nonincreasing; as a positive integer sequence it stabilises, and $E_n=0$
thereafter. The constant $272$ comes from the total mass and not from any single magnitude,
which is why the hypothesis of Theorem~\ref{long243:res:mass} is summability of
$(-E_n)_+/C_n$ and not a bound on it.
\end{example}

\section{Complements and further questions}\label{long243:sec:open}

Absorption, descent and the two finiteness criteria give the following
necessary conditions. Constant or periodic patterns only along the
negative indices of a mixed-sign sequence are not covered by the
all-negative exclusions. None of these results resolves Problem~\#243.

\begin{proposition}[necessary conditions on a counterexample]\label{long243:res:frontier}
For the integer tail attached to any counterexample to
Problem~\ref{long243:res:problem},
\[
 E_n\ne0\quad\hbox{eventually},\qquad
 \frac{|E_n|}{C_n}\longrightarrow0,
\]
and
\[
 \limsup_{\substack{n\to\infty\\E_n<0}}(-E_n)=\infty,
 \qquad
 \sum_{n=0}^{\infty}\frac{(-E_n)_+}{C_n}=\infty .
\]
In particular, negative indices occur infinitely often.  Thus the remaining
regime consists of unbounded negative errors along exact reciprocal tails for which
the sum of relative increases diverges.
\end{proposition}
\leannote{Lean: \lproof{ErdosProblems/Erdos243/PaperCompleteR7/Frontier.lean}{151}{canonical_frontier}.}

\begin{proof}
Koizumi's construction from a rational reciprocal tail gives strict centring and vanishing relative error after a finite shift; the strict-centring hypothesis of
Theorem~\ref{long243:res:bounded} also follows from vanishing relative error with $K=1$.
Absorption then makes $E_n\ne0$ eventually,
since eventual zero would give the Sylvester recurrence.  Descent excludes an
eventually nonnegative error.  Theorem~\ref{long243:res:bounded} excludes a bounded
negative part, and Theorem~\ref{long243:res:mass} excludes convergence of the sum of relative increases.  These statements are unchanged by deleting a finite prefix.
\end{proof}

\noindent Lean checks Proposition~\ref{long243:res:frontier} as
\href{https://github.com/wcook04/plectis-erdos/blob/3d6d938d696fed0fb71dd55115a18a73738ff223/lean/ErdosProblems/Erdos243/PaperCompleteR7/Frontier.lean#L151}{the necessary conditions on a counterexample},
from the rational reciprocal sum, the quadratic growth limit, and the failure
of eventual Sylvester behaviour.  The divergent sum of relative increases is
recorded there as divergence of the partial sums.

Proposition~\ref{long243:res:frontier} does not assert infinitely many
prime divisors of the error magnitudes. Nor are its numerical conditions
alone a reformulation of the problem: the numerator and denominator
must satisfy the exact recurrences. With those recurrences, positivity
and vanishing relative error, Section~\ref{long243:sec:state} supplies
the reciprocal-series realisation. The following exact-orbit question
is therefore equivalent to the original problem.

\begin{problem}[unbounded negative errors with divergent relative sum]\label{long243:res:excursions}
Exclude, or construct, an exact orbit of natural numbers $(a,D,C)$ with $a_n>1$ and
$C_n>0$, whose multipliers satisfy $\lim_n a_{n+1}/a_n^{2}=1$ and whose error satisfies vanishing relative error, and which is
negative infinitely often with magnitudes unbounded along that infinite set and
\[
 \sum_n\frac{(-E_n)_+}{C_n}=\infty .
\]
Such an orbit fails the bounded-negative-part and summable-relative-increase
hypotheses.
By Proposition~\ref{long243:res:frontier} the integer tail of a counterexample is
such an orbit after deletion of a finite prefix, so an exclusion would settle
Problem~\#243. Conversely, a global orbit with the displayed properties
has $C_n/D_n\to0$ by Section~\ref{long243:sec:state}. The
\lword{Erdos243/ReciprocalTailRigidity.lean}{1955}{reciprocalSeries_hasSum_of_tailRatio_tendsto_zero}{reciprocal-series
realisation} then identifies its reciprocal sum as $C_0/D_0$.
After deletion of a finite prefix the multipliers are strictly increasing,
and the infinitely many negative errors rule out a Sylvester tail.
It is therefore a counterexample to Problem~\#243. Finite admissible
prefixes alone do not provide such a construction.
\end{problem}

\begin{remark}
The hypotheses in Problem~\ref{long243:res:excursions} beyond the state recurrence are
not decorative: the recurrence alone admits unbounded divergent-mass
excursions.  Put
\[
 C_n=2^{n},\qquad b_0=2,\qquad b_{n+1}=\tfrac12 b_n(b_n+2),\qquad
 a_n=b_n+2,\qquad D_n=b_nC_n ,
\]
so that $(b_n)=2,4,12,84,3612,\ldots$ and $(a_n)=4,6,14,86,3614,\ldots$.  Every
$b_n$ is a positive even integer, since $b_n=2k$ gives $b_{n+1}=2k(k+1)$, so all
four sequences are $\N$-valued.  The orbit is exact:
$C_{n+1}+D_n=2^{n}(2+b_n)=a_nC_n$ and
$D_{n+1}=b_{n+1}2^{n+1}=b_n(b_n+2)2^{n}=a_nD_n$.  Its error is
\[
 E_n=D_n-(a_n-1)C_n=b_n2^{n}-(b_n+1)2^{n}=-C_n ,
\]
so the negative magnitudes $2^{n}$ are unbounded and $\sum_n(-E_n)_+/C_n$
diverges.  What fails is the centring: $|E_n|=C_n$ at every index, so strict
centring fails at the boundary and vanishing relative error fails outright.  The
multipliers satisfy $a_{n+1}=\tfrac12(a_n^{2}-2a_n+4)$, so
$a_{n+1}/a_n^{2}\to\tfrac12$ and the growth hypothesis fails as well.  The
example therefore says nothing about Problem~\ref{long243:res:problem}, and it does not
isolate the centring hypotheses: $E_n=-C_n$ at every index forces that
multiplier recurrence, so growth fails with them.  It is recorded only to show
that the state recurrence alone is not enough.
\end{remark}

\paragraph{A scalar profile need not satisfy the recurrences.}
The example $C_n=n^2+1$, $E_n=-(2n+1)$ satisfies the numerator update,
$|E_n|/C_n\to0$, $\log C_n=o(n)$ and $\sum_n(E_n/C_n)^2<\infty$.
It is not an exact reciprocal-tail orbit. Indeed, $C_2=5$, $C_3=10$
and $C_4=17$. The first numerator update forces $5\mid D_2$; the
multiplicative denominator update gives $5\mid D_3$; the next numerator
update then forces $5\mid C_4$, a contradiction. This is the failed-route
example referred to at the end of the short note: subexponential size
and a small relative error do not supply arithmetic compatibility.

\subsection{Growth of the factors repeated in the denominator}

To compare repeated prime factors in the denominator with its total
size, define
\[
 L_0=D_0,\qquad L_{n+1}=\lcm(L_n,a_n),\qquad M_n=\frac{D_n}{L_n}.
\]
Since $D_n=D_0\prod_{j<n}a_j$, the quotient is integral and
$M_nL_n=D_n$.  The product counts every occurrence of a prime, whereas the least
common multiple keeps only its largest exponent. The quotient $M_n$
records the remaining factors. Problem~\ref{long243:res:lcmheight}
below asks whether failure of the Sylvester recurrence would force this
quotient to grow exponentially, contradicting its known subexponential
upper bound. We first record a local estimate for cancellation during
a recovery interval; it does not by itself give that global lower bound.

\subsection{Cancellation before a numerator regains its former value}

The LCM quotient $M_n$ need not equal the common divisor $G_n$ removed
by reduction to lowest terms. Since $M_n$ divides both $C_n$ and $D_n$,
we have $M_n\mid G_n$, but equality is not assumed. The following
estimate concerns the successive reduction factors, over an interval
where the reduced numerator regains its starting value.
In the calculation below, $h_n$ is the factor removed at step $n$,
$c_n$ is a positive integer with $c_n^2\mid h_n$, and $\tilde e_n$
denotes the signed reduced error, as in Section~\ref{long243:sec:records}.
For the estimate itself, let $u,h,c:\N\to\N$ and $\tilde e:\N\to\Z$ satisfy
\[
 h_nu_{n+1}=u_n-\tilde e_n,
 \qquad u_n>0,
 \qquad h_n>0.
\]
Fix $K,r,L\in\N$ with $K,L>0$, assume
\[
 K|\tilde e_{r+i}|<u_{r+i}\quad(0\le i<L),
 \qquad u_r\le u_{r+L},
\]
and let $R$ be a finite family of moduli.  Suppose every $m_q$, $q\in R$,
divides $\prod_{r\le n<r+L}c_n$, and that
$c_n^2\mid h_n$ throughout the interval.  Then
\begin{equation}\label{long243:eq:repair-entropy}
 K^L\lcm(m_q:q\in R)^2<(K+1)^L.
\end{equation}
To see this, each relative-error bound gives
$h_nu_{n+1}<(1+1/K)u_n$. Multiplication over the interval and
$u_r\le u_{r+L}$ yield
\[
 \prod_{r\le n<r+L}h_n<(1+1/K)^L.
\]
The least common multiple of the $m_q$ divides the product of the
$c_n$, and $c_n^2\mid h_n$ at every step. Its square therefore divides,
and is at most, the positive product on the left. Multiplying by $K^L$
proves the claimed bound.
This is the
\lword{Erdos243/RepairEntropy.lean}{48}
      {repairedFamily_recovery_energy_divisionFree}
      {bound on cancellation over a recovery interval}.

The square in~\eqref{long243:eq:repair-entropy} comes from the assumption
$c_n^2\mid h_n$. If the least common multiple of the chosen moduli
is at least $2^{|R|}$, then
\[
 K^L4^{|R|}<(K+1)^L.
\]
The formal consequence is
\lword{Erdos243/RepairEntropy.lean}{78}
      {repaired_card_bound_of_independent}
      {the bound on the number of independent moduli}.
Taking logarithms gives the explicit bound
\[
 \frac{|R|}{L}<\frac{\log(1+1/K)}{\log4}
 \le\frac1{K\log4}.
\]
Thus the number of these moduli is small relative to the recovery length
when the relative-error bound is small. This statement requires the lower bound
$2^{|R|}$ on their least common multiple; it does not apply to an arbitrary
family with repeated prime factors.

For a fixed number of steps, the same estimate has a simpler consequence:
a sufficiently late recovery cannot include any cancellation.  If
$K|\tilde e_n|<u_n$ eventually for every $K$, then for each fixed $L>0$ there is an
$N$ such that every recovery $u_r\le u_{r+L}$ with $r\ge N$ satisfies
\[
 \prod_{0\le i<L}h_{r+i}=1.
\]
Indeed, choose $K$ so large that $(1+1/K)^L<2$ and apply the same
product estimate beyond its error threshold. The positive integer
product is then smaller than $2$, so it equals $1$.
The same $K$ works for every length $1\le j\le L$, since
$(1+1/K)^j\le(1+1/K)^L<2$. Thus, after a sufficiently late step with
$h_n>1$, the numerator cannot regain its starting value at any of the
next $L$ indices. This includes a recovery followed by another fall
before the last endpoint. The fixed-length conclusion is formalised as
\lword{Erdos243/RepairEntropy.lean}{107}
      {eventually_recoveryPayment_eq_one_of_fixedLength}
      {absence of late cancellation during recovery intervals of fixed length}.
The theorem does not bound a recovery length that varies with $r$, nor does it
supply a lower bound for the least common multiple of the chosen moduli.
Those are the missing inputs needed to turn~\eqref{long243:eq:repair-entropy} into a
global contradiction.

\begin{problem}[growth of repeated denominator factors]\label{long243:res:lcmheight}
For every rational-tail orbit satisfying the hypotheses but not the conclusion of Problem~\ref{long243:res:problem},
must
\[
 \limsup_{n\to\infty}\frac{\log M_n}{n}>0?
\]
Equivalently, must there be a $K\ge1$ for which
\[
 2^n\le M_n^K
\]
at infinitely many indices?
\end{problem}

The two displayed formulations are equivalent up to changing the
positive constant; the second is not a weaker target.
Section~\ref{long243:sec:lcmrecords} already gives
$1\le M_n\le C_n$, hence $\log M_n/n\to0$ on every orbit under
consideration. A positive answer would therefore be a contradiction,
not a further compatible necessary condition. It would prove
Problem~\#243. The missing step is to force exponential growth of $M_n$
from failure of the Sylvester recurrence.

There is also a more local-looking question whose content is nevertheless the
entire prefix.  Put $A_n=\prod_{j<n}a_j$, so $D_n=D_0A_n$.  From
\[
 D_0A_n=(a_n-1)C_n+E_n
\]
one immediately obtains
\[
 \gcd(A_n,a_n-1)\mid E_n.
\]
The checked
\lword{Erdos243/ReciprocalTailRigidity.lean}{2170}{tailState_subexponential_of_normalizedVanishes}{tail-height estimate}
and vanishing relative error make $|E_n|$ subexponential in $n$.

\begin{problem}[common factors of the prefix and $a_n-1$]\label{long243:res:prefixgcd}
In every rational-tail orbit satisfying the hypotheses but not the conclusion of Problem~\ref{long243:res:problem}, is
\[
 \limsup_{n\to\infty}
 \frac{\log\gcd(A_n,a_n-1)}{n}>0?
\]
Equivalently, do there exist $\eta>0$ and infinitely many $n$ such that
$\gcd(A_n,a_n-1)\ge e^{\eta n}$?
\end{problem}

\noindent By Proposition~\ref{long243:res:frontier}, the error is nonzero
eventually. The displayed divisibility therefore bounds
$\gcd(A_n,a_n-1)$ by $|E_n|$ at every sufficiently late index.
A positive answer contradicts the subexponential bound on that error.
Arbitrarily long locally admissible blocks do
not answer this question: the gcd uses the complete prefix.

\subsection{The direct analytic question}

\begin{problem}[summability from rationality]\label{long243:res:masshyp}
Let $(a_n)$ satisfy the hypotheses of Problem~\ref{long243:res:problem},
and let $(D_n,C_n,E_n)$ be its integer tail. Must
\[
 \sum_{n=0}^{\infty}\frac{(-E_n)_+}{C_n}<\infty ?
\]
\end{problem}

\noindent An affirmative answer closes Problem~\#243 by
Theorem~\ref{long243:res:mass}. The summands are nonnegative and tend to
zero, so $\log(1+t)$ is comparable to $t$ at their values. Convergence
is therefore equivalent to boundedness of the partial products
$\prod_{n<N}(1+(-E_n)_+/C_n)$, the same criterion in multiplicative form.
A negative answer requires a sequence satisfying the full growth and
rationality hypotheses. A locally admissible state orbit is insufficient.
By Theorem~\ref{long243:res:weightedrecord}, it is enough to establish the
finiteness of~\eqref{long243:eq:weightedgrowth} for one nonincreasing weight
with divergent integral and one fixed $B\ge0$. The equivalent sum in
Theorem~\ref{long243:res:weightedrecord} uses only steps where the LCM
numerator exceeds its previous maximum. It subtracts $B$ from each upward
jump, takes the positive part and weights it by $f(U_n)$.

\begin{remark}
One further question is not left open here, since it is answered
in~\cite{koizumi2025}: whether vanishing relative error follows from
$a_{n+1}/a_n^{2}\to1$ and rationality.  It does.  After
deleting a finite prefix, Koizumi's Corollary~3
~\cite[p.~9]{koizumi2025} makes the sequence the pseudo-greedy expansion of its
own reciprocal sum and gives $\varepsilon_n\to0$ under exactly the hypotheses of
Problem~\ref{long243:res:problem}; rationality is not needed for that step, only
summability of $\sum1/a_n$.  The matched-index dictionary in
Section~\ref{long243:sec:priorwork} gives $\varepsilon_n=E_n/C_n$
on the restarted tail. For a rational sum, Lemma~4(2)
~\cite[pp.~11--12]{koizumi2025} also gives
$-c_n/2\le e_n<c_n/2$ at every index of that expansion. Thus
$|E_n|<C_n$ holds beyond the restart, as required by
Theorem~\ref{long243:res:bounded}; it need not hold at every original
index. Eventual strict centring also follows directly from vanishing
relative error with $K=1$.  What remains unsupplied is hypothesis~(5), the bound
on the negative part, which is why the theorem is still conditional.
\end{remark}

\subsection{What changes when upward increments are not bounded}

\begin{proposition}[small increases when the prime moduli are sparse]\label{long243:res:variablerise}
There exist strictly increasing primes $p_i$ and a strictly increasing
positive integer sequence $u_n\to\infty$ such that $\gcd(u_n,p_i)=1$ for
all $i,n$ and
\[
 u_{n+1}-u_n=O\bigl(\sqrt{\log\log(u_n+e^e)}\bigr)
             =o\bigl(\log\log(u_n+3)\bigr).
\]
Thus the bounded-increment hypothesis of
Theorem~\ref{long243:res:barrier} cannot be replaced by an
$o(\log\log u_n)$ bound without a quantitative restriction on the
moduli.
\end{proposition}
\leannote{Lean: \lproof{ErdosProblems/Erdos243/PaperCompleteR21/WindowAvoidance.lean}{862}{exists_sparse_prime_coprime_sequence}.}

\begin{proof}
Choose increasing primes
$p_i>\max\{\exp(\exp((i+2)^2)),2^{i+3}\}$, for $i\ge0$.
Primes of arbitrarily large size suffice; no distribution theorem is used.
Then $\theta=\sum_i1/p_i<1/4$ and
$k(z)=\#\{i:p_i\le z\}\le\sqrt{\log\log z}$ whenever $z$ is large.
For integer $x$ put
$L(x)=\lceil4\sqrt{\log\log(x+e^e)}+8\rceil$.
Since $L(x)=o(x)$, eventually $L(x)>k(x+L(x))/(1-\theta)$.
The elementary window bound of
Proposition~\ref{long243:res:coprimalitycap} leaves an integer in
$[x,x+L(x))$ divisible by no $p_i$. The set of such integers is therefore
unbounded. Enumerate it increasingly as $(u_n)$ and apply the same bound
with $x=u_n+1$ to obtain $u_{n+1}-u_n\le L(u_n+1)$.
For prime moduli nondivisibility is exactly coprimality, proving every
claim. The moduli are deliberately much sparser than the scale used in
the next theorem.
\end{proof}

We now keep the full family of moduli, rather than discarding a prefix
as in Proposition~\ref{long243:res:coprimalitycap}. The resulting
coefficient depends on $\sigma=\prod_j(1-1/m_j)$. For a finite prefix,
the corresponding product is exactly the proportion of admissible
residue classes by the Chinese remainder theorem. This explains the
constant in the statement before the limiting argument is made.

\begin{theorem}[largest gaps between integers avoiding given multiples]\label{long243:res:gapconstant}
Let $m_0<m_1<\cdots$ be pairwise coprime integers at least $2$ with
$\ell(m_j)=j+O(1)$, where $\ell(x)=\log_2\log_2\max(4,x)$.
Let $\sigma=\prod_j(1-1/m_j)>0$ and enumerate the positive integers
divisible by no $m_j$ in increasing order as $(u_n)$. Then
\[
 \limsup_{n\to\infty}\frac{u_{n+1}-u_n}{\ell(u_n)}=\sigma^{-1}.
\]
This is a statement about avoidance of multiples of whole moduli. It is not a statement
about coprimality to composite $m_j$, or about integer tail orbits.
\end{theorem}
\leannote{Lean: \lproof{ErdosProblems/Erdos243/PaperCompleteR21/MaximalGapConstant.lean}{1041}{maximal_gap_limsup_eq_inv_sigma}.}

\begin{proof}
Write $k(z)=\#\{j:m_j\le z\}=\ell(z)+O(1)$. For a fixed prefix length
$T$, set
\[
 M_T=\prod_{j<T}m_j,\qquad
 \sigma_T=\prod_{j<T}(1-1/m_j),\qquad
 \theta_T=\sum_{j\ge T}1/m_j.
\]
Here $\theta_T\to0$. By the Chinese remainder theorem, avoiding the first
$T$ whole moduli selects precisely the fraction $\sigma_T$ of the
residues modulo $M_T$.

For the upper bound, an integer interval $[x,x+L)$ contains at least
$\sigma_TL-M_T$ integers avoiding that prefix. The remaining moduli
cover at most $L\theta_T+k(x+L)$ integers. Choose $T$ large enough that
$\sigma_T>\theta_T$, and then any
$c>(\sigma_T-\theta_T)^{-1}$.
For $L=\lceil c\ell(x)\rceil$, the number left is positive for all
sufficiently large $x$, since $k(x+L)=\ell(x)+O(1)$.
Thus the set being enumerated is unbounded, and every sufficiently late
such interval meets it. Applying this at $x=u_n+1$ gives
\[
 \limsup_n\frac{u_{n+1}-u_n}{\ell(u_n)}
 \le(\sigma_T-\theta_T)^{-1}.
\]
Letting $T\to\infty$ gives the upper bound $\sigma^{-1}$.

For the lower bound, fix $T$ and let the integer length $L$ tend to
infinity. Among the offsets $0\le j<L$, precisely
$K_L=\sigma_TL+O(M_T)$ avoid the first $T$ moduli.
Assign to these offsets distinct moduli
$m_T,\ldots,m_{T+K_L-1}$. Solve simultaneously
$x\equiv0\pmod{M_T}$ and $x\equiv-j$ modulo the modulus assigned to $j$.
Put $Q_L=\prod_{i<T+K_L}m_i$ and choose this solution in $[Q_L,2Q_L)$.
Every integer of $[x,x+L)$ is then divisible by a modulus: either an old
one or its assigned new one. Let $u$ be the greatest admissible integer
below $x$; its successor is at least $x+L$, so the gap is at least $L$.
The scale assumption gives
\[
 \ell(2Q_L)\le T+K_L+O(1),
\]
because $\log_2m_i=2^{i+O(1)}$ and these upper bounds sum geometrically.
Since $u<2Q_L$, the ratio for this gap is at least
$L/(T+K_L+O(1))$, tending to $\sigma_T^{-1}$.
These $u$ tend to infinity: $x\ge Q_L\to\infty$ and admissible integers
are unbounded. Hence the limsup is at least $\sigma_T^{-1}$ for every
$T$. Let $T\to\infty$ to finish.
\end{proof}

For the Fermat moduli $m_j=2^{2^j}+1$, pairwise coprimality and the
finite-product identity
\[
 \prod_{j=0}^{N}\left(1-\frac1{2^{2^j}+1}\right)
 =\frac{2^{2^{N+1}-1}}{2^{2^{N+1}}-1}
\]
give $\sigma=1/2$, so the coefficient is exactly $2$.
These results describe gaps between integers avoiding fixed moduli.
They do not construct multipliers or numerator and denominator sequences
satisfying the exact recurrences. In particular, the static proportion
$\sigma$ must not be substituted for $\varphi(v_T)/v_T$: the first
excludes multiples of the whole moduli, while the second excludes every
prime factor of the reduced denominator.

\subsection*{Formalisation targets}

\paragraph{Erd\H{o}s--Straus.}
Formalise Theorem~3 of~\cite{erdosstraus1964}, with its prefix-LCM quotient
and correctly indexed growth factor, under the published analytic hypotheses.

\paragraph{Duverney.}
Formalise the absolute-convergence form of Corollary~3.2
of~\cite{duverney2001}, including its signed numerators, and recover the
all-positive specialisation used for comparison here. For the printed
condition interpreted as mere signed convergence, first justify the
nonvanishing-product step discussed in
Section~\ref{long243:sec:priorwork}. That interpretation is not treated
as an already verified formalisation target.

\pdfbookmark[1]{Statements and declarations}{statements-declarations}
\subsection*{Statements and declarations}

\paragraph{Artefact and data availability.} The
\href{\sourceurl}{pinned formal-source revision} contains the Lean sources, the
fixed toolchain and library manifest for the formalised results. Each
declaration establishes its stated implication. Results that also need
written arguments are identified separately in the manuscript.

The supplied source index records a successful build of the listed
main-repository modules at revision \texttt{6b78209ab63a}, using the
identical-file evidence described there. The links in this manuscript
retain their original pins, including \texttt{\commitshort}; this build
record is not a new verification of every linked revision. No new Lean run or transitive-axiom audit is claimed here, and the separate
release is not covered as a whole. A source link is not a build receipt.
The external formal-conjecture entry states the problem rather than proving
it. The ordinary proofs and their exposition require mathematical review
independently of the recorded checks.

\paragraph{Funding and competing interests.} This work received no external
funding.  The author declares no competing interests.

\paragraph{Acknowledgements.} I thank Wouter van Doorn for advice on
exposition, including the explanation of restrictive hypotheses and the
removal of unnecessary terminology. The problem numbering follows the
supplied Erd\H{o}s Problems catalogue snapshot~\cite{erdosproblems}.

\appendix
\section{Guide to the formal sources}\label{long243:app:index}

The main-repository links retain revision \texttt{\commitshort};
other links specify their repository and revision. The supplied source
index records a main-repository build at \texttt{6b78209ab63a} and sixty
theorem-specific Comparator checks in the separate release.
The label \emph{release only} distinguishes the repositories and build
records; it does not deny those sixty checks.
The principal state theorem is in \texttt{ReciprocalTailRigidity.lean}.
The scalar convergence theorem is in \texttt{SparseResetRecovery.lean},
and the bounds on cancelled factors are in \texttt{RepairEntropy.lean}.
The rational-tail estimates are separate checked declarations in
\texttt{PaperCompleteR7}, not consequences of the abstract state module
alone. The checked periodic statement assumes $e_n<a_n$;
Section~\ref{long243:sec:periodic} derives this bound eventually for a
periodic positive error magnitude. The bounded-negative statement lists
strict centring and vanishing relative error separately;
the latter implies the former eventually by taking $K=1$.
Sections~\ref{long243:sec:priorwork} and~\ref{long243:sec:bounded} give the
construction from a rational reciprocal sum.

\subsection*{Formalisation coverage and remaining dependencies}
\label{long243:sec:coverage}
\papersectiontarget{coverage}

Every theorem, lemma, proposition and corollary of this record other than the
four listed below has a Lean statement of the same assertion, with the same
hypotheses, checked by the Lean kernel using only the axioms
\texttt{propext}, \texttt{Classical.choice} and \texttt{Quot.sound}, in the
development at revision \texttt{181078b6b009}.

The exceptions are one chain of results with one shared input.  The
development verifies the reduction of Lemma~\ref{long243:res:squarespec},
Proposition~\ref{long243:res:transportsquare},
Theorem~\ref{long243:res:cubicexclusion} and
Theorem~\ref{long243:res:cubicrate} to the following input: for a finite
Galois extension of $\Q$, a prescribed element $\sigma$ of its Galois group
and a nonzero algebraic integer of the extension, there are arbitrarily large
rational primes at which $\sigma$ is a Frobenius element and at which that
integer has nonzero reduction.  The first clause is the Chebotarev density
theorem in the form cited in the proof of
Lemma~\ref{long243:res:squarespec}~\cite[\S3, p.~15]{stevenhagenlenstra1996};
the second is the elementary fact that a nonzero algebraic integer is
divisible by only finitely many rational primes.  A formal proof of that
input is not included in the verified development; these four endpoints
therefore remain conditional on it.  Everything the chain adds to it is
checked: the square forced modulo a prime dividing a middle numerator, the
passage from that congruence to a square in the cubic field, the exclusion of
the proposed cubic orbit, and the irrationality conclusion.

A theorem whose own statement is conditional is formalised exactly as stated,
with its hypothesis intact.  The four results above are asserted
unconditionally in this record, and their Lean counterparts carry the
additional named input described here.

% The exhaustive source manifest is retained in the authored TeX for automated
% validation, and kept outside the printed note.
\iffalse

\begin{tabular}{@{}ll@{}}
informal statement & linked Lean source\\
Sylvester successor & \lrefx{Erdos243/ReciprocalTailRigidity.lean}{37}{sylvesterNext}\\
denominator update & \lrefx{Erdos243/ReciprocalTailRigidity.lean}{41}{nextDenState}\\
tail update & \lrefx{Erdos243/ReciprocalTailRigidity.lean}{45}{nextTailState}\\
error & \lrefx{Erdos243/ReciprocalTailRigidity.lean}{49}{centeredState}\\
Sylvester defect & \lrefx{Erdos243/ReciprocalTailRigidity.lean}{53}{sylvesterDefect}\\
update law & \lrefx{Erdos243/ReciprocalTailRigidity.lean}{57}{nextTailState_eq_sub_centered}\\
denominator equivariance & \lrefx{Erdos243/ReciprocalTailRigidity.lean}{65}{nextDenState_scale}\\
tail equivariance & \lrefx{Erdos243/ReciprocalTailRigidity.lean}{70}{nextTailState_scale}\\
centred equivariance & \lrefx{Erdos243/ReciprocalTailRigidity.lean}{75}{centeredState_scale}\\
normalised exclusion & \lrefx{Erdos243/ReciprocalTailRigidity.lean}{100}{no_normalizedConstantNegative_orbit}\\
scale-primitive exclusion & \lrefx{Erdos243/ReciprocalTailRigidity.lean}{147}{no_scalePrimitiveConstantNegative_orbit}\\
constant exclusion at every scale & \lrefx{Erdos243/ReciprocalTailRigidity.lean}{286}{no_constantNegative_orbit}\\
eventual constant exclusion & \lrefx{Erdos243/ReciprocalTailRigidity.lean}{350}{no_eventuallyConstantNegative_orbit}\\
persistence of a multiplier & \lrefx{Erdos243/ReciprocalTailRigidity.lean}{382}{base_dvd_denState_later}\\
one-step lock & \lrefx{Erdos243/ReciprocalTailRigidity.lean}{402}{divisorLocksTailSucc}\\
persistence of a lock & \lrefx{Erdos243/ReciprocalTailRigidity.lean}{429}{divisorLock_persists}\\
iterated period & \lrefx{Erdos243/ReciprocalTailRigidity.lean}{461}{periodic_add_mul}\\
iterated drift & \lrefx{Erdos243/ReciprocalTailRigidity.lean}{476}{phaseDrift_add_mul}\\
repeated-divisor transport & \lrefx{Erdos243/ReciprocalTailRigidity.lean}{494}{repeatedDivisor_forces_periodicCommonScale}\\
phase-primitive exclusion & \lrefx{Erdos243/ReciprocalTailRigidity.lean}{548}{no_phasePrimitivePeriodicNegative_orbit}\\
periodic exclusion at every scale & \lrefx{Erdos243/ReciprocalTailRigidity.lean}{690}{no_periodicNegative_orbit}\\
eventual periodic exclusion & \lrefx{Erdos243/ReciprocalTailRigidity.lean}{802}{no_eventuallyPeriodicNegative_orbit}\\
shifted consecutive multiples & \lrefx{Erdos243/ReciprocalTailRigidity.lean}{839}{exists_shifted_consecutiveMultiples}\\
obstruction to bounded upward increments & \lrefx{Erdos243/ReciprocalTailRigidity.lean}{903}{no_boundedRise_of_tailAvoidance}\\
step coprimality & \lrefx{Erdos243/ReciprocalTailRigidity.lean}{989}{reducedStep_coprime_currentFactor}\\
pairwise coprimality & \lrefx{Erdos243/ReciprocalTailRigidity.lean}{1016}{reducedTail_pairwiseCoprime}\\
avoidance of multiples of whole moduli & \lrefx{Erdos243/ReciprocalTailRigidity.lean}{1030}{reducedTail_wholeModulusAvoidance}\\
reduced exclusion under bounded upward increments & \lrefx{Erdos243/ReciprocalTailRigidity.lean}{1041}{no_boundedRise_reducedTail}\\
eventual reduced exclusion & \lrefx{Erdos243/ReciprocalTailRigidity.lean}{1066}{no_eventuallyBoundedRise_reducedTail}\\
gcd monotonicity & \lrefx{Erdos243/ReciprocalTailRigidity.lean}{1430}{tailGcd_dvd_succ}\\
gcd at a negative index & \lrefx{Erdos243/ReciprocalTailRigidity.lean}{1587}{tailGcd_eq_gcd_negativeMagnitude}\\
chain stabilisation & \lrefx{Erdos243/ReciprocalTailRigidity.lean}{1605}{dvdChain_eventuallyConstant_of_cofinally_bounded}\\
gcd stabilisation & \lrefx{Erdos243/ReciprocalTailRigidity.lean}{1636}{tailGcd_eventuallyConstant_of_cofinally_boundedNegative}\\
tail-divergence form & \lrefx{Erdos243/ReciprocalTailRigidity.lean}{1661}{no_cofinallyBoundedNegative_of_tailTendsto}\\
divergence from vanishing relative error & \lrefx{Erdos243/ReciprocalTailRigidity.lean}{1739}{tailState_tendsto_atTop_of_nonzero_normalizedVanishes}\\
exclusion from vanishing relative error & \lrefx{Erdos243/ReciprocalTailRigidity.lean}{1754}{no_cofinallyBoundedNegative_of_normalizedVanishes}\\
defect identity & \lrefx{Erdos243/ReciprocalTailRigidity.lean}{1775}{sylvesterDefect_mul_nextTailState}\\
local rigidity step & \lrefx{Erdos243/ReciprocalTailRigidity.lean}{1787}{sylvesterNext_eq_of_centered_zero}\\
eventual Sylvester recurrence & \lrefx{Erdos243/ReciprocalTailRigidity.lean}{1805}{sylvesterNext_eventually_of_centered_zero}\\
eventual constancy & \lrefx{Erdos243/ReciprocalTailRigidity.lean}{1825}{antitone_nat_eventually_constant}\\
stabilisation of the nonnegative error & \lrefx{Erdos243/ReciprocalTailRigidity.lean}{1843}{centeredState_eventually_zero}\\
tail realisation & \lrefx{Erdos243/ReciprocalTailRigidity.lean}{1869}{natTail_eq_nextTailState}\\
denominator realisation & \lrefx{Erdos243/ReciprocalTailRigidity.lean}{1881}{natDen_eq_nextDenState}\\
signed update law & \lrefx{Erdos243/ReciprocalTailRigidity.lean}{1978}{natTail_eq_sub_centeredState}\\
absorption of a vanishing error & \lrefx{Erdos243/ReciprocalTailRigidity.lean}{2227}{centeredState_zero_absorbing}\\
contrapositive along a tail & \lrefx{Erdos243/ReciprocalTailRigidity.lean}{2254}{eventually_nonzero_of_zero_absorbing}\\
vanishing under a bound on the negative part & \lrefx{Erdos243/ReciprocalTailRigidity.lean}{2271}{boundedNegativePart_eventually_zero}\\
eventual vanishing under a bound on the negative part & \lrefx{Erdos243/ReciprocalTailRigidity.lean}{2360}{eventuallyBoundedNegativePart_eventually_zero}\\
bound on cancellation over a recovery interval & \lrefx{Erdos243/RepairEntropy.lean}{48}{repairedFamily_recovery_energy_divisionFree}\\
independent-repair bound & \lrefx{Erdos243/RepairEntropy.lean}{78}{repaired_card_bound_of_independent}\\
fixed-length reset disappearance & \lrefx{Erdos243/RepairEntropy.lean}{107}{eventually_recoveryPayment_eq_one_of_fixedLength}\\
forced numerator & \lrefx{Erdos243/FiniteHorizonResidue.lean}{22}{forcedNumerator}\\
polynomial congruence & \lrefx{Erdos243/FiniteHorizonResidue.lean}{26}{forcedNumerator_modEq}\\
exact-division cancellation & \lrefx{Erdos243/FiniteHorizonResidue.lean}{41}{quotient_modEq_of_modEq_mul}\\
horizon modulus & \lrefx{Erdos243/FiniteHorizonResidue.lean}{53}{horizonModulus}\\
ascending-factorial form & \lrefx{Erdos243/FiniteHorizonResidue.lean}{59}{horizonModulus_eq_ascFactorial}\\
factorial value at the initial index & \lrefx{Erdos243/FiniteHorizonResidue.lean}{69}{horizonModulus_zero_eq_factorial}\\
survival predicate & \lrefx{Erdos243/FiniteHorizonResidue.lean}{75}{ForcedSurvives}\\
quotient transport & \lrefx{Erdos243/FiniteHorizonResidue.lean}{85}{forcedQuotient_modEq}\\
shrinking-modulus transport & \lrefx{Erdos243/FiniteHorizonResidue.lean}{97}{forcedSurvives_iff_of_modEq}\\
factorial residue reduction & \lrefx{Erdos243/FiniteHorizonResidue.lean}{134}{forcedSurvives_iff_of_modEq_factorial}\\
\end{tabular}
\fi

% BEGIN GENERATED CONCORDANCE (claude_lane/build_concordance.py; do not edit by hand)
\paragraph{Concordance of statements and Lean declarations.}
Each result of this record that has a kernel-checked Lean statement of the same assertion is listed
below with the declarations that jointly state it.  Each name links to its declaration at revision
\texttt{181078b6b009}.  Where the Lean statement is stronger than the printed one and implies it by
an immediate specialisation, the entry says so.

\begingroup\small\raggedright
\par\noindent\hangindent=1.5em Lemma~\ref{long243:res:periodicobstruction}: \href{https://github.com/wcook04/plectis-erdos/blob/181078b6b009d905809cf2e007ad309386a3d1e0/lean/ErdosProblems/Erdos243/PaperCompleteR11/DensityTransport.lean\#L156}{\nolinkurl{fixed_offsets_periodic_lowerDensity}}.
\par\noindent\hangindent=1.5em Lemma~\ref{long243:res:gcdshape}: \href{https://github.com/wcook04/plectis-erdos/blob/181078b6b009d905809cf2e007ad309386a3d1e0/lean/ErdosProblems/Erdos243/PaperCompleteR21/CubicProfileGcdShape.lean\#L103}{\nolinkurl{cubic_profile_gcd_stabilisation_and_primitive_shape}}.
\par\noindent\hangindent=1.5em Lemma~\ref{long243:res:reduciblecase}: \href{https://github.com/wcook04/plectis-erdos/blob/181078b6b009d905809cf2e007ad309386a3d1e0/lean/ErdosProblems/Erdos243/PaperCompleteR11/PrimitiveMultiplierSupply.lean\#L244}{\nolinkurl{primitive_zero_density_paper_multiplier_lemma}}.
\par\noindent\hangindent=1.5em Lemma~\ref{long243:res:scaletwelve}: \href{https://github.com/wcook04/plectis-erdos/blob/181078b6b009d905809cf2e007ad309386a3d1e0/lean/ErdosProblems/Erdos243/PaperCompleteR20/ScaleTwelve.lean\#L36}{\nolinkurl{scale_twelve_of_square_in_rootField}}.
\par\noindent\hangindent=1.5em Lemma~\ref{long243:res:modseven}: \href{https://github.com/wcook04/plectis-erdos/blob/181078b6b009d905809cf2e007ad309386a3d1e0/lean/ErdosProblems/Erdos243/PaperCompleteR20/TwoForbiddenWords.lean\#L18}{\nolinkurl{plus_one_forbidden_word}}, \href{https://github.com/wcook04/plectis-erdos/blob/181078b6b009d905809cf2e007ad309386a3d1e0/lean/ErdosProblems/Erdos243/PaperCompleteR20/TwoForbiddenWords.lean\#L38}{\nolinkurl{minus_one_forbidden_word}}.
\par\noindent\hangindent=1.5em Lemma~\ref{long243:res:extraction}: \href{https://github.com/wcook04/plectis-erdos/blob/181078b6b009d905809cf2e007ad309386a3d1e0/lean/ErdosProblems/Erdos243/PaperCompleteR21/RegularRateExtraction.lean\#L874}{\nolinkurl{regular_rate_extraction}}, \href{https://github.com/wcook04/plectis-erdos/blob/181078b6b009d905809cf2e007ad309386a3d1e0/lean/ErdosProblems/Erdos243/PaperCompleteR21/RegularRateExtraction.lean\#L1151}{\nolinkurl{regular_rate_extraction_cubic}}.
\par\noindent\hangindent=1.5em Lemma~\ref{long243:res:tailratio} (the Lean statement is stronger): \href{https://github.com/wcook04/plectis-erdos/blob/181078b6b009d905809cf2e007ad309386a3d1e0/lean/ErdosProblems/Erdos243/PaperCompleteR7/QuantitativeTail.lean\#L204}{\nolinkurl{canonical_tail_ratio_quantitative}}.
\par\noindent\hangindent=1.5em Proposition~\ref{long243:res:update}: \href{https://github.com/wcook04/plectis-erdos/blob/181078b6b009d905809cf2e007ad309386a3d1e0/lean/ErdosProblems/Erdos243/ReciprocalTailRigidity.lean\#L57}{\nolinkurl{nextTailState_eq_sub_centered}}.
\par\noindent\hangindent=1.5em Proposition~\ref{long243:res:scale}: \href{https://github.com/wcook04/plectis-erdos/blob/181078b6b009d905809cf2e007ad309386a3d1e0/lean/ErdosProblems/Erdos243/PaperCompleteR7/Arithmetic.lean\#L31}{\nolinkurl{state_scale}}.
\par\noindent\hangindent=1.5em Theorem~\ref{long243:res:defect}: \href{https://github.com/wcook04/plectis-erdos/blob/181078b6b009d905809cf2e007ad309386a3d1e0/lean/ErdosProblems/Erdos243/ReciprocalTailRigidity.lean\#L1775}{\nolinkurl{sylvesterDefect_mul_nextTailState}}.
\par\noindent\hangindent=1.5em Theorem~\ref{long243:res:step}: \href{https://github.com/wcook04/plectis-erdos/blob/181078b6b009d905809cf2e007ad309386a3d1e0/lean/ErdosProblems/Erdos243/ReciprocalTailRigidity.lean\#L1787}{\nolinkurl{sylvesterNext_eq_of_centered_zero}}.
\par\noindent\hangindent=1.5em Theorem~\ref{long243:res:secondorder}: \href{https://github.com/wcook04/plectis-erdos/blob/181078b6b009d905809cf2e007ad309386a3d1e0/lean/ErdosProblems/Erdos243/PaperCompleteR7/Arithmetic.lean\#L42}{\nolinkurl{reduced_second_order_int}}.
\par\noindent\hangindent=1.5em Theorem~\ref{long243:res:curvature}: \href{https://github.com/wcook04/plectis-erdos/blob/181078b6b009d905809cf2e007ad309386a3d1e0/lean/ErdosProblems/Erdos243/DynamicCancellation.lean\#L309}{\nolinkurl{cancellationFree_curvature_square}}.
\par\noindent\hangindent=1.5em Lemma~\ref{long243:res:oldmodulussaturation}: \href{https://github.com/wcook04/plectis-erdos/blob/181078b6b009d905809cf2e007ad309386a3d1e0/lean/ErdosProblems/Erdos243/PaperCompleteR21/ReducedStepLocalArithmetic.lean\#L39}{\nolinkurl{unit_word_saturates_old_modulus}}.
\par\noindent\hangindent=1.5em Theorem~\ref{long243:res:eventual}: \href{https://github.com/wcook04/plectis-erdos/blob/181078b6b009d905809cf2e007ad309386a3d1e0/lean/ErdosProblems/Erdos243/ReciprocalTailRigidity.lean\#L1805}{\nolinkurl{sylvesterNext_eventually_of_centered_zero}}.
\par\noindent\hangindent=1.5em Theorem~\ref{long243:res:absorb}: \href{https://github.com/wcook04/plectis-erdos/blob/181078b6b009d905809cf2e007ad309386a3d1e0/lean/ErdosProblems/Erdos243/ReciprocalTailRigidity.lean\#L2227}{\nolinkurl{centeredState_zero_absorbing}}.
\par\noindent\hangindent=1.5em Theorem~\ref{long243:res:arithmeticrecord}: \href{https://github.com/wcook04/plectis-erdos/blob/181078b6b009d905809cf2e007ad309386a3d1e0/lean/ErdosProblems/Erdos243/PaperCompleteR11/ArithmeticWeightedRecord.lean\#L264}{\nolinkurl{arithmetic_weighted_record_dichotomy}}.
\par\noindent\hangindent=1.5em Theorem~\ref{long243:res:weightedrecord}: \href{https://github.com/wcook04/plectis-erdos/blob/181078b6b009d905809cf2e007ad309386a3d1e0/lean/ErdosProblems/Erdos243/PaperCompleteR11/CanonicalRecords.lean\#L202}{\nolinkurl{canonical_weighted_record_excess}}.
\par\noindent\hangindent=1.5em Lemma~\ref{long243:res:valuationtransition}: \href{https://github.com/wcook04/plectis-erdos/blob/181078b6b009d905809cf2e007ad309386a3d1e0/lean/ErdosProblems/Erdos243/PaperCompleteR21/ReducedStepLocalArithmetic.lean\#L244}{\nolinkurl{reduced_denominator_valuation_transition}}, \href{https://github.com/wcook04/plectis-erdos/blob/181078b6b009d905809cf2e007ad309386a3d1e0/lean/ErdosProblems/Erdos243/PaperCompleteR21/ReducedStepLocalArithmetic.lean\#L262}{\nolinkurl{reduced_denominator_valuation_le_max}}, \href{https://github.com/wcook04/plectis-erdos/blob/181078b6b009d905809cf2e007ad309386a3d1e0/lean/ErdosProblems/Erdos243/PaperCompleteR21/ReducedStepLocalArithmetic.lean\#L279}{\nolinkurl{reduced_denominator_valuation_strict_loss}}.
\par\noindent\hangindent=1.5em Corollary~\ref{long243:res:powerpersistence}: \href{https://github.com/wcook04/plectis-erdos/blob/181078b6b009d905809cf2e007ad309386a3d1e0/lean/ErdosProblems/Erdos243/PaperCompleteR20/PowerPersistence.lean\#L27}{\nolinkurl{primePower_persists}}.
\par\noindent\hangindent=1.5em Theorem~\ref{long243:res:recorddichotomy}: \href{https://github.com/wcook04/plectis-erdos/blob/181078b6b009d905809cf2e007ad309386a3d1e0/lean/ErdosProblems/Erdos243/PaperCompleteR11/QuantitativeRecordDichotomy.lean\#L340}{\nolinkurl{canonical_quantitative_record_dichotomy}}, \href{https://github.com/wcook04/plectis-erdos/blob/181078b6b009d905809cf2e007ad309386a3d1e0/lean/ErdosProblems/Erdos243/PaperCompleteR11/QuantitativeRecordDichotomy.lean\#L308}{\nolinkurl{canonical_recordTheta_zero_or_gt_one}}, \href{https://github.com/wcook04/plectis-erdos/blob/181078b6b009d905809cf2e007ad309386a3d1e0/lean/ErdosProblems/Erdos243/PaperCompleteR11/QuantitativeRecordDichotomy.lean\#L321}{\nolinkurl{canonical_recordTheta_eq_zero_iff}}, \href{https://github.com/wcook04/plectis-erdos/blob/181078b6b009d905809cf2e007ad309386a3d1e0/lean/ErdosProblems/Erdos243/PaperCompleteR11/InclusiveLimsup.lean\#L54}{\nolinkurl{canonical_recordTheta_gt_one}}.
\par\noindent\hangindent=1.5em Corollary~\ref{long243:res:loglogboundary}: \href{https://github.com/wcook04/plectis-erdos/blob/181078b6b009d905809cf2e007ad309386a3d1e0/lean/ErdosProblems/Erdos243/PaperCompleteR11/InclusiveLimsup.lean\#L129}{\nolinkurl{canonical_inclusive_logLog_criterion}}, \href{https://github.com/wcook04/plectis-erdos/blob/181078b6b009d905809cf2e007ad309386a3d1e0/lean/ErdosProblems/Erdos243/PaperCompleteR11/InclusiveLimsup.lean\#L108}{\nolinkurl{canonical_negativeError_limsup_gt_one}}.
\par\noindent\hangindent=1.5em Theorem~\ref{long243:res:recordamplified}: \href{https://github.com/wcook04/plectis-erdos/blob/181078b6b009d905809cf2e007ad309386a3d1e0/lean/ErdosProblems/Erdos243/PaperCompleteR21/AmplifiedRecordEquivalence.lean\#L975}{\nolinkurl{recordAmplified}}, \href{https://github.com/wcook04/plectis-erdos/blob/181078b6b009d905809cf2e007ad309386a3d1e0/lean/ErdosProblems/Erdos243/PaperCompleteR21/AmplifiedRecordEquivalence.lean\#L931}{\nolinkurl{amp_bddAbove_iff_sylvester}}, \href{https://github.com/wcook04/plectis-erdos/blob/181078b6b009d905809cf2e007ad309386a3d1e0/lean/ErdosProblems/Erdos243/PaperCompleteR21/AmplifiedRecordEquivalence.lean\#L948}{\nolinkurl{Rdelta_bddAbove_iff_amp}}, \href{https://github.com/wcook04/plectis-erdos/blob/181078b6b009d905809cf2e007ad309386a3d1e0/lean/ErdosProblems/Erdos243/PaperCompleteR21/AmplifiedRecordEquivalence.lean\#L587}{\nolinkurl{eventuallySylvester_of_amp_le}}, \href{https://github.com/wcook04/plectis-erdos/blob/181078b6b009d905809cf2e007ad309386a3d1e0/lean/ErdosProblems/Erdos243/PaperCompleteR21/AmplifiedRecordEquivalence.lean\#L187}{\nolinkurl{coprimeMultiplier_cofinal}}, \href{https://github.com/wcook04/plectis-erdos/blob/181078b6b009d905809cf2e007ad309386a3d1e0/lean/ErdosProblems/Erdos243/PaperCompleteR21/AmplifiedRecordEquivalence.lean\#L295}{\nolinkurl{largePrime_coprimeMultiplier}}, \href{https://github.com/wcook04/plectis-erdos/blob/181078b6b009d905809cf2e007ad309386a3d1e0/lean/ErdosProblems/Erdos243/PaperCompleteR21/AmplifiedRecordEquivalence.lean\#L546}{\nolinkurl{primeBlock_supply}}, \href{https://github.com/wcook04/plectis-erdos/blob/181078b6b009d905809cf2e007ad309386a3d1e0/lean/ErdosProblems/Erdos243/PaperCompleteR21/AmplifiedRecordEquivalence.lean\#L462}{\nolinkurl{canc_lt_of_amp_le}}, \href{https://github.com/wcook04/plectis-erdos/blob/181078b6b009d905809cf2e007ad309386a3d1e0/lean/ErdosProblems/Erdos243/PaperCompleteR21/AmplifiedRecordEquivalence.lean\#L738}{\nolinkurl{delta_negPart_comparison}}, \href{https://github.com/wcook04/plectis-erdos/blob/181078b6b009d905809cf2e007ad309386a3d1e0/lean/ErdosProblems/Erdos243/PaperCompleteR21/AmplifiedRecordEquivalence.lean\#L870}{\nolinkurl{R_delta_sub_amp_tendsto_zero}}.
\par\noindent\hangindent=1.5em Corollary~\ref{long243:res:criticalrate}: \href{https://github.com/wcook04/plectis-erdos/blob/181078b6b009d905809cf2e007ad309386a3d1e0/lean/ErdosProblems/Erdos243/PaperCompleteR21/AmplifiedRecordEquivalence.lean\#L1033}{\nolinkurl{criticalRate}}, \href{https://github.com/wcook04/plectis-erdos/blob/181078b6b009d905809cf2e007ad309386a3d1e0/lean/ErdosProblems/Erdos243/PaperCompleteR21/AmplifiedRecordEquivalence.lean\#L1197}{\nolinkurl{criticalRate_counterexample}}, \href{https://github.com/wcook04/plectis-erdos/blob/181078b6b009d905809cf2e007ad309386a3d1e0/lean/ErdosProblems/Erdos243/PaperCompleteR21/AmplifiedRecordEquivalence.lean\#L1016}{\nolinkurl{R_le_of_u_le}}.
\par\noindent\hangindent=1.5em Lemma~\ref{long243:res:oddpowersupply}: \href{https://github.com/wcook04/plectis-erdos/blob/181078b6b009d905809cf2e007ad309386a3d1e0/lean/ErdosProblems/Erdos243/PaperCompleteR21/ReducedDenominatorPrimePowers.lean\#L790}{\nolinkurl{oddPrimePower_supply}}, \href{https://github.com/wcook04/plectis-erdos/blob/181078b6b009d905809cf2e007ad309386a3d1e0/lean/ErdosProblems/Erdos243/PaperCompleteR21/ReducedDenominatorPrimePowers.lean\#L718}{\nolinkurl{oddPrimePower_supply_nat}}.
\par\noindent\hangindent=1.5em Theorem~\ref{long243:res:unitrecord}: \href{https://github.com/wcook04/plectis-erdos/blob/181078b6b009d905809cf2e007ad309386a3d1e0/lean/ErdosProblems/Erdos243/PaperCompleteR21/ReducedDenominatorPrimePowers.lean\#L846}{\nolinkurl{unitRecordIncrement_sylvesterNext}}, \href{https://github.com/wcook04/plectis-erdos/blob/181078b6b009d905809cf2e007ad309386a3d1e0/lean/ErdosProblems/Erdos243/PaperCompleteR21/ReducedDenominatorPrimePowers.lean\#L950}{\nolinkurl{unitRecordIncrement_criterion}}.
\par\noindent\hangindent=1.5em Theorem~\ref{long243:res:epochenergy}: \href{https://github.com/wcook04/plectis-erdos/blob/181078b6b009d905809cf2e007ad309386a3d1e0/lean/ErdosProblems/Erdos243/PaperCompleteR21/ProtectedEpochBarrierCount.lean\#L149}{\nolinkurl{epoch_energy_named_crossing_set}}, \href{https://github.com/wcook04/plectis-erdos/blob/181078b6b009d905809cf2e007ad309386a3d1e0/lean/ErdosProblems/Erdos243/PaperCompleteR21/ProtectedEpochBarrierCount.lean\#L34}{\nolinkurl{mem_barrierIdx_iff}}.
\par\noindent\hangindent=1.5em Theorem~\ref{long243:res:energycriterion}: \href{https://github.com/wcook04/plectis-erdos/blob/181078b6b009d905809cf2e007ad309386a3d1e0/lean/ErdosProblems/Erdos243/PaperCompleteR21/RecordJumpEnergySeries.lean\#L468}{\nolinkurl{energy_criterion}}, \href{https://github.com/wcook04/plectis-erdos/blob/181078b6b009d905809cf2e007ad309386a3d1e0/lean/ErdosProblems/Erdos243/PaperCompleteR21/RecordJumpEnergySeries.lean\#L414}{\nolinkurl{energy_summable_iff}}, \href{https://github.com/wcook04/plectis-erdos/blob/181078b6b009d905809cf2e007ad309386a3d1e0/lean/ErdosProblems/Erdos243/PaperCompleteR21/RecordJumpEnergySeries.lean\#L455}{\nolinkurl{energySqrt_summable_iff}}, \href{https://github.com/wcook04/plectis-erdos/blob/181078b6b009d905809cf2e007ad309386a3d1e0/lean/ErdosProblems/Erdos243/PaperCompleteR21/RecordJumpEnergySeries.lean\#L161}{\nolinkurl{energy_le_two_energySqrt}}, \href{https://github.com/wcook04/plectis-erdos/blob/181078b6b009d905809cf2e007ad309386a3d1e0/lean/ErdosProblems/Erdos243/PaperCompleteR21/RecordJumpEnergySeries.lean\#L263}{\nolinkurl{exists_late_energy_window}}, \href{https://github.com/wcook04/plectis-erdos/blob/181078b6b009d905809cf2e007ad309386a3d1e0/lean/ErdosProblems/Erdos243/PaperCompleteR21/RecordJumpEnergySeries.lean\#L84}{\nolinkurl{energy_window_real_bound}}.
\par\noindent\hangindent=1.5em Theorem~\ref{long243:res:slownegative}: \href{https://github.com/wcook04/plectis-erdos/blob/181078b6b009d905809cf2e007ad309386a3d1e0/lean/ErdosProblems/Erdos243/PaperCompleteR21/ExactOrbitRecordDichotomy.lean\#L559}{\nolinkurl{slowNegative_eventually_zero_and_sylvesterNext_unconditional}}, \href{https://github.com/wcook04/plectis-erdos/blob/181078b6b009d905809cf2e007ad309386a3d1e0/lean/ErdosProblems/Erdos243/PaperCompleteR21/ExactOrbitRecordDichotomy.lean\#L363}{\nolinkurl{exactOrbit_recordTheta_gt_one}}, \href{https://github.com/wcook04/plectis-erdos/blob/181078b6b009d905809cf2e007ad309386a3d1e0/lean/ErdosProblems/Erdos243/PaperCompleteR21/ExactOrbitRecordDichotomy.lean\#L522}{\nolinkurl{exactOrbit_one_le_recordTheta}}, \href{https://github.com/wcook04/plectis-erdos/blob/181078b6b009d905809cf2e007ad309386a3d1e0/lean/ErdosProblems/Erdos243/PaperCompleteR21/ExactOrbitRecordDichotomy.lean\#L288}{\nolinkurl{exactOrbit_unbounded_of_error_not_eventually_zero}}, \href{https://github.com/wcook04/plectis-erdos/blob/181078b6b009d905809cf2e007ad309386a3d1e0/lean/ErdosProblems/Erdos243/PaperCompleteR21/ExactOrbitRecordDichotomy.lean\#L145}{\nolinkurl{tail_multiplier_quadratic_lower}}, \href{https://github.com/wcook04/plectis-erdos/blob/181078b6b009d905809cf2e007ad309386a3d1e0/lean/ErdosProblems/Erdos243/PaperCompleteR21/ExactOrbitRecordDichotomy.lean\#L254}{\nolinkurl{tail_binaryTower_lower}}, \href{https://github.com/wcook04/plectis-erdos/blob/181078b6b009d905809cf2e007ad309386a3d1e0/lean/ErdosProblems/Erdos243/PaperCompleteR21/ExactOrbitRecordDichotomy.lean\#L218}{\nolinkurl{exists_multiplier_ge_four}}.
\par\noindent\hangindent=1.5em Theorem~\ref{long243:res:strausbounded}: \href{https://github.com/wcook04/plectis-erdos/blob/181078b6b009d905809cf2e007ad309386a3d1e0/lean/ErdosProblems/Erdos243/PaperCompleteR7/ProductDefect.lean\#L211}{\nolinkurl{original_coordinate_bounded_defect}}, \href{https://github.com/wcook04/plectis-erdos/blob/181078b6b009d905809cf2e007ad309386a3d1e0/lean/ErdosProblems/Erdos243/PaperCompleteR21/SlowGrowthProductIncrements.lean\#L155}{\nolinkurl{original_coordinate_slow_growth_defect}}, \href{https://github.com/wcook04/plectis-erdos/blob/181078b6b009d905809cf2e007ad309386a3d1e0/lean/ErdosProblems/Erdos243/PaperCompleteR21/SlowGrowthProductIncrements.lean\#L61}{\nolinkurl{recordTheta_le_of_slow_negative}}, \href{https://github.com/wcook04/plectis-erdos/blob/181078b6b009d905809cf2e007ad309386a3d1e0/lean/ErdosProblems/Erdos243/PaperCompleteR21/SlowGrowthProductIncrements.lean\#L111}{\nolinkurl{prefix_ratio_le_canonicalNumerator}}.
\par\noindent\hangindent=1.5em Corollary~\ref{long243:res:onethreshold} (the Lean statement is stronger): \href{https://github.com/wcook04/plectis-erdos/blob/181078b6b009d905809cf2e007ad309386a3d1e0/lean/ErdosProblems/Erdos243/PaperCompleteR21/ProductDefectThresholds.lean\#L92}{\nolinkurl{original_coordinate_strict_one}}, \href{https://github.com/wcook04/plectis-erdos/blob/181078b6b009d905809cf2e007ad309386a3d1e0/lean/ErdosProblems/Erdos243/PaperCompleteR20/InclusiveOne.lean\#L207}{\nolinkurl{original_coordinate_inclusive_one}}, \href{https://github.com/wcook04/plectis-erdos/blob/181078b6b009d905809cf2e007ad309386a3d1e0/lean/ErdosProblems/Erdos243/PaperCompleteR20/InclusiveOne.lean\#L225}{\nolinkurl{original_coordinate_inclusive_one_pointwise}}.
\par\noindent\hangindent=1.5em Proposition~\ref{long243:res:classicalhalfspace}: \href{https://github.com/wcook04/plectis-erdos/blob/181078b6b009d905809cf2e007ad309386a3d1e0/lean/ErdosProblems/Erdos243/PaperCompleteR21/ClassicalHalfspaceSigns.lean\#L206}{\nolinkurl{overlapDebt_dvd_gcd}}, \href{https://github.com/wcook04/plectis-erdos/blob/181078b6b009d905809cf2e007ad309386a3d1e0/lean/ErdosProblems/Erdos243/PaperCompleteR21/ClassicalHalfspaceSigns.lean\#L225}{\nolinkurl{canonicalError_div_overlapDebt}}, \href{https://github.com/wcook04/plectis-erdos/blob/181078b6b009d905809cf2e007ad309386a3d1e0/lean/ErdosProblems/Erdos243/PaperCompleteR21/ClassicalHalfspaceSigns.lean\#L325}{\nolinkurl{growthDefect_eq_neg_relativeError_add_shiftedCorrection}}, \href{https://github.com/wcook04/plectis-erdos/blob/181078b6b009d905809cf2e007ad309386a3d1e0/lean/ErdosProblems/Erdos243/PaperCompleteR21/ClassicalHalfspaceSigns.lean\#L415}{\nolinkurl{canonical_growthDefect_identity}}, \href{https://github.com/wcook04/plectis-erdos/blob/181078b6b009d905809cf2e007ad309386a3d1e0/lean/ErdosProblems/Erdos243/PaperCompleteR21/ClassicalHalfspaceSigns.lean\#L459}{\nolinkurl{canonicalCorrection_pos_and_lt_three_div}}, \href{https://github.com/wcook04/plectis-erdos/blob/181078b6b009d905809cf2e007ad309386a3d1e0/lean/ErdosProblems/Erdos243/PaperCompleteR21/ClassicalHalfspaceSigns.lean\#L380}{\nolinkurl{shiftedCorrectionTerm_pos_and_lt_three_div}}, \href{https://github.com/wcook04/plectis-erdos/blob/181078b6b009d905809cf2e007ad309386a3d1e0/lean/ErdosProblems/Erdos243/PaperCompleteR21/ClassicalHalfspaceSigns.lean\#L527}{\nolinkurl{erdosStraus_lcm_includes_digit_not_denominator}}, \href{https://github.com/wcook04/plectis-erdos/blob/181078b6b009d905809cf2e007ad309386a3d1e0/lean/ErdosProblems/Erdos243/PaperCompleteR21/ClassicalHalfspaceSigns.lean\#L543}{\nolinkurl{erdosStrausQuantity_ne_productDefect_ne_lcmShift}}, \href{https://github.com/wcook04/plectis-erdos/blob/181078b6b009d905809cf2e007ad309386a3d1e0/lean/ErdosProblems/Erdos243/PaperCompleteR21/ClassicalHalfspaceSigns.lean\#L573}{\nolinkurl{erdosStrausQuantity_sign}}, \href{https://github.com/wcook04/plectis-erdos/blob/181078b6b009d905809cf2e007ad309386a3d1e0/lean/ErdosProblems/Erdos243/PaperCompleteR21/ClassicalHalfspaceSigns.lean\#L351}{\nolinkurl{sylvesterTail_shiftedCorrection}}, \href{https://github.com/wcook04/plectis-erdos/blob/181078b6b009d905809cf2e007ad309386a3d1e0/lean/ErdosProblems/Erdos243/PaperCompleteR21/ClassicalHalfspaceSigns.lean\#L135}{\nolinkurl{overlapDebt_mul_lcmClearedNumerator}}, \href{https://github.com/wcook04/plectis-erdos/blob/181078b6b009d905809cf2e007ad309386a3d1e0/lean/ErdosProblems/Erdos243/PaperCompleteR21/ClassicalHalfspaceSigns.lean\#L82}{\nolinkurl{digitProductScale_eq_canonicalDenominator}}.
\par\noindent\hangindent=1.5em Proposition~\ref{long243:res:coprimalitycap}: \href{https://github.com/wcook04/plectis-erdos/blob/181078b6b009d905809cf2e007ad309386a3d1e0/lean/ErdosProblems/Erdos243/PaperCompleteR21/WindowAvoidance.lean\#L201}{\nolinkurl{exists_avoiding_in_window}}, \href{https://github.com/wcook04/plectis-erdos/blob/181078b6b009d905809cf2e007ad309386a3d1e0/lean/ErdosProblems/Erdos243/PaperCompleteR21/WindowAvoidance.lean\#L423}{\nolinkurl{exists_slow_rise_avoiding_sequence}}.
\par\noindent\hangindent=1.5em Theorem~\ref{long243:res:descent}: \href{https://github.com/wcook04/plectis-erdos/blob/181078b6b009d905809cf2e007ad309386a3d1e0/lean/ErdosProblems/Erdos243/ReciprocalTailRigidity.lean\#L1843}{\nolinkurl{centeredState_eventually_zero}}.
\par\noindent\hangindent=1.5em Theorem~\ref{long243:res:constant}: \href{https://github.com/wcook04/plectis-erdos/blob/181078b6b009d905809cf2e007ad309386a3d1e0/lean/ErdosProblems/Erdos243/PaperCompleteR21/NegativeMagnitudeExclusions.lean\#L38}{\nolinkurl{no_constantNegative_shapeEquation}}, \href{https://github.com/wcook04/plectis-erdos/blob/181078b6b009d905809cf2e007ad309386a3d1e0/lean/ErdosProblems/Erdos243/PaperCompleteR21/NegativeMagnitudeExclusions.lean\#L52}{\nolinkurl{no_eventuallyConstantNegative_shapeEquation}}.
\par\noindent\hangindent=1.5em Theorem~\ref{long243:res:periodic}: \href{https://github.com/wcook04/plectis-erdos/blob/181078b6b009d905809cf2e007ad309386a3d1e0/lean/ErdosProblems/Erdos243/PaperCompleteR21/NegativeMagnitudeExclusions.lean\#L68}{\nolinkurl{no_periodicNegative_shapeEquation}}.
\par\noindent\hangindent=1.5em Lemma~\ref{long243:res:crt}: \href{https://github.com/wcook04/plectis-erdos/blob/181078b6b009d905809cf2e007ad309386a3d1e0/lean/ErdosProblems/Erdos243/PaperCompleteR21/ForbiddenBlockCrossing.lean\#L26}{\nolinkurl{exists_shiftedBlock_consecutiveMultiples}}.
\par\noindent\hangindent=1.5em Theorem~\ref{long243:res:barrier}: \href{https://github.com/wcook04/plectis-erdos/blob/181078b6b009d905809cf2e007ad309386a3d1e0/lean/ErdosProblems/Erdos243/PaperCompleteR21/ForbiddenBlockCrossing.lean\#L44}{\nolinkurl{no_boundedRise_coprimeToEarlierModuli}}.
\par\noindent\hangindent=1.5em Proposition~\ref{long243:res:reduced}: \href{https://github.com/wcook04/plectis-erdos/blob/181078b6b009d905809cf2e007ad309386a3d1e0/lean/ErdosProblems/Erdos243/PaperCompleteR7/Reduction.lean\#L16}{\nolinkurl{persistent_coprimality}}.
\par\noindent\hangindent=1.5em Proposition~\ref{long243:res:gcdstab}: \href{https://github.com/wcook04/plectis-erdos/blob/181078b6b009d905809cf2e007ad309386a3d1e0/lean/ErdosProblems/Erdos243/PaperCompleteR7/Reduction.lean\#L103}{\nolinkurl{gcd_stabilises_and_reduces}}.
\par\noindent\hangindent=1.5em Proposition~\ref{long243:res:gcdsparse}: \href{https://github.com/wcook04/plectis-erdos/blob/181078b6b009d905809cf2e007ad309386a3d1e0/lean/ErdosProblems/Erdos243/PaperCompleteR7/Limits.lean\#L97}{\nolinkurl{sparse_gcd_changes}}.
\par\noindent\hangindent=1.5em Theorem~\ref{long243:res:bounded} (the Lean statement is stronger): \href{https://github.com/wcook04/plectis-erdos/blob/181078b6b009d905809cf2e007ad309386a3d1e0/lean/ErdosProblems/Erdos243/ReciprocalTailRigidity.lean\#L2360}{\nolinkurl{eventuallyBoundedNegativePart_eventually_zero}}.
\par\noindent\hangindent=1.5em Corollary~\ref{long243:res:cor} (the Lean statement is stronger): \href{https://github.com/wcook04/plectis-erdos/blob/181078b6b009d905809cf2e007ad309386a3d1e0/lean/ErdosProblems/Erdos243/ReciprocalTailRigidity.lean\#L2412}{\nolinkurl{boundedNegativePart_sylvesterNext_eventually}}.
\par\noindent\hangindent=1.5em Theorem~\ref{long243:res:mass}: \href{https://github.com/wcook04/plectis-erdos/blob/181078b6b009d905809cf2e007ad309386a3d1e0/lean/ErdosProblems/Erdos243/PaperCompleteR7/Frontier.lean\#L84}{\nolinkurl{finite_negative_mass_paper}}.
\par\noindent\hangindent=1.5em Proposition~\ref{long243:res:frontier}: \href{https://github.com/wcook04/plectis-erdos/blob/181078b6b009d905809cf2e007ad309386a3d1e0/lean/ErdosProblems/Erdos243/PaperCompleteR7/Frontier.lean\#L151}{\nolinkurl{canonical_frontier}}.
\par\noindent\hangindent=1.5em Proposition~\ref{long243:res:variablerise} (the Lean statement is stronger): \href{https://github.com/wcook04/plectis-erdos/blob/181078b6b009d905809cf2e007ad309386a3d1e0/lean/ErdosProblems/Erdos243/PaperCompleteR21/WindowAvoidance.lean\#L862}{\nolinkurl{exists_sparse_prime_coprime_sequence}}.
\par\noindent\hangindent=1.5em Theorem~\ref{long243:res:gapconstant}: \href{https://github.com/wcook04/plectis-erdos/blob/181078b6b009d905809cf2e007ad309386a3d1e0/lean/ErdosProblems/Erdos243/PaperCompleteR21/MaximalGapConstant.lean\#L1041}{\nolinkurl{maximal_gap_limsup_eq_inv_sigma}}.
\par\noindent\hangindent=1.5em Theorem~\ref{long243:res:residue}: \href{https://github.com/wcook04/plectis-erdos/blob/181078b6b009d905809cf2e007ad309386a3d1e0/lean/ErdosProblems/Erdos243/PaperCompleteR21/ForcedOrbitResidueHorizon.lean\#L21}{\nolinkurl{forcedOrbit_survives_iff_of_factorial_modEq}}, \href{https://github.com/wcook04/plectis-erdos/blob/181078b6b009d905809cf2e007ad309386a3d1e0/lean/ErdosProblems/Erdos243/FiniteHorizonResidue.lean\#L134}{\nolinkurl{forcedSurvives_iff_of_modEq_factorial}}.
\par\endgroup
% END GENERATED CONCORDANCE

\section{A factorial residue reduction for forced orbits}\label{long243:app:residue}

In the case $m=c=1$ of Section~\ref{long243:sec:constant}, where the shape
equation~\eqref{long243:eq:shape} reads $D_n+1=(a_n-1)(n+1)$, each multiplier is
determined by its predecessor; we call such an orbit \emph{forced}.  At index
$n$ the numerator of the next multiplier is
\[
 \fnum(n,a)=(n+1)a^{2}-(n+2)a+(n+3),
\]
the
\lword{Erdos243/FiniteHorizonResidue.lean}{22}{forcedNumerator}{forced
numerator}, and the divisor is $n+2$.  The orbit continues for another step exactly when this division is
exact. The
\lword{Erdos243/FiniteHorizonResidue.lean}{75}{ForcedSurvives}{formal survival predicate}
requires exact division at each step.  Example~\ref{long243:ex:shape} is the forced orbit from $a=3$ read this
way: $\fnum(0,3)=6$ is divisible by $2$ and gives $a_1=3$, while
$\fnum(1,3)=13$ is not divisible by $3$, so the orbit stops there.  Direct iteration can produce very large intermediate integers. To
decide whether the first $h$ divisions are exact, however, it suffices
to know the initial value modulo $(h+1)!$.
Computing a pseudo-greedy orbit through residues modulo a shrinking product
modulus is Koizumi's method \cite[Remark~2 and Algorithm~1,
pp.~13--14]{koizumi2025}; for the forced numerator that product is a
factorial.

\begin{theorem}[factorial residue reduction]\label{long243:res:residue}
For all $h$ and all integers $a\equiv b \pmod{(h+1)!}$, the orbit from $a$
survives $h$ forced updates exactly when the orbit from $b$ does.
\end{theorem}
\leannote{Lean: \lproof{ErdosProblems/Erdos243/PaperCompleteR21/ForcedOrbitResidueHorizon.lean}{21}{forcedOrbit_survives_iff_of_factorial_modEq}, \lproof{ErdosProblems/Erdos243/FiniteHorizonResidue.lean}{134}{forcedSurvives_iff_of_modEq_factorial}.}

\begin{proof}
Define $M(0,i)=1$ and $M(h+1,i)=(i+2)M(h,i+1)$. Thus
$M(h,0)=(h+1)!$. We prove the stronger statement that, at index $i$,
congruent inputs modulo $M(h,i)$ survive the same $h$ updates.
There is nothing to prove for $h=0$.

For the inductive step, suppose
$a\equiv b\pmod{(i+2)M(h,i+1)}$. Since $\fnum(i,\cdot)$ is an integer
polynomial, its values at $a$ and $b$ are congruent modulo that product.
In particular, one is divisible by $i+2$ if and only if the other is.
If neither is divisible, both orbits stop. Otherwise their quotients
are congruent modulo $M(h,i+1)$, so the induction hypothesis applies
to the remaining $h$ updates at index $i+1$. Taking $i=0$ proves the claim.
\end{proof}

\noindent The formal
\lword{Erdos243/FiniteHorizonResidue.lean}{134}{forcedSurvives_iff_of_modEq_factorial}{factorial
residue reduction} uses the
\lword{Erdos243/FiniteHorizonResidue.lean}{97}{forcedSurvives_iff_of_modEq}{shrinking-modulus
induction},
\lword{Erdos243/FiniteHorizonResidue.lean}{26}{forcedNumerator_modEq}{polynomial
congruence}, and
\lword{Erdos243/FiniteHorizonResidue.lean}{41}{quotient_modEq_of_modEq_mul}{cancellation
after exact division}. Its modulus is identified by the
\lword{Erdos243/FiniteHorizonResidue.lean}{59}{horizonModulus_eq_ascFactorial}{ascending-factorial
formula} and its
\lword{Erdos243/FiniteHorizonResidue.lean}{69}{horizonModulus_zero_eq_factorial}{factorial
value at the initial index}.

At $h=1$ the modulus is $2!=2$, and surviving one update means
$2\mid a^{2}-2a+3$, which holds exactly for odd $a$: for instance
$\fnum(0,3)=6$ but $\fnum(0,4)=11$.  So one step of survival is decided by the
parity of $a$ alone, which is Theorem~\ref{long243:res:residue} at its smallest
nontrivial horizon.

Exact enumeration for $2\le a_0<5000$ gives a maximum of $17$
successful updates, or $18$ multiplier values including $a_0$.
Nine seeds attain it, the least being $1501$. The seed $a_0=1$ is excluded:
it gives $a_n=1$ at every step and violates the hypothesis $a_n\ge2$.
Survival here counts exact divisions, not all the conditions on a positive
integer orbit. The finite search does not prove the infinite exclusion;
Theorem~\ref{long243:res:constant} does so under its stated hypotheses.
The reduction remains useful because its shrinking-modulus argument
applies to other recurrences defined by polynomial division.
% ; - part family_catalogue ; -
\section{Result map and proof dependencies}\label{sec:erdos-243-complete-family-map}
The following map records proof dependencies. A checked component does
not confer that status on an assembled ordinary theorem.
The short note contains its bounded-increment proof in Sections~2--4.
Details kept here rather than repeated there are the signed-series
comparison in Section~\ref{long243:sec:priorwork}, the integer-coefficient
extension at the end of Section~\ref{long243:sec:lcmrecords}, the general
$1+c/n$ comparison in Section~\ref{long243:sec:records}, and the scalar
failed-divisibility example in Section~\ref{long243:sec:open}.
The cubic and double-logarithmic arguments retain their separate proofs
in Section~\ref{long243:sec:cubicrate} and
Theorem~\ref{long243:res:recorddichotomy}, respectively.

\paragraph{Bounded upward increments.}
Vanishing relative error gives eventual absorption. Bounded negative
values at arbitrarily late indices give a stable gcd; CRT then excludes
the unbounded reduced tail with bounded upward increments. The state
theorem and rational-tail construction have supplied main-repository
build evidence. The increment bound is an extra hypothesis.

\paragraph{A convergent sum of relative increases.}
The inequality $C_{n+1}\le C_n(1+(-E_n)_+/C_n)$ turns summability into a
uniform bound on $C_n$. Integrality then permits only finitely many rises,
after which the numerator stabilises.
No vanishing relative error is needed. The scalar and exact-orbit results
have cited main-repository declarations; the linked rational-tail
formulation is outside that build record.

\paragraph{New maxima and prime-power persistence.}
A bound on the negative reduced error scaled by the previous maximum
also bounds every cancellation factor. Primes larger than that bound
persist in the reduced denominator, giving the CRT contradiction in
Theorem~\ref{long243:res:recordamplified}. The double-logarithmic argument
uses the density-one coprimality count to obtain enough moduli, then
controls the height of the CRT block.
The LCM numerator gives weighted first-crossing estimates. A prime power
dividing the reduced denominator persists while the unreduced next
numerator stays below the threshold in
Corollary~\ref{long243:res:powerpersistence}. Crossings can be counted
before that threshold is reached. The large odd prime powers and
convergence equivalences have written proofs. Separate-release integer
lemmas and main-repository conditional record declarations establish
only their stated conclusions.

\paragraph{Cubic rate.}
The tail estimate and finite differences of the Gamma ratio force an
eventual cubic polynomial for the integer numerator. The recurrence and
Chebotar\"ev then force a square in the cubic field. The trace calculation
leaves $m=12$, and both possible signs are excluded modulo seven.
This is an ordinary proof. Finite checks and component sources do not
constitute a formal proof of the complete cubic-rate theorem.

\paragraph{Static barriers.}
For sufficiently sparse prime moduli, a coprime integer sequence can
have unbounded but very small increments. At the specified
double-exponential scale, counting residue classes and applying CRT give
the largest gaps between integers divisible by none of the whole moduli.
These auxiliary results do not construct reciprocal-tail examples.

\paragraph{What remains.}
For the unrestricted problem, one still needs the stated summability or
record bound, or a contradiction to the necessary conditions on a
counterexample. The static examples and local identities supply neither.

% ; - part back ; -
\begin{thebibliography}{99}
\small
\raggedright
\setlength{\itemsep}{1pt}
\bibitem{erdosstraus1964} P.~Erd\H{o}s and E.~G. Straus,
  \href{https://users.renyi.hu/~p_erdos/1964-19.pdf}{\emph{On the irrationality
  of certain Ahmes series}},
  J. Indian Math. Soc. (N.S.) \textbf{27} (1964), 129--133. MR~175848.
\bibitem{duverney2001} D.~Duverney,
  \href{https://www.ms.u-tokyo.ac.jp/journal/pdf/jms080206.pdf}
  {\emph{Irrationality of fast converging series of rational numbers}},
  J. Math. Sci. Univ. Tokyo \textbf{8} (2001), 275--316. MR~1837165.
\bibitem{badea1993} C.~Badea,
  \href{https://matwbn.icm.edu.pl/ksiazki/aa/aa63/aa6342.pdf}{\emph{A theorem
  on irrationality of infinite series and applications}}, Acta Arith.
  \textbf{63} (1993), no.~4, 313--323,
  doi:\href{https://doi.org/10.4064/aa-63-4-313-323}{10.4064/aa-63-4-313-323}.
\bibitem{tijdemanyuan2002} R.~Tijdeman and P.~Yuan,
  \emph{On the rationality of Cantor and Ahmes series},
  Indag. Math. (N.S.) \textbf{13} (2002), no.~3, 407--418,
  doi:\href{https://doi.org/10.1016/S0019-3577(02)80018-0}{10.1016/S0019-3577(02)80018-0}.
\bibitem{erdosgraham1980} P.~Erd\H{o}s and R.~L. Graham,
  \href{https://mathweb.ucsd.edu/~ronspubs/80_11_number_theory.pdf}{\emph{Old
  and New Problems and Results in Combinatorial Number Theory}},
  Monogr. Enseign. Math. 28, Geneva, 1980, p.~64.
\bibitem{erdos1988} P.~Erd\H{o}s,
  \href{https://users.renyi.hu/~p_erdos/1988-22.pdf}{\emph{On the irrationality
  of certain series: problems and results}}, in A.~Baker (ed.), \emph{New
  Advances in Transcendence Theory}, Cambridge UP, 1988, pp.~102--109,
  doi:\href{https://doi.org/10.1017/CBO9780511897184.009}{10.1017/CBO9780511897184.009}.
\bibitem{kovactao2024} V.~Kova\v{c} and T.~Tao,
  \href{https://arxiv.org/abs/2406.17593v4}{\emph{On several irrationality
  problems for Ahmes series}}, Acta Math. Hungar. \textbf{175} (2025), no.~2,
  572--608,
  doi:\href{https://doi.org/10.1007/s10474-025-01528-0}{10.1007/s10474-025-01528-0};
  arXiv:\href{https://arxiv.org/abs/2406.17593v4}{2406.17593v4}.
  Page references are to arXiv:2406.17593v4.
\bibitem{bado2026} I.~O. Bado,
  \emph{Prime-Support Rigidity and Primitive Pseudo-Greedy Dynamics:
  Partial Progress on Erd\H{o}s Problem~\#243}, preprint posted September 2026,
  doi:\href{https://doi.org/10.13140/RG.2.2.36612.08325}
  {10.13140/RG.2.2.36612.08325}.
\bibitem{erdosproblemaday243} P.~White with Claude (Anthropic),
  \href{https://erdosproblemaday.com/report/243}{\emph{Erd\H{o}s \#243:
  working report}}, Erd\H{o}s Problem a Day, page dated 12 August 2026.
  AI-assisted, unrefereed working report.
\bibitem{stevenhagenlenstra1996} P.~Stevenhagen and H.~W. Lenstra, Jr.,
  \href{https://pub.math.leidenuniv.nl/~lenstrahw/PUBLICATIONS/1994c/art.pdf}
  {\emph{Chebotar\"ev and his density theorem}}, Math. Intelligencer
  \textbf{18} (1996), no.~2, 26--37,
  doi:\href{https://doi.org/10.1007/BF03027290}{10.1007/BF03027290}.
  Page references are to the linked author version dated 23 March 1995.
\bibitem{vandoornlitang2026} W.~van Doorn, Y.~Li and Q.~Tang,
  \href{https://arxiv.org/abs/2603.28636}{\emph{Optimal bounds for an
  Erd\H{o}s problem on matching integers to distinct multiples}},
  preprint arXiv:2603.28636v1, 2026.
\bibitem{koizumi2025} J.~Koizumi,
  \href{https://math.colgate.edu/~integers/aa28/aa28.pdf}
  {\emph{Irrationality of the reciprocal sum of doubly exponential sequences}},
  Integers \textbf{26} (2026), Paper No.~A28, 17 pp.,
  doi:\href{https://doi.org/10.5281/zenodo.18714404}{10.5281/zenodo.18714404}.
  Locators refer to this published version.
\bibitem{erdosproblems} T.~F. Bloom,
  \href{https://www.erdosproblems.com/243}{\emph{Erd\H{o}s Problem \#243}},
  supplied snapshot of 28 July 2026; not a live status verification.
\bibitem{formalconjectures243} The Formal Conjectures Authors,
  \href{https://github.com/google-deepmind/formal-conjectures/blob/f776d2f2039351b00737ffcafb9d7d7666e1d9af/FormalConjectures/ErdosProblems/243.lean}
  {\emph{FormalConjectures.ErdosProblems.\texttt{243}}},
  Lean source at commit \texttt{f776d2f}, 2025. Formal problem statement,
  not a proof.
\bibitem{erdosstraus1974} P.~Erd\H{o}s and E.~G. Straus,
  \href{https://msp.org/pjm/1974/55-1/pjm-v55-n1-p08-s.pdf}
  {\emph{On the irrationality of certain series}}, Pacific J. Math.
  \textbf{55} (1974), no.~1, 85--92.
\bibitem{hancltijdeman2008} J.~Han\v{c}l and R.~Tijdeman,
  \href{https://www.impan.pl/shop/en/publication/transaction/download/product/82186}
  {\emph{On the irrationality of polynomial Cantor series}},
  Acta Arith. \textbf{133} (2008), no.~1, 37--52,
  doi:\href{https://doi.org/10.4064/aa133-1-3}{10.4064/aa133-1-3}.
  Locators refer to the published version.
\bibitem{duverneykurosawashiokawa2020} D.~Duverney, T.~Kurosawa and I.~Shiokawa,
  \href{https://danielduverney.fr/documents/theorie-des-nombres/Tsukuba.pdf}
  {\emph{Irrationality exponents of certain fast converging series of rational
  numbers}}, Tsukuba J. Math. \textbf{44} (2020), no.~2, 235--250,
  doi:\href{https://doi.org/10.21099/tkbjm/20204402235}{10.21099/tkbjm/20204402235}.
  Theorem locators follow the linked 14-page author version.
\bibitem{crmarickovac2025} T.~Crmari\'{c} and V.~Kova\v{c},
  \href{https://arxiv.org/abs/2504.18712v1}
  {\emph{On the irrationality of certain super-polynomially decaying series}},
  Colloq. Math. \textbf{179} (2025), 55--68,
  doi:\href{https://doi.org/10.4064/cm9628-5-2025}{10.4064/cm9628-5-2025}.
  Theorem locators follow arXiv:2504.18712v1.
\bibitem{kouli2020} A.~Koutsoukou-Argyraki and W.~Li,
  \href{https://isa-afp.org/entries/Irrational_Series_Erdos_Straus.html}
  {\emph{Irrationality Criteria for Series by Erd\H{o}s and Straus}},
  Archive of Formal Proofs, 12 May 2020. An Isabelle/HOL formalisation;
  the archive entry identifies the results formalised.
\bibitem{dlmf_gamma} National Institute of Standards and Technology,
  \href{https://dlmf.nist.gov/5.11.E12}{\emph{Digital Library of Mathematical
  Functions}, \S5.11(iii), formula 5.11.12}, accessed 16 September 2026.
\bibitem{elabdalaoui2015} E.~H. el Abdalaoui, M.~Lema\'nczyk and T.~de la Rue,
  \emph{A dynamical point of view on the set of $\mathcal B$-free integers},
  Int. Math. Res. Not. IMRN (2015), no.~16, 7258--7286,
  doi:\href{https://doi.org/10.1093/imrn/rnu164}{10.1093/imrn/rnu164}.
  The cited definition is also in \S1.2 of
  arXiv:\href{https://arxiv.org/abs/1311.3752v3}{1311.3752v3}.
\end{thebibliography}

